\documentclass[11pt]{article}
\usepackage[utf8]{inputenc}
\usepackage{amsmath,amsfonts,amssymb}
\usepackage{amsthm}
\usepackage{mdframed}
\usepackage{graphicx}
\graphicspath{{./Photos/}}
\usepackage[english, russian]{babel}
\usepackage{stmaryrd}
\usepackage{hyperref}
\usepackage{setspace}
\pdfimageresolution=1200
\usepackage{hyperref}
\usepackage{doi}
\usepackage{enumitem}
\usepackage{ragged2e} 
\usepackage[numbers]{natbib}
\usepackage[dvipsnames]{xcolor}
\usepackage{multirow}
\usepackage{array}
\usepackage{tikz}
\usepackage{lipsum}
\usetikzlibrary{arrows.meta, positioning, decorations.pathmorphing}
\usepackage[a4paper, margin=0.75in]{geometry}
\usepackage{afterpage}
\usepackage{titlesec}

\usepackage{titlesec}

\usepackage{fancyhdr}

\pagestyle{fancy}
\fancyhf{} % Clear all header/footer fields

% Right-aligned section title in the header
\fancyhead[R]{\rightmark}
\fancyfoot[C]{\thepage}

% Add the horizontal line under the header
\renewcommand{\headrulewidth}{0.4pt}

% Ensure section marks are updated properly
\renewcommand{\sectionmark}[1]{\markright{#1}}

\newcommand\blankpage{%
    \null
    \thispagestyle{empty}%
    \addtocounter{page}{-1}%
    \newpage}


\definecolor{crimson}{rgb}{0.644, 0.109, 0.187}
\newtheorem{theorem}{Theorem}

\newenvironment{boxtheorem}[1]
  {\begin{mdframed}\noindent\textbf{#1}\itshape\space}
  {\end{mdframed}}
  
 \newenvironment{boxt}[1]
  {\begin{mdframed}\noindent\textbf{#1}\normalfont\space}
  {\end{mdframed}}
  
\newcommand{\E}{\mathcal{E}}
\newcommand{\N}{\mathbb{N}}
\newcommand{\Z}{\mathbb{Z}}
\newcommand{\R}{\mathbb{R}}
\newcommand{\sto}{\text{stop}}
\newcommand{\im}{\text{Im }}



\title{Delocalization of Eigenvectors in Random Matrices}

\author{Georgi Ivanov\\Advisor: Dr. Kevin Yang}
\date{}
\doublespacing
\begin{document}
\renewcommand{\refname}{References}
\renewcommand{\contentsname}{Contents}


\pagenumbering{roman}
% Title Page
\begin{titlepage}
    \begin{center}
        \vspace*{4cm}
        
        \Huge{\textcolor{crimson}{Delocalization of Eigenvectors \\ \vspace{-0.5cm} in Random Matrices}}
        
        \vspace{2cm}
        
        \Large{A \textsc{thesis presented\\
        \vspace{-0.5cm}
		by\\
		\vspace{-0.5cm}
        Georgi Iliyanov Ivanov\\
        \vspace{-0.5cm}
			to\\ \vspace{-0.5cm}
        The Department of Mathematics}}
        
        \vspace{1cm}
        
        \Large{\textsc{in partial fulfillment of the requirements\\ \vspace{-0.5cm}
        for the degree of\\ \vspace{-0.5cm}
        Bachelor of Arts\\ \vspace{-0.5cm}
        in the subjects of\\ \vspace{-0.5cm}
        Mathematics and Statistics}}
        
        \vspace{0.5cm}
                \begin{center}
    \Large{\textsc{Advised by Dr. Kevin Yang}}
    \end{center}
    \vspace{0.5cm}
        \includegraphics[scale=1.25]{harvard_logo.png}
        
        \Large{\textsc{Harvard University\\ \vspace{-0.5cm}
        Cambridge, MA\\ \vspace{-0.5cm}
        March 2025}}
        
    \end{center}
    \afterpage{\blankpage}
\end{titlepage}

% Abstract


    \begin{center}
    
    \vspace*{2cm}
    \LARGE{\textcolor{crimson}{Delocalization of Eigenvectors in Random Matrices}}
    

    \end{center}
    \vspace{1cm}
    \begin{center}
        {\Large{\textsc{Abstract}}}
    \end{center}
    \begin{center}
    
    
	 \justify{  
	 In this thesis, we investigate the delocalization properties of eigenvectors in random matrices with a particular focus on the non-mean field model of the random band matrix. We provide an overview of recent advancements that employ the stochastic flow method, placing them in the context of universality classes and Anderson localization. The approach enables rigorous bounds on eigenvector statistics using resolvent identities and diagrammatic perturbation methods. We include detailed estimates on the drift terms of the dynamical $T$-equation and applications of martingale theory to the control of the quadratic variation. The final section explores the extension of these results beyond the Gaussian distribution. We prove a Five Moment Theorem that addresses the universality of delocalization behavior, using the Lindeberg exchange strategy. 
	 }
    \end{center}
  



% Acknowledgments

\thispagestyle{empty}
\clearpage

\begin{center}
    \Large{\textbf{Acknowledgments}}\vspace{0.1cm}
\end{center}
\begin{center}
\justify{\textsc{First, and foremost}, I would like to thank Kevin Yang, without whom this thesis would not be possible. I could not have asked for a more supportive advisor - one who allowed me to explore ideas freely, whilst nudging me towards the concrete steps that  gave them shape. I am also grateful to the Office of Undergraduate Research and Fellowship for supporting my work with a HCRP grant. \vspace{0.25 cm} 

\noindent \textsc{I want to thank all those from whom I learned over the last 4 years}:\\
My first teacher at Harvard, Hakim Walker, for inspiring me to leap into world of pure mathematics. \\
\noindent My concentration advisor, Denis Auroux, for giving me the confidence to venture into graduate courses.\\ 
\noindent Joe Blitzstein, for making me fall in love with statistics.\\
Yum-Tong Siu, for showing me that I can handle the true intensity of mathematics. \vspace{0.25cm}

\noindent \textsc{I am grateful to my mentors, without whom I wouldn't be here:}\\
\noindent Luis Lomeli, for bringing out the best of me, teaching me to be a bulldog. \\
\noindent Nikola Karavasilev, who gave me a taste of real science, when it was nowhere to be found around me. \\
\noindent My high school teacher, Elena Kiselova, for sowing the seeds of trees I am now in the shade of. \\
\noindent Wasalu Jaco, whose lived experience gave me clarity amidst the heat. \vspace{0.25 cm}

\noindent \textsc{I want to thank the friends who made this foreign land hospitable:}\\
Jordan, the most gracious person I know, for being a strong positive influence in my life. \\
Tega, for being there when things were most difficult. \\
Jan, for bringing brotherly competition and pushing me to be better. \\
John, for introducing me to powerlifting, which has become the bedrock I fall back on. \\
My blocking mates, Cameron and Leah, for all the late night conversations. \\
Kubra and Bolu, for making me feel seen and heard since day one on campus.
\vspace{0.25 cm}

\noindent \textsc{I am grateful to my friends back home} - Stani, Eli, Lily and Georgi - for being by my side throughout the years, unwavering in their belief and unapologetic in their support.\vspace{0.25 cm}

\noindent \textsc{I am thankful to my entire family for their patience and guidance.} I am grateful to my father, Iliyan, for teaching me discipline and my mother, Irina, for the gift of intuition. I want to thank my brothers, Vladislav and Kiril, and my sisters, Dimana and Rozena, for keeping me rooted and never letting me forget who I am doing this for. Last, but not least, I want to thank our people, who through culture, stories, and celebrations have instilled in me my highest value: Sophia.\vspace{0.25 cm}

\noindent \textsc{This thesis is dedicated to the memory of grandma Rumi, who was the first to show me the beauty of mathematics.}}

\end{center}

\newpage

\vspace*{10 cm}
\begin{center}
\Large{\textsc{В памет на Баба Руми}}
\thispagestyle{empty}
\end{center}

\newpage
\tableofcontents
\clearpage

\pagenumbering{arabic}
\setcounter{page}{1}
\section{Introduction}
\label{sec:intro}
The topic of random matrix theory (RMT) was originally motivated by the study of complex quantum systems, described by a Hamiltonian  - the operator that models the interactions between constituent particles. Even when all couplings are known, the daunting and often impossible task of diagonalizing the operator to obtain its spectral decomposition prevents calculating its exact energy levels. To tackle this, Eugene Wigner introduced in his seminal 1957 paper \cite{wigner} the phenomenological model of a random matrix with independent and identically distributed (i.i.d) entries. His goal was to explain the empirical observation that the spectral gaps of large nuclei follow the same statistics, regardless of the material. 

\vspace{0.5cm}
\noindent The core idea was to ignore all physical details of the system except for the constraint on its symmetry type, defined by the presence of time-reversal and rotational symmetry, or the lack-thereof. The result, now known as the Wigner semicircle law, states that the empirical eigenvalue distribution of \( N \times N \) self-adjoint random matrices with independent entries $H_{ij}\sim [0, 1/N]$, up to the symmetry constraint $H_{ij} = \overline{H_{ji}}$, is given by $\varrho_{SC}(x)=\frac{1}{2\pi} \sqrt{(4-x^2)_+}$ \cite{dynamic}. What Wigner proved is a type of law of large numbers that holds in the limit \( N \to \infty \) and is independent of the distribution of individual entries. The prediction that the spacing of individual eigenvalues $\lambda_1(H)\leq \cdots \leq \lambda_n(H)$ follows universal gap statistics was later refined into the Wigner-Dyson-Mehta (WDM) universality conjecture \cite{mehta}, which asserts the latter hypothesis is true for \textit{Wigner matrices} \cite{taoWDM}. This RMT ensemble defines a general class of mean-field models and is part of the active effort of understanding the different universality classes. 

\vspace{0.5 cm}
\noindent Another subject of great importance is the \textit{Poisson universality class}, closely related to the phenomenon of localized eigenstates. While Wigner's original theory had no mention of eigenvectors \cite{7}, their properties are essential in studying the behavior of more general random operators. One historically relevant example is the Anderson localization model, introduced by Philip W. Anderson in 1958, which explores how disorder affects quantum transport \cite{anderson}. The author considered the behavior of a quantum particle moving on a lattice $\Z^d$ with random on-site energies, represented by the random Schrödinger operator \( H_{RS}=-\Delta + \lambda V \), where \( \Delta \) is the discrete Laplacian and \( \lambda V \) is a random diagonal perturbation \cite{randomoperators}. In his work, Anderson posited that there exists a \textit{mobility edge} that separates a transition between localized and delocalized eigenvector statistics. 
\vspace{0.5 cm}

\noindent Localization for large values of $\lambda$ was first rigorously demonstrated using multiscale analysis \cite{24} with alternative approaches, such as the fractional moment method \cite{1}, later ensuing. The conjectured delocalization for small $\lambda$ is supported by extensive numerical evidence, yet a rigorous proof has remained out of reach \cite{bandSDE}. The primary reason for the scarcity of rigorous results has been the general lack of exactly solvable models of sufficient generality. \cite{kravtsov2009}.  To address these gaps, it is crucial to explore systems that interpolate between localized and delocalized regimes while preserving analytical tractability. Two examples that we will present are the non-invariant Gaussian random matrix ensemble and the Random Band Matrix model (RBM), which will be the primary focus of this thesis. 

\subsection*{Non-Invariant Gaussian Matrix}
The non-invariant Gaussian random matrix ensemble, called also the critical random matrix ensemble (CRMT) \cite{mirlin1996}, is defined for Hermitian matrices \( H \) with independent Gaussian-distributed off-diagonal elements \cite{kravtsov2009}:$$
\langle H_{nm} \rangle = 0, \quad \langle |H_{nm}|^2 \rangle = 
\begin{cases} 
\beta^{-1}, & n = m \\ 
\frac{1}{2} \left[ 1 + \left( \frac{n - m}{b} \right)^2 \right]^{-1}, & n \neq m 
\end{cases}
,$$
where \( \beta = 1, 2, 4 \) corresponds to the universality class of orthogonal (GOE), unitary (GUE), and symplectic (GSE) matrices, respectively, and \( b > 0 \) is a tunable parameter. The CRMT emerged in the study of the 3D Anderson model as the critical value  instantiation of a Power-law RBM (defined below). By mapping the system onto a nonlinear $\sigma$-model with nonlocal interaction and using renormalization group (RG) methods, the authors found that for $a = 1$, the model reaches a critical point with multifractal eigenstate behavior, and the spectral statistics exhibit an intermediate regime. The critical nature of CRMT is encoded in the decay \( |n - m|^{-2} \), reminiscent of the Anderson model's localization in coordinate space. The parameter \( b \) influences the spectrum of fractal dimensions \( d_n(b) \), governed by:
\[
d_2 = 
\begin{cases} 
1 - c_\beta B^{-1}, & B \gg 1 \\ 
c_\beta B, & B \ll 1 
\end{cases}
\]
where \( B = b \pi \beta/2 \), and \( c_\beta \) is a constant specific to the symmetry class. Notably, this leads to a duality relationship \( d_2(B) + d_2(B^{-1}) = 1 \) that has been numerically verified with high precision \cite{kravtsov2009}. The level spacing distribution \( P(s) \) of CRMT exhibits hybrid behavior, combining a Poisson tail with Wigner-Dyson statistics in the bulk. While this model effectively captures key aspects of multifractality and the coexistence of localized and delocalized phases, it lacks the rigor necessary to fully characterize the critical behaviors and universal properties of these transitions. Particularly helpful should be the aforementioned fact that the CRMT is a critical value instantiation of a random band matrix with power-law entires. Given the rich and active topic of RBMs, we now turn to our main object of study.  

\subsection*{Random Band Matrices (RBM)}
The Random band matrix model is of great interest by itself because it serves as an interpolation between Wigner matrices and the random Schrödinger operator $H_{RS}$ \cite{dynamic}. A RBM $(H_{xy})_{x, y\in \Gamma}$, with centered complex random variables (r.v.), independent up to symmetry $H_{ij} = \overline{H_{ji}}$, can represent a $d$-dimensional quantum system on a graph $\Gamma = \llbracket 1, N\rrbracket ^d$ with the effective distance being of order $W<\frac{N}{2}$, defined to be the band width of the model \cite{7, 21}.  What this means is that for $|i-j|_{N} = \min\left([x-y]_N, [y-x]_N\right)>W$ we have $H_{ij}=0$. It is standard to also normalize the covariance matrix $S_{xy} = \mathbb{E} |H_{xy}|^2$, s.t. $\sum_x S_{xy}=1$ for any $x\in \Gamma$. The model is conjectured to exhibit both localization for $W\ll W_c$ and delocalization for $W\gg W_c$ w.r.t a critical transition band $W_c$ that depends on the band width $W$ and dimension $d$ \cite{7}:  \[
W_c = 
\begin{cases} 
\sqrt{N} & \text{for } d=1 \\ 
\sqrt{\log N} & \text{for } d=2 \\
O(1) & \text{for } d\geq 3

\end{cases}
\]
The Anderson model and the RBM are expected to have the same properties when $\lambda  \approx \frac{1}{W}$ \cite{7}, an observation supported with extensive numerical evidence \cite{arbm1, arbm2}. For localization, the current best bound is up to $W\ll N^{1/4}$ \cite{13}. On the other side of the transition, there is a long line of work of iterative improvements.\vspace{0.255 cm}

\noindent Delocalization and quantum diffusion, explained in detailed below, have been rigorously established for the case of $d\geq 7$ \cite{8, 9, 40, highdim} using diagrammatic representations. The first two papers also 
concluded a strong form of delocalization and GUE statistics for  $W\gg N^{3/4}$ \cite{8, 9} in the one-dimensional case. More recently, Dubova and Yang improved on the latter bound for $1d$ Gaussian random band matrices, assuming $W\gg N^{8/11}$ \cite{bandSDE}, by using the flow method \cite{RP, soosten}. Sections \hyperref[sec:flow]{2} through \hyperref[sec:drift]{7} are dedicated exclusively to their work, which analyzed the truncation of the observable dynamic hierarchy. Yau and Yin were later able to analyze the full hierarchy, establishing delocalization for $W\geq N^{{1/2}+\epsilon}$ in the one-dimensional Gaussian case \cite{yauyin}. The continued success of this approach has allowed for proving delocalization in two dimension when $W\geq N^\epsilon$ \cite{twodim}. The latter sequence of notable results is the reason for dedicating the bulk of this thesis to the study of the tools, utilized by the stochastic flow method. Furthermore, given the underlying Gaussian assumption in all previous work, the question of the non-Gaussian case is a highly motivated one. As such, the remainder of our work (\hyperref[sec:non-gaussian]{Section 8}) will be focused on relaxing the Gaussian assumption by proving a Five Moment Matching theorem. With this said, we can start with our model of interest:

\newpage

\subsection*{The Model}
\label{model}
The matrix model $(H_{ij})_{i, j\in \Gamma}$ that will be the subject of our examination throughout this thesis is defined as follows: identify $\Z_n$ with $\Gamma=\{1, \cdots, N\}$ and equip it with the aforementioned periodic distance $|i-j|_N$, s.t. $\forall_{|i-j|_N>W} H_{ij}=0$. Now, for a symmetric and compactly supported probability density function (PDF) on $\mathbb{R}$, define the doubly stochastic matrix $(S_{xy})_{x, y\in \Gamma}$ as: $$\mathbb{E}|H_{xy}|^2=S_{xy} := Z_{N, W}^{-1}f\left(\frac{|x-y|_N}{W}\right),$$
where $Z_{N, W}\asymp W$ is a normalizing constant bounded deterministically above and below by $W$. A technical, but necessary assumption is for $S$ to admit a matrix square root $S^{1/2}$, satisfying the same properties for a similarly symmetric and compactly supported PDF $f'$. \vspace{0.75 cm}
\begin{center}
\begin{tikzpicture}
    % Define parameters
    \def\n{4} % Matrix dimension
    \def\w{0.7} % Half-width of the band
    
    % Draw the proper matrix brackets
    \draw[line width=1.5pt] (-0.1,-0.1) -- (-0.1,\n+0.1);
    \draw[line width=1.5pt] (\n+0.1,-0.1) -- (\n+0.1,\n+0.1);
    % Left bracket extensions (CORRECTED DIRECTION)
    \draw[line width=1.5pt] (-0.1,-0.1) -- (0.4,-0.1);
    \draw[line width=1.5pt] (-0.1,\n+0.1) -- (0.4,\n+0.1);
    % Right bracket extensions (CORRECTED DIRECTION)
    \draw[line width=1.5pt] (\n+0.1,-0.1) -- (\n-0.4,-0.1);
    \draw[line width=1.5pt] (\n+0.1,\n+0.1) -- (\n-0.4,\n+0.1);
    
    % Shade the band region along the diagonal
    \fill[gray!50] (0,\n) -- (\w,\n) -- (\n,\w) -- (\n,0) -- (\n-\w,0) -- (0,\n-\w) -- cycle;
    
	\fill[gray!50] (0,0) -- (\w,0) -- (0,\w) -- cycle; % Lower left triangle
    \fill[gray!50] (\n-\w,\n) -- (\n,\n) -- (\n,\n-\w) -- cycle; % Upper right triangle    
    
    % Write the label
    \node at (-0.8,\n/2) {$H =$};
    
    % Add text label for band width
    \draw[<->] (\n/2+\w/2,\n/2+\w/2) -- (\n/2-\w/2,\n/2-\w/2) node[midway,sloped,above] {$2W$};
    
    % Add zeros in the correct regions (outside the band)
    \node at (\n/4*3+0.5,\n/4*3-0.5) {$0$};
     \node at (\n/4*3-0.5,\n/4*3+0.5) {$0$};
    \node at (\n/4-0.5,\n/4+0.5) {$0$};
        \node at (\n/4+0.5,\n/4-0.5) {$0$};
\end{tikzpicture}
\end{center}
\vspace{0.255 cm}

\noindent In our analysis, we will utilize the flow method (\hyperref[sec:flow]{Section 3}), where the spectral parameter $z=E+i\eta$ of the resolvent of $H$ (defined in \hyperref[sec:prereq]{Section 2}) is varied at constant-speed in the upper-half plane $\mathbb{H}$ ($E\in\mathbb{R}$ and $\eta>0$) as the entries of $H$ are realized as Brownian motion \cite{bandSDE}. We will make two key assumptions. The first is a diffusion time scale of $\eta \asymp W^2 N^{-2}$ and the second is that $\eta\geq W^{-3/4+\eta}$ for some $\eta>0$. Those will be discussed in detail later on in Section 3.3.


\newpage
\section{Prerequisites}
\label{sec:prereq}
The main tool used to study both Anderson localization and random matrix statistics is the \textit{Green's function} or \textit{resolvent}. For a random matrix $H$, it is defined as: $$G(z) := (H-zI_N)^{-1} = (H-z)^{-1},$$
where $z\in \mathbb{C}\backslash \{\lambda_1, \cdots, \lambda_N\}$ is the spectral parameter ranging over the complement of the eigenvalues $\lambda_1(H)\leq \cdots \leq \lambda_N(H)$ of $H$. Observe that for any eigenvector $\mathbf{u}_\alpha$ we have:
$$H\mathbf{u}_\alpha = \lambda_\alpha \mathbf{u}\Leftrightarrow (H - z)\mathbf{u}_\alpha = (\lambda_\alpha-z)\mathbf{u}_\alpha$$
\begin{equation}
\Rightarrow G(z)\mathbf{u}_\alpha = (H-z)^{-1}\mathbf{u}_\alpha = \frac{1}{\lambda_\alpha-z}\frac{(\lambda_\alpha-z)\mathbf{u}_\alpha }{H-z}=\frac{1}{\lambda_\alpha -z}\mathbf{u}_\alpha\end{equation}

\noindent This is equivalent to $\mathbf{u}_\alpha$ being eigenvectors for the resolvent with $\frac{1}{\lambda_\alpha-z}$ as the eigenvalues. This also means that we have by spectral decomposition: 
\begin{equation*}G(z) =\sum_{i=1}^N \frac{\mathbf{u}_\alpha\mathbf{u}_\alpha^*}{\lambda_\alpha-z}.\tag{1.1}\end{equation*}
Now, recall that their empirical measure for $H$ is defined as \cite{dynamic}: $$\varrho_N(dx):=\frac{1}{N}\sum_{j=1}^N \delta(x-\lambda_j)dx,$$
Let $m(z)=m_N(z) = \int_\R \frac{\varrho_N(x)}{x-z}dx$ be its Stieltjes transform for $z = E+i\eta$, where $E\in \R$ and $\eta>0$. By the relationship we found above, we have the following  identity, which follows from the definition of the Dirac delta function: 
\begin{equation*}\scalebox{1}{$m(z)= \int_\R \frac{\varrho_N(x)}{x-z}dx = \int_\mathbb{R} \frac{1}{N}\frac{\sum_{j=1}^N \delta(x-\lambda_j)dx}{x-z}dx = \frac{1}{N}\sum_{j=1}^N\int_\R\frac{\delta(x-\lambda_j)}{x-z}dx =\frac{1}{N} \sum_{j=1}^N \frac{1}{\lambda_j-z} = \frac{1}{N}\text{Tr}(G(z))$}\end{equation*}
From the moments method \cite{dynamic} we know that the Stieltjes transform $m_N(z)$ of the empirical distribution converges in probability to the Stieltjes transform $m_{SC} = \frac{1}{2}(-z+\sqrt{z^2-4})$ of the semicircle law as $N\rightarrow\infty$. This means that in the limit, with probability one the self-consistent equation holds: 
\label{self-consistent}
\begin{equation}
m(z) = -\frac{1}{z+m(z)} \Leftrightarrow m(z) + \frac{1}{m(z)} + z=0
\end{equation}


\subsection*{Stieltjes transform identities}
Let us establish several facts and lemmas that we will need throughout our later work. We can begin by observing that: $$m(z)(m(z)+z) = \frac{1}{4}(-z+\sqrt{z^2-4})(z+\sqrt{z^2-4})=\frac{1}{4}(-z^2+z^2-4)=-1$$
$$\Rightarrow |m(z)| = \frac{1}{|m(z)+z|^{-1}}$$
Now, take the imaginary part of $(2)$: $$\text{Im}(m(z)) + \text{Im}\left(\frac{1}{m(z)}\right)+\eta=0$$
We can calculate the second term by simply expanding: $$\frac{1}{m(z)} = \frac{1}{\text{Re}(m(z))+\text{Im}(m(z))i} = \frac{\text{Re}(m(z))-\text{Im}(m(z))i}{|m(z)|^2}\Leftrightarrow \text{Im}\left(\frac{1}{m(z)}\right) = -\frac{\text{Im}(m(z))}{|m(z)|^2}$$
$$\Rightarrow\text{Im}(m(z)) + \text{Im}\left(\frac{1}{m(z)}\right)+\eta=\text{Im}(m(z)) -\frac{\text{Im}(m(z))}{|m(z)|^2}+\eta=0$$
$$\Rightarrow \text{Im}(m(z))\left(1-\frac{1}{|m(z)|^2}\right)=-\eta$$
And since $\text{Im}(m(z))>0$ and $\eta>0$, we must have:
\begin{equation*}\label{2.1} 1-\frac{1}{|m(z)|^2}<0\quad \Leftrightarrow\quad |m(z)|<1, \tag{2.1}
\end{equation*}
which is also equivalent to $\text{Im }m \asymp 1$, where $\asymp$ means being bounded above and below up to a fixed, positive factor. Furthermore, let us rewrite the latter equation as the following two identities: \begin{equation*}\label{2.2}1-|m(z)|^2 = \frac{\eta|m(z)|^2}{\text{Im}(m(z))}
\tag{2.2}
\end{equation*}
\begin{equation*}\label{2.3}\text{Im}(m(z))= \frac{\eta|m(z)|^2}{1-|m(z)|^2 }
\tag{2.3}
\end{equation*}
Now, let $E\in [-10, 10]$ and $\eta\in (0, 10]$. We claim that $|m(z)|\geq c>0$ uniformly. Indeed, $m(z)$ is continuous, this is equivalent to $m(z)\neq 0$, which can be verified as follows: $$m(z) = \frac{-z+\sqrt{z^2-4}}{2}=0 \Leftrightarrow \sqrt{z^2-4}=z$$
$$\Rightarrow z^2-4=z^2\Leftrightarrow -4=0$$ 





\subsection*{Resolvent Identities}
We can define the augmented minor Green's function for any index $\alpha$ as $G(z)^{(\alpha)} = (H^{(\alpha)}-w_t)^{-1}$, where: $$\left(H^{(\alpha)}\right):=H_{ab}\cdot \mathbf{1}(a\neq \alpha)\cdot \mathbf{1}(b\neq \alpha),$$ 
so the augmented matrix is still $N\times N$ with the $\alpha$-row and -column set to zero \cite{dynamic}. As a consequence: 
\begin{equation*}
G_{xy}^{(\alpha)} = \left\{\begin{array}{cc}
  (-z)^{-1}  & x=y=i\\
  0 & x=\alpha\ \text{XOR } y = \alpha\\
  G_{xy}^{[\alpha]} = (H^{[\alpha]}-z)_{xy}^{-1} & x\neq \alpha, y\neq \alpha
  \end{array}
  \right.
\end{equation*}
Before we derive the estimate, let us first establish several identities:

\begin{boxt}{(Matrix Identities).}
Let $A$ and $B$ be arbitrary matrices. Provided that the inverses exist, we can derive the following identities: 
$$\frac{1}{A}=A^{-1} =A^{-1}(A+B)(A+B)^{-1}=(A+B)^{-1}+A^{-1}B(A+B)^{-1} = \frac{1}{A+B}-\frac{1}{A}B\frac{1}{A+B}$$
\begin{equation*}\Rightarrow \frac{1}{A+B}=\frac{1}{A}-\frac{1}{A}B\frac{1}{A+B} = \frac{1}{A}-\frac{1}{A+B}B\frac{1}{A}\tag{2.4}\end{equation*}
where the last equality follows from simply swapping the order $A^{-1}=A^{-1}(A+B)^{-1}(A+B)$. \end{boxt}

\noindent Using this, we can now set $A = H_t^{(\alpha)}-z$ and $B  = H_t-H_t^{(\alpha)}$, hence by substituting above, we have: $$\frac{1}{H_t-z} = G_{t}(z) = \frac{1}{H_t^{\alpha}-z}-\frac{1}{H_t-z}\left(H_t-H_t^{(\alpha)}\right)\frac{1}{H_t^{(\alpha)}-z} = G^{(\alpha)}_t-G_t(H_t-H_t^{(\alpha)})G_t^{(\alpha)}$$
But observe that by definition of the augmented minor, the $xy$ entry of our identity above simplifies as: 

\begin{equation*}
G_{xy}=G_{xy}^{(\alpha)}-\sum_{j=1}^N\sum_{k=1}^NG_{xj}(H_{jk}-H^{(\alpha)}_{jk})G_{ky}^{(\alpha)} = G_{xy}^{(\alpha)} - G_{x\alpha}\sum_{k\neq \alpha} H_{\alpha k}G_{ky}^{(\alpha)} = G_{xy}^{\alpha}+\frac{G_{x\alpha}G_{\alpha y}}{G_{\alpha\alpha}}
\end{equation*}
Similarly, we can derive our second resolvent identity, by applying again the same properties, with the special case of $\alpha = x$, s.t. we get:

\begin{equation*}
G_{xy} = G_{xy}^{(x)} - G_{xx}\sum_{k\neq x}H_{xk}G_{ky}^{(x)} =- G_{xx}\sum_{k\neq x}H_{xk}G_{ky}^{(x)}\quad x\neq y
\end{equation*}
Lastly, let us derive the derivatives of the resolvent w.r.t. to the matrix entry $H_{ij}$ using the limit definition. Let us define the perturbed matrix $H' = H(H_{ij}\mapsto H_{ij}+\epsilon)$ at entry $H_{ij}$, where $G' = (H'-z)^{-1}$ and let $\epsilon_{ij} = \mathbf{0}(0_{ij}\mapsto \epsilon)$ be the perturbed zero matrix, which we can write as $\mathbf{1}_{ij} = \mathbf{e}_i^*\mathbf{e}_j$ when $\epsilon=1$. By our matrix identities, we have: $$\partial_{H_{ij}}G(z)_{xy}=\lim_{\epsilon\rightarrow 0} \frac{G_{xy}'-G_{xy}}{\epsilon} = \lim_{\epsilon\rightarrow 0}\frac{\left[G'(H-H')G\right]_{xy}}{\epsilon} = -\lim_{\epsilon\rightarrow 0}\frac{\left[G'\epsilon_{ij}G\right]_{xy}}{\epsilon} = -[G'\mathbf{1}_{ij}G]_{xy} = -G_{xi}G_{jy}$$
\begin{equation*}\label{resolventderivative}
\partial_{H_{ij}}\overline{G(z)_{xy}}= \lim_{\epsilon\rightarrow 0}\frac{\left[\overline{G'}(H^\intercal-H'^\intercal)\overline{G}\right]_{xy}}{\epsilon} =\\\\= -\lim_{\epsilon\rightarrow 0}\frac{\left[\overline{G'}\epsilon_{ji}\overline{G}\right]_{xy}}{\epsilon} = -[\overline{G'}\mathbf{1}_{ji}\overline{G}]_{xy} = -\overline{G_{xj}}\ \overline{G_{iy}}
\end{equation*}
By extension of the Hermitian property, we have $\partial_{\overline{H}_{ij}} \overline{G(z)_{xy}} = \partial_{{H}_{ji}} \overline{G(z)_{xy}}  = -\overline{G_{xi}}\ \overline{G_{jy}}$
\noindent Lastly, the following result is ubiquitous in RMT literature (8.3.(iii) \cite{dynamic}), and as such has its own name: 
\begin{boxt}{Ward identity: }\label{ward} $$\sum_j |G_{ij}|^2 = \frac{1}{\eta}\im G_{ii}$$\end{boxt}
The proof follows directly from the spectral decomposition of the resolvent we showed in $(1)$, by which:
$$G_{ij} = \sum_{\alpha} \frac{\mathbf{u}_\alpha(i) \mathbf{u}_\alpha^*(j)}{\lambda_\alpha - z}\quad \Leftrightarrow \sum_j G_{ij} G_{ji}^* = \sum_{j, \alpha, \beta}\frac{\mathbf{u}_\alpha(i) \mathbf{u}_\alpha^*(j)}{\lambda_\alpha - z} \frac{\mathbf{u}_\beta(j) \mathbf{u}_\beta^*(i)}{\lambda_\beta - z^*}
$$
And since the eigenvectors are orthogonal $\sum_j \mathbf{u}_\alpha^*(j) \mathbf{u}_\beta(j) = \delta_{\alpha\beta}$, we have: $$\sum_j |G_{ij}|^2=\sum_j G_{ij} G_{ji}^* = \sum_\alpha \frac{|\mathbf{u}_\alpha(i)|^2}{|\lambda_\alpha - z|^2}$$ 
We can also take the imaginary part of the first identity:
$$\text{Im}[G_{ii}] = \text{Im}\left[\sum_\alpha \frac{|\mathbf{u}_\alpha(i)|^2}{\lambda_\alpha - z}\right]
= \sum_\alpha \text{Im}\left[\frac{|\mathbf{u}_\alpha(i)|^2}{\lambda_\alpha - E - i\eta}\right]=\sum_\alpha \text{Im}\left[\frac{|\mathbf{u}_\alpha(i)|^2(\lambda_\alpha - E + i\eta)}{(\lambda_\alpha - E)^2 + \eta^2}\right]= $$
$$= \sum_\alpha \text{Im}\left[\frac{|\mathbf{u}_\alpha(i)|^2(\lambda_\alpha - E)}{(\lambda_\alpha - E)^2 + \eta^2} + \frac{|\mathbf{u}_\alpha(i)|^2 \cdot i\eta}{(\lambda_\alpha - E)^2 + \eta^2}\right] 
= \sum_\alpha \frac{|\mathbf{u}_\alpha(i)|^2 \cdot \eta}{(\lambda_\alpha - E)^2 + \eta^2} = \frac{1}{\eta} \sum_\alpha \frac{|\mathbf{u}_\alpha(i)|^2}{|\lambda_\alpha - z|^2}$$
$$\Rightarrow \sum_j |G_{ij}|^2 = \frac{1}{\eta}\im G_{ii}$$




\newpage
\section{Flow Method}
\label{sec:flow}
At the heart of the recent advances in the theory of delocalization lies the flow method, a dynamic technique that will be the primary tool of our exploration. First appearing in the paper of von Soosten and Warzel \cite{RP}, it was used to prove non-ergodic delocalization in the Rozenzweig-Porter model \cite{porter1960}. Subsequent work \cite{soosten} extended the methodology to general Wigner matrices \cite{dynamic}, allowing for a simplified proof of the local semicricle law. The primary idea is to trace the evolution of the Green's function (resolvent) along random characteristic curves, enabling direct concentration estimates and more refined spectral control. More recently, it was used to prove delocalization for the one- and two-dimensional band matrices \cite{yauyin, twodim}\vspace{0.255cm}.
\subsection{Construction}
\noindent Consider the Brownian perturbation of the standard deviation matrix defined as $H(t)_{xy} = \sqrt{S_{xy}}b(t)_{xy}$, where $b(t)_{xy}\sim \mathcal{N}(0, t)$ represents the component Brownian motion. For this system, the matrix resolvent is augmented as $G(t) = (H(t)-w(t))^{-1}$ , where: \begin{equation*}
w(t) = -\frac{1}{m(z)}-tm(z) = z+(1-t)m(z) \tag{3}
\end{equation*}
by the self-consistent equation \hyperref[self-consistent]{(2)}. In order to study this system and its transition from $t=0$ with $G(0) = m(z)$, to $t=1$ with $G(1)=G(z)$, we will employ stochastic calculus. For this purpose, observe that our differentials are $dH(t)_{xy} = \sqrt{S_{xy}}db_{xy}(t)$ and $\partial_tw(t) = -m(z)$. Recall the Îto equation, namely, the statement that a function $f(\vec{b}(t), t)$, dependent on a Brownian random walk  $\vec{b}(t) = (b(t)_1, \cdots, b(t)_d)$ w.r.t to time $t$, has derivative:
$$df(\vec{b}(t), t) = \frac{\partial f}{\partial t}dt + \sum_{i=1}^d \frac{\partial f}{\partial b_i}db_i(t) + \frac{1}{2}\sum_{i=1}^d \frac{\partial^2 f}{\partial b_i^2} dt$$
In our case, by taking the function to be $G(t)$ and using restriction of a symmetric matrix and the latter differentials, we get the following expression:

$$dG(t)_{xy} = \frac{\partial G(t)_{xy}}{\partial t}dt +\sqrt{S_{ij}}\sum_{i\leq j}\frac{\partial G(t)_{xy}}{\partial b(t)_{ij}}db(t)_{ij}
+\frac{S_{ij}}{2}\sum_{i\leq j}\frac{\partial^2 G(t)_{xy}}{\partial b(t)_{ij}^2}(db(t)_{ij})^2$$
To calculate the individual terms, we will use the definition of derivative:
$$\frac{\partial G(t)_{xy}}{\partial t}=\lim_{\epsilon\rightarrow 0}\frac{(H(t)-w(t+\epsilon))^{-1}_{xy}-(H(t)-w(t))^{-1}{xy}}{\epsilon}.$$
The matrix identity $(2.4)$ can help us approximate for small perturbations $\Delta w(t) = w(t+\epsilon)-w(t)$:
$$\left[(H(t)-w(t+\epsilon))^{-1}\right]_{xy}-\left[(H(t)-w(t))^{-1}\right ]_{xy}= \left[ G(t)\Delta w(t) G(t)\right]_{xy}$$$$\Rightarrow \frac{\partial G(t)_{xy}}{\partial t} =\lim_{\epsilon\rightarrow 0} \frac{ \left[ G(t)\Delta w(t) G(t)\right]_{xy}}{\epsilon}=\sum_{j}G(t)_{xj}\frac{dw(t)}{dt}G(t)_{jy}= -m(z)\sum_j G(t)_{xj}G(t)_{jy}$$
Similarly, we can carry out the calculation for the second term, using the independence of entries:

$$\frac{\partial G(t)_{xy}}{\partial b(t)_{ij}}= \lim_{\epsilon\rightarrow 0}\frac{(\sqrt{S_{xy}}b(t+\epsilon)_{ij}+w(t))^{-1}_{xy}-(\sqrt{S_{xy}}b(t)_{ij}+w(t))^{-1}_{xy}}{\epsilon}=$$
$$ =\lim_{\epsilon\rightarrow 0}\frac{[G(t)\sqrt{S}(b(t+\epsilon)-b(t))G(t)]_{xy}}{\epsilon}=\sum_{\alpha, \beta}G(t)_{x\alpha}\sqrt{S_{\alpha \beta}}db(t)_{\alpha\beta}G(t)_{\beta y}$$
Lastly, using the Brownian motion fact that $|db(t)|^2=dt$ and carrying out the summation, we get:
$$dG(t)_{xy}=-m(z)\sum_j G_{xj}G_{jy}dt+\sum_{\alpha, \beta}G_{x, \alpha}\sqrt{S_{\alpha, \beta}}db_{\alpha, \beta}(t)G_{\beta, t}+\sum_j G_{xj}\left(\sum_yS_{jy}G_{yy}\right)G_{jy}dt=$$


$$=\sum_{\alpha, \beta}G_{x\alpha}\sqrt{S_{\alpha \beta}}db_{\beta(t)}G_{\beta y}+\sum_j G_{xj}S_{jy}\left[G_{yy}-m(z)\right]G_{jy}dt = d \mathrm{I}(t) + d \mathrm{II}(t)dt$$
This provides us with the decomposition of a stochastic term $d \mathrm{I}(t)$ and a deterministic term $d \mathrm{II}(t)dt$. 


\subsection{Dynamical $T$-equation}
Having found the differential of the resolvent, we will now proceed with defining the time-dependent $T$-matrix as:
$$T(t)_{ab} = \sum_{x,y}S_{ax}^{\frac{1}{2}}F(t)_{xy}S_{yb}^{\frac{1}{2}},$$
where $F(t)_{xy} = |G(t)_{xy}|^2 = G(t)_{xy}\overline{G(t)_{xy}}$. 
Applying Itô's formula to \( F(t)_{xy} \), we have:
\[
dF(t)_{xy}= d|G_{xy}|^2= G(t)_{xy} d\overline{G(t)}_{xy} + \overline{G(t)}_{xy} dG(t)_{xy} + dG(t)_{xy} d\overline{G(t)}_{xy}.
\]
$$d|G_{xy}|^2 = G_{xy}d\overline{G_{xy}} + \overline{G}_{xy}dG_{xy} + d\left[G_{xy}, \overline{G}_{xy}\right]$$
 $$d|G_{xy}|^2 = G_{xy}d\overline{G_{xy}} + \overline{G}_{xy}dG_{xy} +d\left[\sum_{\alpha, \beta} G_{x\alpha}d H_{\alpha \beta}\overline{G}_{xy}\right]$$
$$d|G_{xy}|^2 = G_{xy}d\overline{G_{xy}} + \overline{G}_{xy}dG_{xy} +\sum_{\alpha,\beta, \gamma, \delta}G_{x\alpha}G_{\beta\gamma}\overline{G_{xy}}\overline{G_{\delta \gamma}}$$
In order to find the derivative of the $T$-matrix, we need to substitute for the earlier terms we found, namely \( dG(t)_{xy} = d\mathrm{I}(t)_{xy} + d\mathrm{II}(t)_{xy} dt \) and its conjugate. By observing that \( d\mathrm{I}(t)_{xy} d\overline{\mathrm{I}(t)}_{xy} \) contributes to the quadratic variation, we can obtain the following expression:
\[
dF(t)_{xy} = dM(t)_{xy} + \Omega(t)_{xy} dt,
\]
where:
\begin{equation*}
dM(t)_{xy} = \overline{G(t)}_{xy} d\mathrm{I}(t)_{xy} + G(t)_{xy} \overline{d\mathrm{I}(t)}_{xy} \end{equation*}
\begin{equation*}
\Omega(t)_{xy} = G_{xy}\frac{\overline{d\mathrm{II}(t)}_{xy}}{dt}+\overline{G_{xy}}\frac{d\mathrm{II}(t)_{xy}}{dt}+\frac{d\left[G_{xy}, \overline{G_{xy}}\right]}{dt}
\end{equation*}
When expanded w.r.t to $d\mathrm{I}(t)$ and $d\mathrm{II}(t)dt$, those terms are:
$$dM(t)_{xy} = \sum_{\alpha, \beta}\overline{G_{xy}}G_{x\alpha}{S_{\alpha\beta}}^{1/2}G_{\beta y}dB_{\beta(t)}+\sum_{\alpha, \beta}{G}_{xy}\overline{G_{x\alpha}}{S_{\alpha\beta}}^{1/2}\,\overline{G_{\beta y}}dB_{\beta(t)}$$


$$\Omega(t)_{xy}=G_{xy}(t)\left(\overline{\sum_j G_{xj}(t)S_{jy}(G_{yy}(t)-m(z))G_{jy}(t)}\right)+$$
$$+\overline{G_{xy}(t)}\left(\sum_j G_{xj}(t)S_{jy}(G_{yy}(t)-m(z))G_{jy}(t)\right)+\sum_{\alpha, \beta}|G_{xa}(t)|^2S_{\alpha, \beta}|G_{x\beta}(t)|^2$$
Hence, by Itô the $T$-equation becomes: 
\begin{equation*}
d{T} = {T}^2 dt + S^{1/2} dM_t(z)S^{1/2} + S^{1/2} \Omega_t(z) S^{1/2}dt \tag{3.1}
\end{equation*}
Now, define the \textit{time-dependent diffusion profile} as:$$\Theta_t = \frac{|m(z)|^2S}{1-t|m(z)|^2S}$$
Let us take its derivative w.r.t to $t$: 
$$\frac{d\Theta_t}{dt} = \frac{d}{dt}\frac{|m(z)|^2S}{1-t|m(z)|^2S} = \frac{-|m(z)|^2S(-|m(z)|^2S)}{(1-t|m(z)|^2S)^2} = \Theta_t^2 \Leftrightarrow d\Theta_t = \Theta_t^2 dt$$
Consider the fluctuation term $\mathcal{E}_t(z)$, i.e the difference between our $T$-matrix and the diffusion profile:
$$\E_t(z):T_t(z)-\Theta_t$$
We can compute its evolution equation $d\E_t(z) = dT_t(z)-d\Theta_t$ by substituting for what we found before:
$$\Rightarrow d\E_t(z) = T_t(z)^2dt - \Theta_t^2dt-S^{1/2}dM_t(z)S^{1/2}+S^{1/2}\Omega_t(z)S^{1/2}dt$$
But observe that:
$$T_t(z)^2-\Theta_t^2 =(\E_t(z)+\Theta_t(z))^2- \Theta_t(z)^2=\E_t(z)^2+\Theta_t\E_t(z)+\E_t(z)\Theta_t$$
\begin{equation*}\Rightarrow d\E_t(z)=\{\Theta_t \E_t(z)+\E_t(z)\Theta_t\}dt +\E_t^2(z)dt - S^{1/2} dM_t(z)S^{1/2} + S^{1/2}
 \Omega_t(z) S^{1/2}dt\tag{3.2}\end{equation*}
The first term $\{\Theta_t \E_t(z)+\E_t(z)\Theta_t\}dt$ makes this equation a matrix-valued linear SDE with a nonlinear stochastic term $W(t) = \E_t^2(z)dt - S^{1/2} dM_t(z)S^{1/2} + S^{1/2} \Omega_t(z) S^{1/2}dt$. This means that we have to apply a method of variation of parameters (a matrix-valued Duhamel formula \cite{bandSDE}). For this, let us identify an integrating factor $U(t, s)$. Observe that: 
$$\partial_{t}\{\text{Id}+(t-s)\Theta_t\} = \Theta_t+(t-s)\Theta_t^2=\Theta_t\frac{1-t|m(z)|^2S+(t-s)|m(z)|^2S}{1-t|m(z)|^2S}=$$
$$=\Theta_t \frac{1-s|m(z)|^2S}{1-t|m(z)|^2S}=\Theta_t\{\text{Id}+(t-s)\Theta_t\}$$
This is equivalent to $\partial_t U(t, s) = \Theta_t U(t, s)$, which makes $U(t, s)$ our evolution operator, since also $U(t, t) = \text{Id}$. We can then define the integral form: $\E_t(z) = \int_0^t U(t, s) W(s) U(t, s)ds,$ and verify that it satisfies our SDE by using Leibniz's integral rule, we get:
$$d\E_t(z) = U(t, t) W(t) U(t, t)dt + \int_0^t d\left[U(t, s)W(s)U(t, s)\right]ds=$$
$$=W(t)dt + \int_0^t\left[\Theta_t U(t, s) W(s) U(t, s) + U(t, s)W(s)U(t, s)\Theta_t\right]dtds =$$
$$= \E_t^2(z)dt - S^{1/2} dM_t(z)S^{1/2} + S^{1/2} \Omega_t(z) S^{1/2}dt + \{\Theta_t \E_t(z)+\E_t(z)\Theta_t\}dt$$
Therefore, by the pathwise uniqueness of solutions to SDEs (Thm. 9.1 \cite{steele}), the following expression satisfies the same SDE:
\begin{equation*}\E_t(z)=\E_t^{D}(z)+\E_t^{M}(z)+\E_t^{S}(z),\tag{3.3}\end{equation*}
where the drift term $\E^D_t(z)$, the martingale term $\E^M_t(z)$, and the squared term $\E^S_t(z)$ are defined as:
$$\E_t^{D}(z)=\int_0^t \{\text{Id} + (t-s)\Theta_t\} S^{1/2} \Omega_s(z) S^{1/2} \{\text{Id} + (t-s)\Theta_t\}ds$$
$$\E_t^{M}(z)=-\int_0^t \{\text{Id} + (t-s)\Theta_t\}S^{1/2} dM_s(z)S^{1/2}\{\text{Id} +(t-s)\Theta_t\}$$
$$\E_t^{S}(z)=\int_0^t \{\text{Id} + (t-s)\Theta_t\}\E_s^2(z) \{\text{Id} + (t-s)\Theta_t\}ds.$$

\noindent It is important to note that while we called $\E^M_t(z)$ a martingale, it is, in fact, not one, as neither are the remaining terms. The reason is that if we consider the mapping $t\mapsto \int_s^t U(t, s)dM_s$, we can see there is dependence on future information $t\geq s$, which is not adapted to the natural filtration of $M_s$. As such, we need a stopping time argument for the latter dynamics \cite{bandSDE}:

\subsection{Stopping times}
For a fixed $\delta_\sto>0$ and $D\lesssim 1$ ($a\lesssim b\equiv a = O(b))$, independent of $N$, define:
\begin{equation*}
\tau_{\text{stop}, 1}=\inf \left\{s\geq 0: \max_{a, b} |\E_s(z)_{ab}|\geq W^{\delta_{\sto}}W^{-\frac{3}{4}}|\im w_s|^{-1}\cdot W^{-1}|\im w_s|^{-\frac{1}{2}}\right\}\wedge 1
\tag{3.4.1}
\end{equation*}

\begin{equation*}
\tau_{\text{stop}, 2}=\inf \left\{s\geq 0: \max_{a, b} 
\frac{|G_s(z)_{ab}-\delta_{ab}m(z)|^2}{(S^{1/2}T_s(z)S^{1/2})_{ab}+S_{ab}^{1/2}+W^{-D}}\geq W^{\delta_{\sto/10}}
\right\}\wedge 1
\tag{3.4.2}
\end{equation*}
\begin{equation*}
\tau_\sto = \tau_{\sto, 1}\wedge \tau_{\sto, 2}\tag{3.4.2}
\end{equation*}
\noindent While the particular choices for constants may seem arbitrary at the moment, they are informed by the bounds we get on the maximal entries of the diffusion profile $\Theta_t$ (\hyperref[diffusion]{Sec. 4.1}), the assumption of $\eta\gg W^{-3/4}$, as well as the technical steps in our analytic argument. That said, we can still gain some intuition as to the meaning of each. With $\tau_{\text{stop},1}$ we have control over the size the error term $\E_s(z) = T_s(z) - \Theta_s$, whereas $\tau_{\text{stop},2}$ is w.r.t to the deviation of the resolvent entries from the semicircle law $m(z)$, normalized by appropriate terms. The reasoning behind the latter definitions is simple - informally, our goal is to establish a bootstrapping mechanism that will allow us to propagate forward the stopping time through a self-reinforcement (continuity) argument. For this purpose, we will instead study the stopped version of $\E_t(z)$, namely the solution $\E^\sto$  to the following equation: $$d\E_t^{\sto}(z) = \left(\Theta_t\E_t^{\sto}(z)+\E_t^{\sto}(z)\Theta_t\right)dt + \mathbf{1}_{t\leq \tau_\sto}\E_t^\sto (z)^2dt$$
$$-\mathbf{1}_{t\leq \tau_\sto}S^{\frac{1}{2}}dM_t(z)S^{\frac{1}{2}}+\mathbf{1}_{t\leq \tau_\sto} S^{\frac{1}{2}}\Omega_t(z)S^{\frac{1}{2}}dt,$$
where the same reasoning as before (Duhamel formula for the variation of parameters), we have that: \begin{equation*}\E_t(z)=\E_t^{D}(z)+\E_t^{M}(z)+\E_t^{S}(z),\tag{3.5}\end{equation*}
$$\E_t^{D, \sto}(z)=\int_0^{t\wedge\tau_\sto} \{\text{Id} + (t-s)\Theta_t\} S^{1/2} \Omega_s(z) S^{1/2} \{\text{Id} + (t-s)\Theta_t\}ds$$
$$\E_t^{M, \sto}(z)=-\int_0^{t\wedge\tau_\sto} \{\text{Id} + (t-s)\Theta_t\}S^{1/2} dM_s(z)S^{1/2}\{\text{Id} +(t-s)\Theta_t\}$$
$$\E_t^{S, \sto}(z)=\int_0^{t\wedge\tau_\sto} \{\text{Id} + (t-s)\Theta_t\}\E_s^2(z) \{\text{Id} + (t-s)\Theta_t\}ds.$$

\noindent Since our stopping times are defined w.r.t. to the natural Brownian filtration, the aforementioned uniqueness theorem guarantees that $\E_t^\sto(z) = \E_t(z)$ for all $\leq \tau_\sto$.  The self-reinforcing argument works as follows:

\begin{enumerate}

\item Initially, we know $\E_0(z) = 0$, so $\tau_{\text{stop}} > 0$ with probability 1.


\item For any $s \leq \tau_{\text{stop}}$, we can control the error terms in the flow SDE by means of martingale inequalities, since we will be working with the stopped martingales $\E_t^{D, \sto}$, $\E_t^{M, \sto}$, and $\E_t^{S, \sto}$. 


\item This will allow us to prove the following theorem in \hyperref[proof-stop]{Section 7}:

\begin{boxtheorem}{Theorem 1: (Stopping time).} $\mathbb{P}[\tau_{\text{stop}, i} \neq  1] \lesssim_D N^{-D}$ for any $i=1, 2$ and $D > 0$.
\end{boxtheorem}

\item Then we can prove that $\forall s \in [0,1]$, $\tau_\sto =1$ with high probability. 

\end{enumerate}

\noindent The particular delocalization bound for $W \gg N^{8/11}$ is so we get the combined error term $W^{-\frac{7}{4}}\eta^{-\frac{3}{2}}$ in equation $(3.4.1)$ to be much smaller than the typical size of the $\Theta$ entries of $O( W^{-1}\eta^{-\frac{1}{2}})$, while maintaining he diffusion time scale $\eta \approx W^2N^{-2}$ (\hyperref[diffusion]{Sec. 4.1}) :
	
$$W^{-\frac{7}{4}}\eta^{-\frac{3}{2}} \ll W^{-1}\eta^{-\frac{1}{2}}$$
$$W^{-\frac{7}{4}}(W^2N^{-2})^{-\frac{3}{2}} \ll W^{-1}(W^2N^{-2})^{-\frac{1}{2}}$$
$$\Rightarrow W^{-\frac{7}{4}}W^{-3}N^3 \ll W^{-1}W^{-1}N$$
$$\quad W^{-\frac{11}{4}} \ll N^{-2}\quad\Leftrightarrow W \gg N^{\frac{8}{11}}$$

\noindent Therefore, by demonstrating that $\tau_{\sto}=1$, we will have established control on the diffusion profile of the band matrix. We will now continue on to the next section, where we present the technical definition of delocalization and explore in more detail the nature of quantum diffusion and the origin of the corresponding bounds on $\Theta_t$ and the size of $\eta$. 


\newpage
\section{Quantum Diffusion and Delocalization}
\label{sec:deloc}

We will need the following definition for the statement of delocalization and quantum diffusion:
\begin{boxt}{Definition 1: (Stochastic Domination).}\label{def-1}
\noindent Consider two sequences of r.v.s parametrized by $s\in S_N$:\vspace{-0.25 cm}
$$X = \{X_N(s): N\in \mathbb{Z}_+, s\in S_N\}, \quad Y = \{Y_N(s): N\in \mathbb{Z}_+, s\in S_N\},$$
Then if for any $\epsilon, D>0$, $\exists N_{\epsilon, D}$, s.t. $\sup_{s\in S_N}\mathbb{P}\left(X_N(s)>N^\epsilon Y_N(s)\right)<N^{-D}$ for $N\geq N_{\epsilon, D}$, 
we write $X \prec Y$ and say that $X$ is \textbf{stochastically dominated} by $Y$ uniformly in $s$. 
\end{boxt}

\subsection{Quantum Diffusion}
\label{diffusion}
Recall our time-dependent $T$-equation $T(t)$ and diffusion profile $\Theta_t$. The standard $T$-matrix and diffusion profile $\Theta$ are just the latter two at the end of the spectral curve when $t=1$, i.e: $$T(z)_{ab} = T_{ab} = \sum_{x,y}S_{ax}^{\frac{1}{2}} |G_{xy}|^2S_{yb}^{\frac{1}{2}}, \quad \Theta := \frac{|m(z)|^2S}{1-|m(z)|^2S}$$
The first theorem established by Dubova and Yang (Theorem 2 \cite{bandSDE}) shows that in the bulk $T\approx \Theta$:
\begin{boxtheorem}{Theorem 2: (Quantum Diffusion)}
For $|E|<2$ fixed, assume $\exists \nu>0$, s.t. $\eta \asymp W^2 N^{-2}$ and $W\geq W^{8/11+\nu}$. Then: $$\max_{x, y}|T_{xy}-\Theta_{xy}|\prec W^{-\frac{7}{4}}\eta^{-\frac{3}{2}}$$
\end{boxtheorem}
This result has a physical interpretation. First, observe
that by \hyperref[2.1]{(2.1)} and  \hyperref[lemma-a11]{Lemma A.1.1}, $1-|m(z)|^2\asymp \eta\Leftrightarrow \exists \alpha>0, 1-|m(z)|^2\geq \alpha\eta$, we have that: 
 $$\Theta =\frac{|m(z)|^2S}{1-|m(z)|^2S} = \frac{|m(z)|^2S}{1-|m(z)|^2-|m(z)|^2(S-\text{Id})}\sim \frac{S}{\alpha\eta-(S-\text{Id})}$$
If we recall the phenomenology of our \hyperref[model]{model}, we can interpret the band width $W$ as the range of interactions in the particle system $\Z_N = \{1, \cdots, N\}$. In this context, since the doubly stochastic matrix is normalized as $\sum_x S_{xy}=1$ and $\sum_y S_{xy}=1$ for fixed $x, y\in \Gamma$, it has a unique interpretation as the transition matrix for a random walk on $\Z_N$. In this context, $S-\text{Id}$ is then clearly its generator, which allows us to consider $\Theta$ as its resolvent. The spectral gap of this random walk on $\Gamma$ with steps of variance $W^2$ is of order $W^2N^{-2}$ \cite{bandSDE}. Standard bounds for diffusion \hyperref[lemma-a2]{Lemma A.2} imply that $\max_{a,b}\Theta_{ab} \lesssim W^{-1}\eta^{-1/2}$, whereas the latter theorem reveals a distinct scaling  $W^{-\frac{7}{4}}\eta^{-\frac{3}{2}}\ll W^{-1}\eta^{-1/2}$. This, in fact, is precisely what characterized quantum diffusion, because instead of following standard resolvent bounds, the maximal entry of the diffusion profile $\Theta$ exhibits different scaling behavior, akin to the phenomenon of interference in quantum systems. Moreover, this quantum correction persists until the relaxation time $\eta^{-1}\asymp N^2W^{-2}$ (Thouless time \cite{17, 32, 33}), allowing us to prove delocalization (\hyperref[proof-deloc]{Theorem 2}).

\subsection{Delocalization}
The following (Thm. 4 \cite{bandSDE}) is a direct consequence of Theorem 2, discussed at length in \hyperref[proof-deloc]{Section 7.3}:

\begin{boxtheorem}{Theorem 3}
\noindent Assume $|E|<2$ is fixed and that $\exists \nu>0$, s.t. $\eta \asymp W^2 N^{-2}$ and $W\geq N^{8/11+\nu}$. Then:
\begin{center}$\max_{x, y}|G_{xy} - \delta_{xy}m(z)|^2\prec W^{-1}\eta^{-\frac{1}{2}}$\end{center}
\end{boxtheorem}

\noindent A direct corollary of this result is the "complete delocalization of (bulk) eigenvectors" \cite{bandSDE, 21}, for which we need some notation. Firstly, for any index $x$ and integer $\ell\geq 1$,  define $P_{x, \ell}(y):=\mathbf{1}[|x-y|\geq \ell]$, i.e the component projection for any indices $x, y$ on the complement of the open ball $B_\ell(x) = \{y:|x-y|<\ell \}$. Given any $\epsilon, \kappa>0$, define the labeling set for the localized to scale $\ell$ eigenvectors in the bulk as: $$\mathcal{A}_{\epsilon, \ell, \kappa} = \left\{\alpha: \lambda_\alpha \in [-E+\kappa, E-\kappa]: \sum_x |\mathbf{u}_\alpha(x)|\|P_{x, \ell}\mathbf{u}_\alpha\|\leq \epsilon\right\}$$
As we will see in \hyperref[proof-deloc]{Section 7.3}, Theorem 3 this implies directly the following (Corollary 5 \cite{bandSDE}): 

\begin{boxtheorem}{Theorem 4: (Delocalization).}
For any $\ell\ll N$ and fixed $\epsilon, \kappa, c>0$: $$\frac{|\mathcal{A}_{\epsilon, \ell, \kappa}|}{N}\lesssim \sqrt{\epsilon}+\mathcal{O}(N^{-c})$$
\end{boxtheorem}
The notation for $\mathcal{A}_{\epsilon, \ell, \kappa}$ contains all indices that are of exponentially localized eigenvectors in $B_{O(\ell)}(x)$ (Remark 7.2 \cite{21}). As such, what the result above implies is that the set of such vectors is vanishingly small up to an error term $\sqrt \epsilon$, i.e all vectors in the bulk are "delocalized". 



\newpage
\section{Estimates }
\label{sec:estimates}

In this section we will state all of the inequalities and stochastic bounds we will need for control over our SDE terms. The \hyperref[sec:appendix]{Appendix} includes the proofs of some, and the references for the rest. From now on, for the sake of brevity we will define $\eta_s = |\im w_s|$ and $\eta_t = |\im w_t|$, the notation of which includes the indexing with any particular time instantiation $t_0\in [0, 1]$. With this, let us state

\begin{boxt}{Estimates}

\begin{enumerate}[label=(E\arabic*)]

\item  $\sup_x\sum_{y}|G_{xy}|\lesssim W^{\frac{\delta\sto}{20}} W^{\frac{1}{2}+\epsilon}\eta_{s}^{-\frac{3}{4}}$ per \hyperref[lemma-b1]{Lemma B.1}.

\item $\max_{x,y}\left\vert G_{xy} - \delta_{xy}m(z)\right\vert\lesssim W^{\frac{\delta_\sto}{20}}W^{-\frac{1}{2}}\eta_s^{-\frac{1}{4}}$ per  \hyperref[lemma-b2]{Lemma B.2}. 


\item  The entries of $\Theta_{t}$ are  $O\left(W^{-1}\eta_{t}^{-\frac{1}{2}}\right)$ per \hyperref[lemma-a2]{Lemma A.2}
    
\item Per the covariance normalization, $S_{ab}\lesssim W^{-1}$
\item $\sum_{y}|B_{xy}|$, $\sum_{y}|S_{xy}|$ have size $O(1)$ per \hyperref[lemma-a4]{Lemma A.4}.

\item $\sum_{x}\left[\{\text{Id} + (t-s)\Theta_t\}S^{1/2}\right]_{ax}\lesssim \eta_t^{-1}\eta_s$ per \hyperref[cor-a3]{Corollary A.3}

\item $(t-s)\lesssim \eta_s$ per \hyperref[lemma-a12]{Lemma A.1.2}. 
\end{enumerate}
\end{boxt}

\newpage
\section{Bounds on Flow Terms}
\label{sec:drift}
In this section, our goal is to control the three stopped terms of the fluctuations between the $T_{xy}$ and $\Theta_{xy}$ matrices. The squared $\mathcal{E}_t^{S, \sto}$ and martingale term $\mathcal{E}_t^{M, \sto}$ of the flow SDE can be handled using the resolvent bounds we stated in the last section, along with standard martingale calculus. However, the drift term $\mathcal{E}^{D, \sto}_t$ requires extra care. For it, we will follow Dubova and Yang's derivation closely, which employs the graphical perturbation methods developed by Bourgade, Yau and Yin, in their self-energy renormalization paper \cite{38}. More specifically, we will use Definition 17  \cite{bandSDE} for the diagrammatic notation of a standard oriented graph:
\vspace{0.255 cm}
\begin{boxt}{Definition 2: (Diagrammatic notation)}\label{def2}
\begin{itemize}
\item Each vertex will be assigned a label, corresponding to a matrix index $\alpha \in \{1, \cdots, N\}$. 
\item Each blue-colored loop \textcolor{blue}{$\circlearrowright$} at $\alpha$, regardless of direction, represents a term $G_s(z)_{\alpha\alpha} - m(z)$;
\item A solid blue edge  $\alpha\textcolor{blue}{\longrightarrow}\beta$ represents a factor of $G_s(z)_{\alpha\beta}$, whereas a red $\alpha\textcolor{red}{\longrightarrow}\beta$ one is $\overline{G_s(z)_{\alpha\beta}}$;
\item A black wavy edge $\alpha\leftrightsquigarrow \beta$ represents $S_{\alpha\beta}$, whereas $\alpha\textcolor{blue}{\leftrightsquigarrow} \beta$ is $B_{\alpha\beta} = (I-sm(z)^2S)^{-1}$; 
\item A double edge $\alpha =\!=\beta$ represents $\{\text{Id}+(t-s)\Theta_t\}S^{1/2}$;
\end{itemize}
\end{boxt}
Having defined the diagrammatic notation, we can now present the main operation that will allow for the bounding of the drift term, namely Lemma 16 \cite{bandSDE}, which lists the loop expansion (Lemma 3.5 \cite{38}) and the regular vertex expansion (Lemma 3.14 \cite{38}): 

\subsection*{Loop Expansion}
Consider a differentiable function $f: \mathbb{C}^{N^2}\rightarrow\mathbb{C}$, which represents the remaining subgraph per our diagrammatic notation above. For fixed time $s\in [0, 1]$, we have the following identity of subgraphs, not accounting for the renormalization term: 
\begin{equation*}
(G_s(z))_{vv} - m) f(G_s(z)) = {sm} \sum_{\alpha, \beta=1}^N B_{v\alpha} S_{\alpha\beta} \big((G_s(z))_{\alpha\alpha} - m\big)\big((G_s(z))_{\beta\beta} - m\big) f(G_s(z))
\end{equation*}
\begin{equation*}
 - sm \sum_{\alpha, \beta=1}^N B_{v\alpha} S_{\alpha\beta} {(G_s(z))_{\beta\alpha}} \partial_{H_{s,\beta\alpha}} f(G_s(z)) = {sm}\mathcal{G}_{1, s}(z)_{ab}+sm\mathcal{G}_{2, s}(z)_{ab} + sm\mathcal{G}_{3, s}(z)_{ab} + sm\mathcal{G}_{4, s}(z)_{ab}.
\end{equation*}
Here, $\mathcal{G}_{1, s}(z)_{ab}$ is equivalent to the first term, whereas subgraphs $\mathcal{G}_{2, s}(z)_{ab}$ through $\mathcal{G}_{4, s}(z)_{ab}$ are the result of expanding the partial derivative $\partial_{H_{s, \alpha\beta}}f(G_s(z))$, using the fact that: $\partial_{H_{\beta\alpha}} G_{xy} = -G_{x\beta}G_{\alpha y}$ and $\partial_{H_{\beta\alpha}}\overline{G}_{xy} = -\overline{G}_{x\alpha}\overline{G}_{\beta y}$.  Similarly, the regular vertex expansion for a connected pair of edges is:
\subsection*{Regular Vertex Expansion}
For a differentiable function $f: \mathbb{C}^{N^2}\rightarrow\mathbb{C}$ and fixed time $s\in [0, 1]$, we have the following identity, not accounting for the renormalization term: \begin{equation*}
G_{xu}G_{uy}f(G) = mB_{uy}G_{xy} f(G) + sm\sum_{\gamma\delta=1}^N B_{u\gamma} S_{\gamma\delta} G_{x\gamma}(G_{\delta\delta}-m)G_{\gamma y} f(G)
\end{equation*}
$$+sm\sum_{\gamma\delta=1}^N B_{u\gamma} S_{\gamma\delta} G_{x\delta}(G_{\gamma\gamma}-m)G_{\delta y}f(G)-sm\sum_{\gamma\beta=1}^N B_{u\gamma}S_{\gamma\delta}G_{x\gamma}G_{\delta y}\partial_{H_{\delta \gamma}}(f(G))$$
$$= m\mathcal{G}_{i0, s}+sm\mathcal{G}_{i1,s}+sm\mathcal{G}_{i2, s}-sm\mathcal{G}_{i3,s}^{(1)}-sm\mathcal{G}_{i3,s}^{(2)}-sm\mathcal{G}_{i3,s}^{(3)},$$
where $\mathcal{G}_{i0, s}$, $\mathcal{G}_{i1, s}$ and $\mathcal{G}_{i2, s}$ represent the first three terms, respectively, and $\mathcal{G}_{i3, s}^{(i)}$ are the resulting elements from applying 
the resolvent derivatives w.r.t. to $H_{\beta\alpha}$. To be clear, the exact number of the latter do not follow from the lemma itself and are rather a feature of our particular diagrammatic structure. The explicit results will be verified when we apply the lemma to each respective term.

\subsection{Estimate for $\mathcal{E}_t^{D, \text{stop}}(z)$} 
The goal for this subsection is to bound the drift term, which we can expand and rewrite as follows: 
$$\mathcal{E}_t^{D, \sto}(z)_{ab}=\int_0^{t\wedge \tau_\sto}\{\text{Id}+(t-s)\Theta_t\}S^{\frac{1}{2}}\Omega_s(z)S^{\frac{1}{2}}\{\text{Id}+(t-s)\Theta_t\}ds$$
$$=\int_0^{t\wedge \tau_\sto} \{\text{Id}+(t-s)\Theta_t\}\left\{
\overline{G_s(z)}_{xy} G_s(z)_{xu} S_{uv} 
\left[ G_s(z)_{vv} - m(z) \right] G_s(z)_{uy}
\right\}\{\text{Id}+(t-s)\Theta_t\}ds+
$$
$$
\int_0^{t\wedge \tau_\sto} \{\text{Id}+(t-s)\Theta_t\}\left\{
 G_s(z)_{xy} \overline{G_s(z)}_{xu} S_{uv} 
\left[ \overline{G_s(z)}_{vv} - m(z) \right] \overline{G_s(z)}_{uy} 
\right\} \{\text{Id}+(t-s)\Theta_t\}ds = 
$$
Observe that we can represent the first term with the diagrammatic notation defined earlier, whereas its complex conjugate has the exact same structure, with only the colors of the straight edges being flipped. Since this doesn't change the underlying estimates we have from the previous section, WLOG we can pick the first term and write it out as follows:

\begin{center}\label{graph}

\begin{tikzpicture}[
    thick, % Line thickness
    wavy/.style={decorate, decoration={snake, amplitude=2pt, segment length=5pt}},
    ]
\node (a) at (0,0) {a};
\node (x) at (2, 0) {x};
\node (u) at (4, 0) {u};
\node (v) at (4, 1) {v};
\node (y) at (6, 0) {y};
\node (b) at (8, 0) {b};

\draw[thick, double] (a) -- (x);
\draw[blue, ->] (x) to (u);
\draw[wavy] (u) -- (v);
\draw[blue, ->] (u) -- (y);
\draw[thick, double] (y) -- (b);
\draw[red, ->] (x) to[bend right=47	] (y);
\draw[blue, ->] (v) to[out=200, in=150, loop] (v);
\end{tikzpicture}
\begin{equation*}\left[\{\text{Id} + (t-s)\Theta_t\}S^{1/2}\right]_{ax}\overline{G_s(z)}_{xy} G_s(z)_{xu} S_{uv} 
\left[ G_s(z)_{vv} - m(z) \right] G_s(z)_{uy}\left[\{\text{Id} + (t-s)\Theta_t\}S^{1/2}\right]_{ax}  \tag{6}
\end{equation*}
\end{center}
By applying the loop expansion at edge $v$, we get the following four subgraphs: 
\begin{center}
\begin{tikzpicture}[
    thick, % Line thickness
    wavy/.style={decorate, decoration={snake, amplitude=2pt, segment length=5pt}},
    ]
\node (a) at (0,0) {a};
\node (x) at (2, 0) {x};
\node (u) at (4, 0) {u};
\node (v) at (4, 1) {v};
\node (y) at (6, 0) {y};
\node (b) at (8, 0) {b};
\node (alpha) at (4, 2) {$\alpha$};
\node (beta) at (6, 2) {$\beta$};

\draw[thick, double] (a) -- (x);
\draw[blue, ->] (x) to (u);
\draw[wavy] (u) -- (v);
\draw[blue, ->] (u) -- (y);
\draw[thick, double] (y) -- (b);
\draw[red, ->] (x) to[bend right=47	] (y);
\draw[blue, ->] (alpha) to[out=200, in=150, loop] (alpha);
\draw[blue, ->] (beta) to[out=0, in=40, loop] (beta);
\draw[wavy, blue] (v) to (alpha);
\draw[wavy] (alpha) to (beta);
\end{tikzpicture}
\begin{equation*}\scalebox{0.95}{$
\mathcal{G}_{1, s}(z)_{ab} = \left[\{\text{Id} + (t-s)\Theta_t\}S^{1/2}\right]_{ax}\overline{G_{xy}}G_{xu}S_{uv}B_{v\alpha}(G_{\alpha\alpha}-m)S_{\alpha\beta}(G_{\beta\beta}-m)G_{uy}\left[\{\text{Id} + (t-s)\Theta_t\}S^{1/2}\right]_{yb}$}
\end{equation*}

\begin{tikzpicture}[
    thick, % Line thickness
    wavy/.style={decorate, decoration={snake, amplitude=2pt, segment length=5pt}},
    ]
\node (a) at (0,0) {a};
\node (x) at (2, 0) {x};
\node (u) at (4, 0) {u};
\node (v) at (4, 1) {v};
\node (y) at (6, 0) {y};
\node (b) at (8, 0) {b};
\node (alpha) at (4, 2) {$\alpha$};
\node (beta) at (6, 2) {$\beta$};

\draw[thick, double] (a) -- (x);
\draw[blue, ->] (x) to (u);
\draw[wavy] (u) -- (v);
\draw[blue, ->] (u) -- (y);
\draw[thick, double] (y) -- (b);
\draw[wavy, blue] (v) to (alpha);
\draw[wavy] (alpha) to (beta);
\draw[red, ->] (beta) to (y);
\draw[red, ->] (x) to[bend left=20] (alpha);
\draw[blue, ->] (beta) to[bend left=30] (alpha);
\end{tikzpicture}
\begin{equation*}
\mathcal{G}_{2, s}(z)_{ab} = \left[\{\text{Id} + (t-s)\Theta_t\}S^{1/2}\right]_{ax}\overline{G_{x\alpha}} \overline{G_{\beta y}}G_{xu} S_{uv} B_{v\alpha}S_{\alpha\beta}G_{\beta\alpha}G_{uy}\left[\{\text{Id} + (t-s)\Theta_t\}S^{1/2}\right]_{yb}
\end{equation*}
\begin{tikzpicture}[
    thick, % Line thickness
    wavy/.style={decorate, decoration={snake, amplitude=2pt, segment length=5pt}},
    ]
\node (a) at (0,0) {a};
\node (x) at (2, 0) {x};
\node (u) at (4, 0) {u};
\node (v) at (4, 1) {v};
\node (y) at (6, 0) {y};
\node (b) at (8, 0) {b};
\node (alpha) at (4, 2) {$\alpha$};
\node (beta) at (6, 2) {$\beta$};

\draw[thick, double] (a) -- (x);
\draw[wavy] (u) -- (v);
\draw[blue, ->] (u) -- (y);
\draw[thick, double] (y) -- (b);
\draw[wavy, blue] (v) to (alpha);
\draw[wavy] (alpha) to (beta);
\draw[blue, <-] (alpha) to[bend right=40] (beta);
\draw[blue, ->] (x) to[bend left=60] (beta);
\draw[blue, ->] (alpha) to[bend right=40] (u);
\draw[red, ->] (x) to [bend right=20] (y);
\end{tikzpicture}
\begin{equation*}
\mathcal{G}_{3, s}(z)_{ab} = \left[\{\text{Id} + (t-s)\Theta_t\}S^{1/2}\right]_{ax}\overline{G_{xy}} G_{x\beta} G_{\alpha u} S_{uv} B_{v\alpha}S_{\alpha\beta}G_{\beta\alpha}G_{uy}\left[\{\text{Id} + (t-s)\Theta_t\}S^{1/2}\right]_{yb}
\end{equation*}
\end{center}

\begin{center}
\begin{tikzpicture}[
    thick, % Line thickness
    wavy/.style={decorate, decoration={snake, amplitude=2pt, segment length=5pt}},
    ]
\node (a) at (0,0) {a};
\node (x) at (2, 0) {x};
\node (u) at (4, 0) {u};
\node (v) at (4, 1) {v};
\node (y) at (6, 0) {y};
\node (b) at (8, 0) {b};
\node (alpha) at (4, 2) {$\alpha$};
\node (beta) at (6, 2) {$\beta$};

\draw[thick, double] (a) -- (x);
\draw[blue, ->] (x) to (u);
\draw[wavy] (u) -- (v);
\draw[thick, double] (y) -- (b);
\draw[wavy, blue] (v) to (alpha);
\draw[wavy] (alpha) to (beta);
\draw[blue, ->] (u) to (beta);
\draw[blue, ->] (alpha) to (y);
\draw[red, ->] (x) to [bend right=20] (y);
\draw[blue, <-] (alpha) to [bend left=40] (beta);
\end{tikzpicture}
\begin{equation*}
\mathcal{G}_{4, s}(z)_{ab} = \left[\{\text{Id} + (t-s)\Theta_t\}S^{1/2}\right]_{ax}\overline{G_{xy}} G_{xu}  S_{uv} B_{v\alpha}S_{\alpha\beta}G_{\beta\alpha}G_{u\beta}G_{\alpha y}\left[\{\text{Id} + (t-s)\Theta_t\}S^{1/2}\right]_{yb}
\end{equation*}
\end{center}
Observe that in all cases we have a vertex $u$ that has one incoming and one outgoing blue edge, implying that we can again apply the regular vertex expansion w.r.t to $u$. For both $\mathcal{G}_{1, s}$ and $\mathcal{G}_{2, s}$ the incoming and outgoing edges are the same, starting and ending at $x$ and $y$, respectively. For $i=3$ and $i=4$ we have a different incoming or outgoing edge, so the corresponding expansion will require slight reformulation. Accounting for this fact, we get the following identities: 
\begin{itemize}
\item For $\mathcal{G}_{1, s}$ and $\mathcal{G}_{2, s}$: \begin{equation*}
G_{xu}G_{uy}f(G) = mB_{uy}G_{xy} f(G) + sm\sum_{\gamma\delta=1}^N B_{u\gamma} S_{\gamma\delta} G_{x\gamma}(G_{\delta\delta}-m)G_{\gamma y} f(G)
\end{equation*}
$$+sm\sum_{\gamma\delta=1}^N B_{u\gamma} S_{\gamma\delta} G_{x\delta}(G_{\gamma\gamma}-m)G_{\delta y}f(G)-sm\sum_{\gamma\beta=1}^N B_{u\gamma}S_{\gamma\delta}G_{x\gamma}G_{\delta y}\partial_{H_{\delta \gamma}}(f(G))$$
\item For $\mathcal{G}_{3, s}$:
\begin{equation*}
G_{\alpha u}G_{uy}f(G) = mB_{uy}G_{\alpha y} f(G) + sm\sum_{\gamma\delta=1}^N B_{u\gamma} S_{\gamma\delta} G_{\alpha\gamma}(G_{\delta\delta}-m)G_{\gamma y} f(G)
\end{equation*}
$$+sm\sum_{\gamma\delta=1}^N B_{u\gamma} S_{\gamma\delta} G_{\alpha\delta}(G_{\gamma\gamma}-m)G_{\delta y}f(G)-sm\sum_{\gamma\delta=1}^N B_{u\gamma}S_{\gamma\delta}G_{\alpha\gamma}G_{\delta y}\partial_{H_{\delta \gamma}}(f(G))$$
\item For $\mathcal{G}_{4, s}$:
\begin{equation*}
G_{xu}G_{u\beta}f(G) = mB_{u\beta}G_{x\beta} f(G) + sm\sum_{\gamma\delta=1}^N B_{u\gamma} S_{\gamma\delta} G_{x\gamma}(G_{\delta\delta}-m)G_{\gamma \beta} f(G)
\end{equation*}
$$+sm\sum_{\gamma\delta=1}^N B_{u\gamma} S_{\gamma\delta} G_{x\delta}(G_{\gamma\gamma}-m)G_{\delta \beta}f(G)-sm\sum_{\gamma\delta=1}^N B_{u\gamma}S_{\gamma\delta}G_{x\gamma}G_{\delta \beta}\partial_{H_{\delta \gamma}}(f(G))$$
\end{itemize}
As we noted earlier, each of the former identities can be represented as the sum of independent subraphs, such that for $i=1, 2, 3, 4$, we have:
$$\mathcal{G}_{i, s} = m\mathcal{G}_{i0, s}+sm\mathcal{G}_{i1,s}+sm\mathcal{G}_{i2, s}-sm\mathcal{G}_{i3,s}^{(1)}-sm\mathcal{G}_{i3,s}^{(2)}-sm\mathcal{G}_{i3,s}^{(3)},$$
We will begin by focusing our attention on the first three terms for a fixed $i=1, 2, 3, 4$, namely $\mathcal{G}_{ij}$ for $j=0, 1, 2$. Using the estimates we defined in Section 5, we can simplify each of the resulting subgraph, starting with $j=0$:\newpage
\subsubsection*{Size estimates for $\mathcal{G}_{10, s}$}
Following diagrammatic notation, each of the $B_{uy_i}G_{x_iy_i}f(G)$ terms becomes:
\begin{center}
\begin{tikzpicture}[
    thick, % Line thickness
    wavy/.style={decorate, decoration={snake, amplitude=2pt, segment length=5pt}},
    ]
\node (a) at (0,0) {a};
\node (x) at (2, 0) {x};
\node (u) at (4, 0) {u};
\node (v) at (4, 1) {v};
\node (y) at (6, 0) {y};
\node (b) at (8, 0) {b};
\node (alpha) at (4, 2) {$\alpha$};
\node (beta) at (6, 2) {$\beta$};

\draw[thick, double] (a) -- (x);
\draw[wavy] (u) -- (v);
\draw[wavy, blue] (u) -- (y);
\draw[thick, double] (y) -- (b);
\draw[red, ->] (x) to[bend right=15	] (y);
\draw[blue, ->](x) to [bend right=35] (y);
\draw[blue, ->] (alpha) to[out=200, in=150, loop] (alpha);
\draw[blue, ->] (beta) to[out=0, in=40, loop] (beta);
\draw[wavy, blue] (v) to (alpha);
\draw[wavy] (alpha) to (beta);
\end{tikzpicture}
\begin{equation*}
\mathcal{G}_{10, s}(z)_{ab} = \left[\{\text{Id} + (t-s)\Theta_t\}S^{1/2}\right]_{ax}\overline{G_{xy}}S_{uv}B_{v\alpha}(G_{\alpha\alpha}-m)S_{\alpha\beta}(G_{\beta\beta}-m)B_{uy}\left[\{\text{Id} + (t-s)\Theta_t\}S^{1/2}\right]_{yb}
\end{equation*}
\end{center}
We can apply estimate $(E2)$ to the two loops around $\alpha$ and $\beta$, as well as estimate $(E1)$ to the red edge connecting $x$ and $y$. From this, we receive three multiplicative terms $O(W^{\delta_\sto/20}W^{-1/2}\eta_s^{-1/4})$, i.e combined it yields a term $O(W^{3\delta_\sto/20}W^{-3/2}\eta_s^{-3/4})$. By estimate $(E5)$, we can sum all the wavy lines out in order $\beta, \alpha, v,$ and finally $u$ for a $O(1)$ term, keeping our earlier term. We can stop here by leaving the connected component that is a double line from $a$ to $x$, a blue line from $x$ to $y$ and another double line from $y$ to $b$, i.e: $\mathcal{G}_{10, s}\lesssim W^{3\delta_\sto/20}W^{-3/2}\eta_s^{-3/4}\times $
\begin{tikzpicture}[
    thick, % Line thickness
    wavy/.style={decorate, decoration={snake, amplitude=2pt, segment length=5pt}},
    ]
\node (a) at (0,0) {a};
\node (x) at (1, 0) {x};
\node (y) at (2, 0) {y};
\node (b) at (3, 0) {b};

\draw[thick, double] (a) -- (x);
\draw[thick, double] (y) -- (b);
\draw[blue, ->](x) to (y);

\end{tikzpicture}

\subsubsection*{Size estimates for $\mathcal{G}_{20, s}$}

\begin{center}
\begin{tikzpicture}[
    thick, % Line thickness
    wavy/.style={decorate, decoration={snake, amplitude=2pt, segment length=5pt}},
    ]
\node (a) at (0,0) {a};
\node (x) at (2, 0) {x};
\node (u) at (4, 0) {u};
\node (v) at (4, 1) {v};
\node (y) at (6, 0) {y};
\node (b) at (8, 0) {b};
\node (alpha) at (4, 2) {$\alpha$};
\node (beta) at (6, 2) {$\beta$};

\draw[thick, double] (a) -- (x);
\draw[wavy] (u) -- (v);
\draw[blue, wavy] (u) -- (y);
\draw[blue, ->] (x) to [bend right=20] (y);
\draw[thick, double] (y) -- (b);
\draw[wavy, blue] (v) to (alpha);
\draw[wavy] (alpha) to (beta);
\draw[red, ->] (beta) to (y);
\draw[red, ->] (x) to[bend left=20] (alpha);
\draw[blue, ->] (beta) to[bend left=30] (alpha);
\end{tikzpicture}
\begin{equation*}
\mathcal{G}_{20, s}(z)_{ab} = \left[\{\text{Id} + (t-s)\Theta_t\}S^{1/2}\right]_{ax}\overline{G_{x\alpha}}G_{xy} \overline{G_{\beta y}} S_{uv} B_{v\alpha}S_{\alpha\beta}G_{\beta\alpha}B_{uy}\left[\{\text{Id} + (t-s)\Theta_t\}S^{1/2}\right]_{yb}
\end{equation*}
\end{center}
Repeating the same process as before, we bound all three red lines using estimate $(E2)$ and sum over the wavy edges in the same order $(\beta, \alpha, v$, and $u$), s.t. we get the same bound as above:\\
$\mathcal{G}_{20, s}\lesssim W^{3\delta_\sto/20}W^{-3/2}\eta_s^{-3/4}\times $
\begin{tikzpicture}[
    thick, % Line thickness
    wavy/.style={decorate, decoration={snake, amplitude=2pt, segment length=5pt}},
    ]
\node (a) at (0,0) {a};
\node (x) at (1, 0) {x};
\node (y) at (2, 0) {y};
\node (b) at (3, 0) {b};

\draw[thick, double] (a) -- (x);
\draw[thick, double] (y) -- (b);
\draw[blue, ->](x) to (y);

\end{tikzpicture}

\subsubsection*{Size estimates for $\mathcal{G}_{30, s}$}


\begin{center}
\begin{tikzpicture}[
    thick, % Line thickness
    wavy/.style={decorate, decoration={snake, amplitude=2pt, segment length=5pt}},
    ]
\node (a) at (0,0) {a};
\node (x) at (2, 0) {x};
\node (u) at (4, 0) {u};
\node (v) at (4, 1) {v};
\node (y) at (6, 0) {y};
\node (b) at (8, 0) {b};
\node (alpha) at (4, 2) {$\alpha$};
\node (beta) at (6, 2) {$\beta$};

\draw[thick, double] (a) -- (x);
\draw[wavy] (u) -- (v);
\draw[blue, wavy] (u) to (y);
\draw[blue, ->] (alpha) to (y);
\draw[thick, double] (y) -- (b);
\draw[wavy, blue] (v) to (alpha);
\draw[wavy] (alpha) to (beta);
\draw[blue, <-] (alpha) to[bend right=30] (beta);
\draw[blue, ->] (x) to[bend left=60] (beta);
\draw[red, ->] (x) to [bend right=20] (y);
\end{tikzpicture}
\begin{equation*}
\mathcal{G}_{30, s}(z)_{ab} = \left[\{\text{Id} + (t-s)\Theta_t\}S^{1/2}\right]_{ax}\overline{G_{xy}} G_{x\beta}S_{uv} B_{v\alpha}  S_{\alpha\beta}G_{\beta\alpha}G_{uy} B_{uy}\left[\{\text{Id} + (t-s)\Theta_t\}S^{1/2}\right]_{yb}
\end{equation*}
\end{center}
Here, we eliminate first the edges  $\beta\textcolor{red}{\rightarrow}\alpha$ and $\alpha\textcolor{blue}{\rightarrow} y$ using estimate $(E2)$, giving us a $O(W^{\delta_\sto/10}W^{-1}\eta_s^{-1/2})$ term. Due to the edge from $x\textcolor{blue}{\rightarrow}\beta$, we are cannot sum $\alpha\leadsto\beta$ over $S_{ab}$ entries because it serves as a multiplicative term to the blue edge $x\textcolor{blue}{\rightarrow}\beta$. For this reason, we will apply the individual $S_{ab}\lesssim W^{-1}$ bound, s.t now summing over $x\textcolor{blue}{\rightarrow}\beta$ with estimate $(E1)$ gives us combined $O(W^{\delta_\sto/20}W^{\frac{1}{2}+\epsilon}\eta_s^{-3/4})$. Finally, by summing over the remaining wavy edges, starting with $\alpha, v$ and then $u$, due to the $O(1)$ term (estimate $(E5)$), we get an estimate of the form: $\mathcal{G}_{30, s}\lesssim W^{3\delta_\sto/20}W^{-\frac{1}{2}+\epsilon}\eta_s^{-5/4}\times $
\begin{tikzpicture}[
    thick, % Line thickness
    wavy/.style={decorate, decoration={snake, amplitude=2pt, segment length=5pt}},
    ]
\node (a) at (0,0) {a};
\node (x) at (1, 0) {x};
\node (y) at (2, 0) {y};
\node (b) at (3, 0) {b};

\draw[thick, double] (a) -- (x);
\draw[thick, double] (y) -- (b);
\draw[red, ->](x) to (y);

\end{tikzpicture}


\subsubsection*{Size estimates for $\mathcal{G}_{40, s}$}

\begin{center}
\begin{tikzpicture}[
    thick, % Line thickness
    wavy/.style={decorate, decoration={snake, amplitude=2pt, segment length=5pt}},
    ]
\node (a) at (0,0) {a};
\node (x) at (2, 0) {x};
\node (u) at (4, 0) {u};
\node (v) at (4, 1) {v};
\node (y) at (6, 0) {y};
\node (b) at (8, 0) {b};
\node (alpha) at (4, 2) {$\alpha$};
\node (beta) at (6, 2) {$\beta$};

\draw[thick, double] (a) -- (x);
\draw[wavy] (u) -- (v);
\draw[thick, double] (y) -- (b);
\draw[wavy, blue] (v) to (alpha);
\draw[wavy] (alpha) to (beta);
\draw[blue, wavy] (u) to (beta);
\draw[blue, ->] (alpha) to (y);
\draw[red, ->] (x) to [bend right=20] (y);
\draw[blue, <-] (alpha) to [bend left=15] (beta);
\draw[blue, ->] (x) to[bend left=60] (beta);
\end{tikzpicture}
\begin{equation*}
\mathcal{G}_{4, s}(z)_{ab} = \left[\{\text{Id} + (t-s)\Theta_t\}S^{1/2}\right]_{ax}\overline{G_{xy}} G_{x\beta}  S_{uv} B_{v\alpha}S_{\alpha\beta}G_{\beta\alpha}B_{u\beta}G_{\alpha y}\left[\{\text{Id} + (t-s)\Theta_t\}S^{1/2}\right]_{yb}
\end{equation*}
\end{center}
Here, we bound first the inner edges $\alpha\textcolor{blue}{\rightarrow} y$ and $\beta\textcolor{red}{\rightarrow}\alpha$, we get by estimate $(E2)$ a combined term $\lesssim O(W^{\delta_\sto/10}W^{-1}\eta_{s}^{-1/2})$. Again, by the same logic as above, we have to account for the individual term $\alpha\leadsto\beta$ instead of summing over it, retaining a $O(W^{-1})$ term. By applying estimate $(E1)$ for the edge $x\textcolor{blue}{\rightarrow}\beta$ and summing over the remaining wavy lines, we get the same bound as above, namely: 
$\mathcal{G}_{40, s}\lesssim  W^{3\delta_\sto/20}W^{-\frac{1}{2}+\epsilon}\eta_s^{-5/4}\times $
\begin{tikzpicture}[
    thick, % Line thickness
    wavy/.style={decorate, decoration={snake, amplitude=2pt, segment length=5pt}},
    ]
\node (a) at (0,0) {a};
\node (x) at (1, 0) {x};
\node (y) at (2, 0) {y};
\node (b) at (3, 0) {b};

\draw[thick, double] (a) -- (x);
\draw[thick, double] (y) -- (b);
\draw[red, ->](x) to (y);

\end{tikzpicture}

\newpage
\subsubsection*{Size estimates for $\max_{i=1, \cdots, 4}\mathcal{G}_{40, s}$}

In order to combine all the latter estimates, we can represent with a pink arrow either blue or red:\begin{center}

$\max_{i=1, \cdots, 4}\left\vert\mathcal{G}_{40, s}\right\vert\lesssim  W^{3\delta_\sto/20} W^{-\frac{3}{2}+\epsilon}\eta_s^{-5/4}\times $
\begin{tikzpicture}[
    thick, % Line thickness
    wavy/.style={decorate, decoration={snake, amplitude=2pt, segment length=5pt}},
    ]
\node (a) at (0,0) {a};
\node (x) at (1, 0) {x};
\node (y) at (2, 0) {y};
\node (b) at (3, 0) {b};

\draw[thick, double] (a) -- (x);
\draw[thick, double] (y) -- (b);
\draw[pink, ->](x) to (y);

\end{tikzpicture}
\end{center}
We need to consider two subcases by splitting the double line $\left[\{\text{Id} + (t-s)\Theta_t\}S^{1/2}\right]_{yb}$:
\begin{itemize}
\item $(t-s)\left[\Theta_t\right]_{yb} \lesssim (t-s)W^{-1}\eta_t^{-1/2}\lesssim W^{-1}\eta_t^{-1/2}\eta_s$ by estimate $(E3)$ and $(E7)$. From summing over $y$ with estimate $(E1)$, we get $ O(W^{3\delta_\sto/20}W^{\frac{1}{2}+\epsilon}\eta_s^{-3/4})$. By summing over $x$ with estimate $(E6)$, we have $O(\eta_t^{-1}\eta_s)$, making the combined term $O((t-s) W^{\delta_\sto/20}W^{-\frac{1}{2}+\epsilon}\eta_t^{-3/2}\eta_s^{1/4})$
\item $\text{Id}_{yb}$ is of course $O(1)$, hence we don't have a varying term. This means that we can bound the edge ending at $y$ with estimate $(E2)$, getting a $O( W^{2\delta_\sto/20}W^{-1/2}\eta_s^{-1/4})$ term. Combining it with $O(\eta_t^{-1}\eta_s)$ from estimate $(E6)$, we have $O( W^{\delta_\sto/20}W^{-1/2}\eta_t^{-1}\eta_s^{3/4})$. With this, the total bound is: 
\end{itemize}

\begin{center}
\begin{tikzpicture}[
    thick, % Line thickness
    wavy/.style={decorate, decoration={snake, amplitude=2pt, segment length=5pt}},
    ]
\node (a) at (0,0) {a};
\node (x) at (1, 0) {x};
\node (y) at (2, 0) {y};
\node (b) at (3, 0) {b};

\draw[thick, double] (a) -- (x);
\draw[thick, double] (y) -- (b);
\draw[pink, ->](x) to (y);

\end{tikzpicture}
$\lesssim  (t-s) W^{\delta_\sto/20}W^{-\frac{1}{2}+\epsilon}\eta_t^{-3/2}\eta_s^{1/4} + W^{\delta_\sto/20}W^{-1/2}\eta_t^{-1}\eta_s^{3/4}\lesssim$
\end{center}
\begin{equation*}\lesssim W^{\delta_\sto/20}W^{-\frac{1}{2}+\epsilon}\eta_t^{-3/2}\eta_s^{5/4} + W^{\delta_\sto/20}W^{-1/2}\eta_t^{-1}\eta_s^{3/4}\tag{6.1}\end{equation*}
This means that if we combine the diagram bound with the result above, we get: 
\begin{equation*}\max_{i=1, 2, 3, 4}|\mathcal{G}_{i0, s}|\lesssim W^{\delta_\sto/5}W^{-2+2\epsilon}\eta_t^{-3/2} + W^{\delta_\sto/5}W^{-2+2\epsilon}\eta_t^{-1}\eta_s^{1/2}\end{equation*}
Lastly, we need to integrate over $s\in [0, t]$ by using a change of variable $\sigma =\eta_s$, s.t. the bounds of the integral are $\sigma\in [\eta_t, \eta_0]$:
$$\max_{i=1, 2, 3, 4} \int_0^{t\wedge \tau_\sto}|\mathcal{G}_{i0, s}|ds\lesssim W^{\frac{\delta_\sto}{5}+2\epsilon} W^{-2}\eta_t^{-3/2} \int_{\eta_t}^{\eta_0}d\sigma + W^{\frac{\delta_\sto}{5}+2\epsilon}W^{-2}\eta_t^{-1}\int_{\eta_t}^{\eta_0}\sigma^{-1/2}d\sigma $$
\begin{equation*}\max_{i=1, 2, 3, 4} \int_0^{t\wedge \tau_\sto}|\mathcal{G}_{i0, s}|ds\lesssim W^{\frac{\delta_\sto}{5}+2\epsilon} W^{-2}\eta_t^{-\frac{3}{2}}\lesssim W^{-\frac{3}{4}}\eta_t^{-1}W^{-1}\eta_t^{-\frac{1}{2}},\tag{6.2}\end{equation*}
since $\frac{\delta_\sto}{5}+2\epsilon$ is arbitrarily small. 
\newpage
\subsubsection*{Size estimates for $\mathcal{G}_{11, s}$}
In order to represent $B_{u\gamma} S_{\gamma\delta} G_{x_i\gamma}(G_{\delta\delta}-m)G_{\gamma y_i} f(G)$, we need to add the two new vertices $\gamma$ and $\delta$:
\begin{center}
\begin{tikzpicture}[
    thick, % Line thickness
    wavy/.style={decorate, decoration={snake, amplitude=2pt, segment length=5pt}},
    ]
\node (a) at (0,0) {a};
\node (x) at (2, 0) {x};
\node (gamma) at (4, 0) {$\gamma$};
\node (u) at (4, 1) {u};
\node (v) at (5, 1) {v};
\node (y) at (6, 0) {y};
\node (b) at (8, 0) {b};
\node (alpha) at (4, 2) {$\alpha$};
\node (beta) at (6, 2) {$\beta$};
\node (delta) at (3, 1) {$\delta$};

\draw[thick, double] (a) -- (x);
\draw[wavy, blue] (u) -- (gamma);
\draw[thick, double] (y) -- (b);
\draw[red, ->] (x) to[bend right=15	] (y);
\draw[blue, ->] (alpha) to[out=200, in=150, loop] (alpha);
\draw[blue, ->] (beta) to[out=0, in=40, loop] (beta);
\draw[wavy, blue] (v) to (alpha);
\draw[wavy] (alpha) to (beta);
\draw[blue, ->] (x) to (gamma);
\draw[blue, ->] (gamma) to (y);
\draw[wavy] (u) to (v);
\draw[wavy](gamma) to (delta);
\draw[blue, ->] (delta) to[out=200, in=150, loop] (delta);
\end{tikzpicture}
\begin{equation*}\scalebox{0.85}{$
\left[\{\text{Id} + (t-s)\Theta_t\}S^{1/2}\right]_{ax}\overline{G_{xy}}B_{u\gamma} S_{\gamma \delta} G_{x\gamma}(G_{\delta\delta}-m)G_{\gamma y}S_{uv}B_{v\alpha}(G_{\alpha\alpha}-m)S_{\alpha\beta}(G_{\beta\beta}-m)B_{uy}\left[\{\text{Id} + (t-s)\Theta_t\}S^{1/2}\right]_{yb}
$}\end{equation*}
\end{center}
\noindent To estimate this graph, we will bound all loops $\textcolor{blue}{\circlearrowright}$ and $x\textcolor{red}{\rightarrow} y$, using estimate $(E2)$. This yields in total $O(W^{\delta_\sto/5}W^{-2}\eta_s^{-1})$. Then, we sum the wavy lines over $\beta, \alpha, v, u, \delta$ in that order, which is $O(1)$ by estimate $(E5)$.The combined bound is:
$\mathcal{G}_{11, s}\lesssim W^{\delta_\sto/5}W^{-2}\eta_s^{-1}\times $
\begin{tikzpicture}[
    thick, % Line thickness
    wavy/.style={decorate, decoration={snake, amplitude=2pt, segment length=5pt}},
    ]
\node (a) at (0,0) {a};
\node (x) at (1, 0) {x};
\node (gamma) at (2, 0) {$\gamma$};
\node (y) at (3, 0) {y};
\node (b) at (4, 0) {b};

\draw[thick, double] (a) -- (x);
\draw[thick, double] (y) -- (b);
\draw[blue, ->](x) to (gamma);
\draw[blue, ->] (gamma) to (y);

\end{tikzpicture}
\subsubsection*{Size estimates for $\mathcal{G}_{21, s}$}
\begin{center}
\begin{tikzpicture}[
    thick, % Line thickness
    wavy/.style={decorate, decoration={snake, amplitude=2pt, segment length=5pt}},
    ]
\node (a) at (0,0) {a};
\node (x) at (2, 0) {x};
\node (gamma) at (4, 0) {$\gamma$};
\node (u) at (4, 1) {u};
\node (v) at (5, 1) {v};
\node (y) at (6, 0) {y};
\node (b) at (8, 0) {b};
\node (alpha) at (4, 2) {$\alpha$};
\node (beta) at (6, 2) {$\beta$};
\node (delta) at (3, 1) {$\delta$};

\draw[thick, double] (a) -- (x);
\draw[wavy] (u) -- (v);
\draw[thick, double] (y) -- (b);
\draw[wavy, blue] (v) to (alpha);
\draw[wavy] (alpha) to (beta);
\draw[red, ->] (beta) to (y);
\draw[red, ->] (x) to[bend left=40] (alpha);
\draw[blue, ->] (beta) to[bend left=20] (alpha);
\draw[wavy, blue] (u) -- (gamma);
\draw[blue, ->] (x) to (gamma);
\draw[blue, ->] (gamma) to (y);
\draw[wavy](gamma) to (delta);
\draw[blue, ->] (delta) to[out=50, in=10, loop] (delta);
\end{tikzpicture}

\begin{equation*}\scalebox{0.9}{$
\mathcal{G}_{21, s}(z)_{ab}= \left[\{\text{Id} + (t-s)\Theta_t\}S^{1/2}\right]_{ax}\overline{G_{x\alpha}} G_{x\gamma}G_{\gamma y}\overline{G_{\beta y}}B_{u\gamma}S_{\gamma\delta}(G_{\delta\delta}-m) S_{uv} B_{v\alpha}S_{\alpha\beta}G_{\beta\alpha}\left[\{\text{Id} + (t-s)\Theta_t\}S^{1/2}\right]_{yb}$}
\end{equation*}
\end{center}

\noindent Here, we will repeat an equivalent procedure as above, bounding the loop at $\delta$ and the edges $x\textcolor{red}{\rightarrow}\alpha$, $\beta\textcolor{blue}{\rightarrow}\alpha$ and $\beta\textcolor{red}{\rightarrow}y$ , getting a bound $O(W^{\delta_\sto/5}W^{-2}\eta_s^{-1})$. Then, we can sum the wavy lines with estimate $(E5)$, getting an equivalent bound as above:
$\mathcal{G}_{21, s}\lesssim W^{\delta_\sto/5}W^{-2}\eta_s^{-1}\times $
\begin{tikzpicture}[
    thick, % Line thickness
    wavy/.style={decorate, decoration={snake, amplitude=2pt, segment length=5pt}},
    ]
\node (a) at (0,0) {a};
\node (x) at (1, 0) {x};
\node (gamma) at (2, 0) {$\gamma$};
\node (y) at (3, 0) {y};
\node (b) at (4, 0) {b};

\draw[thick, double] (a) -- (x);
\draw[thick, double] (y) -- (b);
\draw[blue, ->](x) to (gamma);
\draw[blue, ->] (gamma) to (y);

\end{tikzpicture}

\subsubsection*{Size estimates for $\mathcal{G}_{31, s}$}

\begin{center}
\begin{tikzpicture}[
    thick, % Line thickness
    wavy/.style={decorate, decoration={snake, amplitude=2pt, segment length=5pt}},
    ]
\node (a) at (0,0) {a};
\node (x) at (2, 0) {x};
\node (gamma) at (4, 0) {$\gamma$};
\node (u) at (4, 1) {u};
\node (v) at (5, 1) {v};
\node (y) at (6, 0) {y};
\node (b) at (8, 0) {b};
\node (alpha) at (4, 2) {$\alpha$};
\node (beta) at (2, 2) {$\beta$};
\node (delta) at (3, 1) {$\delta$};

\draw[thick, double] (a) -- (x);
\draw[wavy] (u) -- (v);
\draw[blue, wavy] (gamma) to (u);
\draw[thick, double] (y) -- (b);
\draw[wavy, blue] (v) to (alpha);
\draw[wavy](gamma) to (delta);
\draw[blue, ->] (delta) to[out=250, in=200, loop] (delta);
\draw[wavy] (alpha) to (beta);
\draw[blue, <-] (alpha) to[bend right=20] (beta);
\draw[blue, ->] (x) to (beta);
\draw[red, ->] (x) to [bend right=20] (y);
\draw[blue, ->] (gamma) to (y);
\draw[blue, ->] (alpha) to [bend right=28](gamma);
\end{tikzpicture}
\begin{equation*}\scalebox{0.92}{$
\mathcal{G}_{31, s}(z)_{ab} = \left[\{\text{Id} + (t-s)\Theta_t\}S^{1/2}\right]_{ax}\overline{G_{xy}} G_{x\beta}B_{u\gamma}S_{\gamma \delta}G_{\alpha \gamma}(G_{\delta\delta}-m)G_{\gamma y}S_{uv} B_{v\alpha}  S_{\alpha\beta}G_{\beta\alpha}\left[\{\text{Id} + (t-s)\Theta_t\}S^{1/2}\right]_{yb}$}
\end{equation*}
\end{center}

\noindent Let us start by bounding the loop at $\delta$, as well as the $x\textcolor{red}{\rightarrow}y$, $\beta\textcolor{blue}{\rightarrow}\alpha$ and $\alpha\textcolor{blue}{\rightarrow}\gamma$. In total, this contributes $O(W^{\delta_\sto/5}W^{-2}\eta_s^{-1})$. However, here we don't have the same structure as before, since $\gamma$ no longer has an ingoing edge from $x$. We can abstract away the wavy lines (blue and red) with an orange wavy line, and handle the individual terms, which will be contributing only with $O(1)$ per estimate $(E5)$. As such, the bound is of the form: $\mathcal{G}_{31, s}\lesssim W^{\delta_\sto/5} W^{-2}\eta_s^{-1}\times $
\begin{tikzpicture}[
    thick, % Line thickness
    wavy/.style={decorate, decoration={snake, amplitude=2pt, segment length=5pt}},
    ]
\node (a) at (0,0) {a};
\node (x) at (1, 0) {x};
\node (beta) at (2, 0) {$\beta$};
\node (gamma) at (3,0) {$\gamma$};
\node (y) at (4, 0) {y};
\node (b) at (5, 0) {b};

\draw[thick, double] (a) -- (x);
\draw[thick, double] (y) -- (b);
\draw[wavy, orange] (beta) to (gamma);
\draw[blue, ->](x) to (beta);
\draw[blue, ->] (gamma) to (y);

\end{tikzpicture}

\subsubsection*{Size estimate for $\mathcal{G}_{41, s}$}
\begin{center}
\begin{tikzpicture}[
    thick, % Line thickness
    wavy/.style={decorate, decoration={snake, amplitude=2pt, segment length=5pt}},
    ]
\node (a) at (0,0) {a};
\node (x) at (2, 0) {x};
\node (gamma) at (4, 0) {$\gamma$};
\node (v) at (3, 1) {v};
\node (u) at(4, 1) {u};
\node (y) at (6, 0) {y};
\node (b) at (8, 0) {b};
\node (alpha) at (4, 2) {$\alpha$};
\node (beta) at (2, 2) {$\beta$};
\node (delta) at (5, 1) {$\delta$};

\draw[thick, double] (a) -- (x);
\draw[wavy] (u) -- (v);
\draw[thick, double] (y) -- (b);
\draw[wavy, blue] (v) to (alpha);
\draw[wavy] (alpha) to (beta);
\draw[blue, ->] (alpha) to[bend left = 20] (y);
\draw[blue, <-] (alpha) to [bend right=20] (beta);
\draw[red, ->] (x) to [bend right=20] (y);
\draw[blue, ->] (gamma) to[bend left=30] (beta);
\draw[wavy, blue] (gamma) to (u);
\draw[blue, ->] (x) to (gamma);
\draw[wavy] (gamma) to (delta);
\draw[blue, ->] (delta) to[out=320, in=280, loop] (delta);

\end{tikzpicture}
\begin{equation*}\scalebox{0.92}{$
\mathcal{G}_{41, s}(z)_{ab} = \left[\{\text{Id} + (t-s)\Theta_t\}S^{1/2}\right]_{ax}\overline{G_{xy}} B_{u\gamma}S_{\gamma \delta}G_{x\gamma}(G_{\delta\delta}-m)G_{\gamma\beta} S_{uv} B_{v\alpha}S_{\alpha\beta}G_{\beta\alpha}G_{\alpha y}\left[\{\text{Id} + (t-s)\Theta_t\}S^{1/2}\right]_{yb}$}
\end{equation*}
\end{center}
\noindent Here the logic is the same as above, namely, after eliminating the loop at $\delta$ and the edges $x\textcolor{red}{\rightarrow} y$, $\gamma\textcolor{blue}{\rightarrow} \beta$, and $\beta\textcolor{blue}{\rightarrow} \alpha$ we get a term $O(W^{\delta_\sto/5}W^{-2}\eta_s^{-1})$. The remaining graph is of similar form as the one before, with the wavy edges being abstracted away with the orange one: $\mathcal{G}_{41, s}\lesssim W^{\delta_\sto/5}W^{-2}\eta_s^{-1}\times $
\begin{tikzpicture}[
    thick, % Line thickness
    wavy/.style={decorate, decoration={snake, amplitude=2pt, segment length=5pt}},
    ]
\node (a) at (0,0) {a};
\node (x) at (1, 0) {x};
\node (gamma) at (2,0) {$\gamma$};
\node (alpha) at (3, 0) {$\alpha$};
\node (y) at (4, 0) {y};
\node (b) at (5, 0) {b};

\draw[thick, double] (a) -- (x);
\draw[thick, double] (y) -- (b);
\draw[wavy, orange] (alpha) to (gamma);
\draw[blue, ->](x) to (gamma);
\draw[blue, ->] (alpha) to (y);

\end{tikzpicture}

\subsubsection*{Size estimate for $\max_{i=1 ,\cdots, 4}\mathcal{G}_{i1, s}$}

As such, all our graphs are bounded by:
 
\begin{center}
$\max_{i=1, \cdots, 4} \mathcal{G}_{i1, s}\lesssim W^{\delta_\sto/5} W^{-2}\eta_s^{-1}\times \left\{\begin{array}{l}
\begin{tikzpicture}[
    thick, % Line thickness
    wavy/.style={decorate, decoration={snake, amplitude=2pt, segment length=5pt}},
    ]
\node (a) at (0,0) {a};
\node (x) at (1, 0) {x};
\node (gamma) at (2, 0) {$\gamma$};
\node (y) at (3, 0) {y};
\node (b) at (4, 0) {b};

\draw[thick, double] (a) -- (x);
\draw[thick, double] (y) -- (b);
\draw[blue, ->](x) to (gamma);
\draw[blue, ->] (gamma) to (y);
\end{tikzpicture}
\\[1cm]
\begin{tikzpicture}[
    thick, % Line thickness
    wavy/.style={decorate, decoration={snake, amplitude=2pt, segment length=5pt}},
    ]
\node (a) at (0,0) {a};
\node (x) at (1, 0) {x};
\node (gamma) at (2,0) {$\gamma$};
\node (alpha) at (3, 0) {$\alpha$};
\node (y) at (4, 0) {y};
\node (b) at (5, 0) {b};

\draw[thick, double] (a) -- (x);
\draw[thick, double] (y) -- (b);
\draw[wavy, orange] (alpha) to (gamma);
\draw[blue, ->](x) to (gamma);
\draw[blue, ->] (alpha) to (y);

\end{tikzpicture}
\end{array}\right.$
\end{center}
Observe that the latter subgraphs are almost identical to the one we calculate for the $\mathcal{G}_{i0, s}$ estimates with the addition of a solid line. This means that we have to multiply our $(6.1)$ bound by an additional term $O(W^{\delta_\sto/20}W^{\frac{1}{2}+\epsilon}\eta_s^{-3/4})$ from estimate $(E1)$ to get a bound on the graph that is
$\lesssim W^{\delta_\sto/10}W^{2\epsilon}\eta_t^{-3/2}\eta_s^{\frac{1}{2}} + W^{\delta_\sto/10}W^{\epsilon}\eta_t^{-1}$
Multiplying this by $W^{\delta_\sto/5} W^{-2}\eta_s^{-1}$ results in the following bound:
$$\max_{i=1, \cdots, 4}| \mathcal{G}_{i1, s}|\lesssim W^{\frac{3\delta_\sto}{10}}W^{-2+2\epsilon}\eta_t^{-3/2}\eta_s^{-1/2}+W^{\frac{3\delta_\sto}{10}}W^{-2+\epsilon}\eta_t^{-1}\eta_s^{-1}$$
Following the same change of variables and integration, while also noting the fact that $\eta_0\lesssim 1$, gives us:
$$\Rightarrow \max_{i=1, \cdots, 4}\int_{0}^{t\wedge \tau_\sto}|\mathcal{G}_{i1, a}|ds\lesssim W^{\frac{3\delta_\sto}{10}}W^{-2+2\epsilon}\eta_t^{-3/2}\int_{\eta_t}^{\eta_0}\sigma^{-1/2}d\sigma + W^{\frac{3\delta_\sto}{10}+\epsilon}W^{-2}\eta_t^{-1}\int_{\eta_t}^{\eta_0}\sigma^{-1}d\sigma\lesssim$$
$$\lesssim W^{\frac{3\delta_\sto}{10}}W^{-2+2\epsilon}\eta_t^{-3/2}\eta_0^{-1}-W^{\frac{3\delta_\sto}{10}}W^{-2+2\epsilon}\eta_t^{-1}+W^{\frac{3\delta_\sto}{10}+\epsilon}W^{-2}\eta_t^{-1}\ln \eta_0-W^{\frac{3\delta_\sto}{10}+\epsilon}W^{-2}\eta_t^{-1}\ln \eta_t$$
$$\lesssim W^{\delta'}W^{-2}\eta_t^{-3/2}\lesssim W^{-\frac{3}{4}}\eta_t^{-1}W^{-1}\eta_t^{-\frac{1}{2}},$$
since $\delta' = \frac{3\delta_\sto}{10}+2\epsilon>0$ is arbitrarily small. 
\subsubsection*{Size estimates for $\mathcal{G}_{12, s}$}
Following the regular vertex expansion, we have to now represent each of the subgraphs that has the form $\mathcal{G}_{i2, s}=B_{u\gamma} S_{\gamma \delta} G_{x_i\delta} (G_{\gamma\gamma}-m)G_{\delta y_i} f(G)$. By applying the graphical representation we get: 
\begin{center}
\begin{tikzpicture}[
    thick, % Line thickness
    wavy/.style={decorate, decoration={snake, amplitude=2pt, segment length=5pt}},
    ]
\node (a) at (0,0) {a};
\node (x) at (2, 0) {x};
\node (delta) at (4, 0) {$\delta$};
\node (u) at (4, 1) {u};
\node (v) at (5, 1) {v};
\node (y) at (6, 0) {y};
\node (b) at (8, 0) {b};
\node (alpha) at (4, 2) {$\alpha$};
\node (beta) at (6, 2) {$\beta$};
\node (gamma) at (3, 1) {$\gamma$};

\draw[thick, double] (a) -- (x);
\draw[wavy] (u) -- (v);
\draw[blue, wavy] (gamma) to (u);
\draw[thick, double] (y) -- (b);
\draw[wavy, blue] (v) to (alpha);
\draw[wavy](gamma) to (delta);
\draw[wavy] (alpha) to (beta);
\draw[red, ->] (x) to [bend right=20] (y);
\draw[blue, ->] (delta) to (y);
\draw[blue, ->] (alpha) to [out=100, in=135, loop] (alpha);
\draw[blue, ->] (beta) to [out=100, in=70, loop] (beta);
\draw[blue, ->] (x) to (delta);
\draw[blue, ->] (gamma) to [out=100, in=135, loop] (gamma);
\end{tikzpicture}
\end{center}
\begin{equation*}\scalebox{0.85}{$
\mathcal{G}_{12, s} = 
\left[\{\text{Id} + (t-s)\Theta_t\}S^{1/2}\right]_{ax}\overline{G_{xy}} B_{u\gamma} S_{\gamma \delta} G_{x\delta} (G_{\gamma\gamma}-m)G_{\delta y} S_{uv} B_{v\alpha} (G_{\alpha\alpha}-m)S_{\alpha\beta}(G_{\beta\beta}-m)\left[\{\text{Id} + (t-s)\Theta_t\}S^{1/2}\right]_{yb}$}
\end{equation*}
Canceling all three loops $\textcolor{blue}{\circlearrowright}$ and the edge $x\textcolor{red}{\rightarrow}y$ gives us a combined term $O(W^{\delta_\sto/5}W^{-2}\eta_s^{-1})$ per estimate $(E2)$. Subsequently, we sum out all the wavy lines, which gives us the same bound as before:
$\mathcal{G}_{12, s}\lesssim W^{\delta_\sto/5}W^{-2}\eta_s^{-1}\times $\begin{tikzpicture}[
    thick, % Line thickness
    wavy/.style={decorate, decoration={snake, amplitude=2pt, segment length=5pt}},
    ]
\node (a) at (0,0) {a};
\node (x) at (1, 0) {x};
\node (delta) at (2, 0) {$\delta$};
\node (y) at (3, 0) {y};
\node (b) at (4, 0) {b};

\draw[thick, double] (a) -- (x);
\draw[thick, double] (y) -- (b);
\draw[blue, ->](x) to (delta);
\draw[blue, ->] (delta) to (y);

\end{tikzpicture}
\subsubsection*{Size estimates for $\mathcal{G}_{22, s}$}


\begin{center}
\begin{tikzpicture}[
    thick, % Line thickness
    wavy/.style={decorate, decoration={snake, amplitude=2pt, segment length=5pt}},
    ]
\node (a) at (0,0) {a};
\node (x) at (2, 0) {x};
\node (delta) at (4, 0) {$\delta$};
\node (u) at (4, 1) {u};
\node (v) at (5, 1) {v};
\node (y) at (6, 0) {y};
\node (b) at (8, 0) {b};
\node (alpha) at (4, 2) {$\alpha$};
\node (beta) at (6, 2) {$\beta$};
\node (gamma) at (3, 1) {$\gamma$};

\draw[thick, double] (a) -- (x);
\draw[wavy] (u) -- (v);
\draw[blue, wavy] (gamma) to (u);
\draw[thick, double] (y) -- (b);
\draw[wavy, blue] (v) to (alpha);
\draw[wavy](gamma) to (delta);
\draw[wavy] (alpha) to (beta);
\draw[blue, ->] (delta) to (y);
\draw[blue, ->] (x) to (delta);
\draw[blue, ->] (gamma) to [out=210, in=240, loop] (gamma);
\draw[blue, ->] (beta) to [bend right=30] (alpha);
\draw[red, ->] (x) to [bend left = 30] (alpha);
\draw[red, ->] (beta) to (y);
\end{tikzpicture}
\end{center}
\begin{equation*}\scalebox{0.95}{$
\mathcal{G}_{22, s} = 
\left[\{\text{Id} + (t-s)\Theta_t\}S^{1/2}\right]_{ax}
\overline{G_{x\alpha}} B_{u\gamma} S_{\gamma \delta} G_{x\delta} (G_{\gamma\gamma}-m)G_{\delta y}\overline{G_{\beta y}} S_{uv} B_{v\alpha} S_{\alpha\beta}G_{\beta\alpha}\left[\{\text{Id} + (t-s)\Theta_t\}S^{1/2}\right]_{yb}$}
\end{equation*}
Here we cancel the loop at $\gamma$ and the directed edges $x\textcolor{red}{\rightarrow}\alpha$, $\beta\textcolor{red}{\rightarrow}y$ and $\beta\textcolor{blue}{\rightarrow}\alpha$. With this, we are again left with a term $O(W^{\delta_\sto/5}W^{-2}\eta_s^{-1})$, s.t. after summing out the wavy lines using estimate $(E5)$, we get the same bound:
$\mathcal{G}_{22, s}\lesssim W^{\delta_\sto/5} W^{-2}\eta_s^{-1}\times $\begin{tikzpicture}[
    thick, % Line thickness
    wavy/.style={decorate, decoration={snake, amplitude=2pt, segment length=5pt}},
    ]
\node (a) at (0,0) {a};
\node (x) at (1, 0) {x};
\node (delta) at (2, 0) {$\delta$};
\node (y) at (3, 0) {y};
\node (b) at (4, 0) {b};

\draw[thick, double] (a) -- (x);
\draw[thick, double] (y) -- (b);
\draw[blue, ->](x) to (delta);
\draw[blue, ->] (delta) to (y);

\end{tikzpicture}

\subsubsection*{Size estimates for $\mathcal{G}_{32, s}$}

\begin{center}
\begin{tikzpicture}[
    thick, % Line thickness
    wavy/.style={decorate, decoration={snake, amplitude=2pt, segment length=5pt}},
    ]
\node (a) at (0,0) {a};
\node (x) at (2, 0) {x};
\node (delta) at (4, 0) {$\delta$};
\node (y) at (6, 0) {y};
\node (b) at (8, 0) {b};
\node (gamma) at (4, 1) {$\gamma$};
\node (u) at (5, 1) {u};
\node (v) at (6, 1) {v};
\node (alpha) at (4, 2) {$\alpha$};
\node (beta) at (2, 2) {$\beta$};


\draw[thick, double] (a) -- (x);
\draw[wavy] (u) -- (v);
\draw[blue, wavy] (gamma) to (u);
\draw[thick, double] (y) -- (b);
\draw[wavy, blue] (v) to (alpha);
\draw[wavy](gamma) to (delta);
\draw[wavy] (alpha) to (beta);
\draw[blue, ->] (delta) to (y);
\draw[blue, ->] (x) to (beta);
\draw[blue, ->] (beta) to [bend left=30] (alpha);
\draw[blue, ->] (gamma) to [out=40, in=70, loop] (gamma);
\draw[blue, ->] (alpha) to [bend right=30] (delta);
\draw[red, ->] (x) to [bend right=20] (y);
\end{tikzpicture}
\end{center}
\begin{equation*}\scalebox{0.95}{$
\mathcal{G}_{32, s} = \left[\{\text{Id} + (t-s)\Theta_t\}S^{1/2}\right]_{ax}
\overline{G_{xy}}G_{x\beta} B_{u\gamma} S_{\gamma \delta} G_{\alpha \delta} (G_{\gamma\gamma}-m)G_{\delta y}S_{uv}B_{v\alpha}S_{\alpha\beta}G_{\beta \alpha} 
\left[\{\text{Id} + (t-s)\Theta_t\}S^{1/2}\right]_{yb}$}
\end{equation*}
For this subgraph, we start by eliminating the loop at $\gamma$ and the edges $\beta\textcolor{blue}{\rightarrow}\alpha$, $\alpha\textcolor{blue}{\rightarrow}\delta$, and $x\textcolor{red}{\rightarrow}y$. This means that get a combined term $O(W^{\delta_\sto/5}W^{-2}\eta_s^{-1})$. With this, we have a similar subgraph remaining with a few extra wavy edges that will sum out to $O(1)$ per estimate $(E5)$, hence the total estimate will be equivalent to:
$\mathcal{G}_{32, s}\lesssim W^{\delta_\sto/5}W^{-2}\eta_s^{-1}\times $\begin{tikzpicture}[
    thick, % Line thickness
    wavy/.style={decorate, decoration={snake, amplitude=2pt, segment length=5pt}},
    ]
\node (a) at (0,0) {a};
\node (x) at (1, 0) {x};
\node (delta) at (2, 0) {$\delta$};
\node (y) at (3, 0) {y};
\node (b) at (4, 0) {b};

\draw[thick, double] (a) -- (x);
\draw[thick, double] (y) -- (b);
\draw[blue, ->](x) to (delta);
\draw[blue, ->] (delta) to (y);

\end{tikzpicture}

\subsubsection*{Size estimate for $\mathcal{G}_{42, s}$}

\begin{center}
\begin{tikzpicture}[
    thick, % Line thickness
    wavy/.style={decorate, decoration={snake, amplitude=2pt, segment length=5pt}},
    ]
\node (a) at (0,0) {a};
\node (x) at (2, 0) {x};
\node (gamma) at (4, 0) {$\gamma$};
\node (v) at (4, 1) {v};
\node (u) at (5, 1) {u};
\node (y) at (6, 0) {y};
\node (b) at (8, 0) {b};
\node (alpha) at (4, 2) {$\alpha$};
\node (beta) at (2, 2) {$\beta$};
\node (delta) at (3, 1) {$\delta$};

\draw[thick, double] (a) -- (x);
\draw[wavy] (u) -- (v);
\draw[blue, wavy] (gamma) to (u);
\draw[thick, double] (y) -- (b);
\draw[wavy, blue] (v) to (alpha);
\draw[wavy](gamma) to (delta);
\draw[blue, ->] (gamma) to[out=180, in=200, loop] (gamma);
\draw[wavy] (alpha) to (beta);
\draw[blue, <-] (alpha) to[bend right=20] (beta);
\draw[blue, ->] (delta) to (beta);
\draw[blue, ->] (x) to (delta);
\draw[red, ->] (x) to [bend right=20] (y);
\draw[blue, ->] (alpha) to [bend left=30](y);
\end{tikzpicture}
\end{center}
\begin{equation*}\scalebox{0.95}{$
\mathcal{G}_{42, s} = \left[\{\text{Id} + (t-s)\Theta_t\}S^{1/2}\right]_{ax}
\overline{G_{xy}}B_{u\gamma} S_{\gamma \delta}  G_{x\delta} (G_{\gamma \gamma}-m)G_{\delta \beta}S_{uv} B_{v\alpha} S_{\alpha \beta} G_{\beta \alpha} G_{\alpha y}
\left[\{\text{Id} + (t-s)\Theta_t\}S^{1/2}\right]_{yb}$}
\end{equation*}
We eliminate the loop at $\gamma$ and the edges $\delta\textcolor{blue}{\rightarrow}\beta$, $\beta\textcolor{blue}{\rightarrow}\alpha$ and $x\textcolor{red}{\rightarrow}y$, all contributing $O(W^{\delta_\sto/20}W^{-1/2}\eta_s^{-1/4})$. By summing over $\beta\rightsquigarrow \alpha$ for an $O(1)$ term, we again have a bound that is equivalent to the previous ones, since all the wavy lines sum up to contribute $O(1)$. As such, we have a bound, equivalent to:
$\mathcal{G}_{42, s}\lesssim W^{\delta_\sto/5}W^{-2}\eta_s^{-1}\times $\begin{tikzpicture}[
    thick, % Line thickness
    wavy/.style={decorate, decoration={snake, amplitude=2pt, segment length=5pt}},
    ]
\node (a) at (0,0) {a};
\node (x) at (1, 0) {x};
\node (delta) at (2, 0) {$\delta$};
\node (y) at (3, 0) {y};
\node (b) at (4, 0) {b};

\draw[thick, double] (a) -- (x);
\draw[thick, double] (y) -- (b);
\draw[blue, ->](x) to (delta);
\draw[blue, ->] (delta) to (y);

\end{tikzpicture}. But these are exactly the same bounds as for $\mathcal{G}_{i1, s}$, which means that: $$\max_{i=1, \cdots, 4}| \mathcal{G}_{i2, s}|\lesssim W^{\delta_\sto/4+2\epsilon}W^{-2}\eta_t^{-3/2}\lesssim  W^{-\frac{3}{4}}\eta_t^{-1}W^{-1}\eta_t^{-\frac{1}{2}}$$


\subsubsection*{Size estimate for $\mathcal{G}_{13, s}^{(1)}$}
The last non-regularization term from the regular vertex expansion contains a partial derivative: $\mathcal{G}_{i3, s} =B_{u\gamma} S_{\gamma \delta} G_{x_i\gamma} G_{\delta y_i} \partial_{H_{\delta \gamma}}(f(G))$. Recall the resolvent perturbation by which we have that $\partial_{H_{\delta \gamma}}\overline{G_{x_iy_i}} = -\overline{G_{x_i\gamma}G_{\delta y_i}}$ and $\partial_{H_{\delta \gamma}}G_{x_iy_i} = -G_{x_i\delta}G_{\gamma y_i}$. Applying the latter identities results in summation of 3 separate subgraphs for each $G_{i3, s}$. We can again represent each of them individually for $i=1, 2, 3, 4$:
\begin{center}
\begin{tikzpicture}[
    thick, % Line thickness
    wavy/.style={decorate, decoration={snake, amplitude=2pt, segment length=5pt}},
    ]
\node (a) at (0,0) {a};
\node (x) at (2, 0) {x};
\node (delta) at (4, 0) {$\delta$};
\node (y) at (6, 0) {y};
\node (b) at (8, 0) {b};
\node (gamma) at (4, 1) {$\gamma$};
\node (u) at (5, 1) {u};
\node (v) at (6, 1) {v};
\node (alpha) at (6, 2) {$\alpha$};
\node (beta) at (4, 2) {$\beta$};


\draw[thick, double] (a) -- (x);
\draw[wavy] (u) -- (v);
\draw[blue, wavy] (gamma) to (u);
\draw[thick, double] (y) -- (b);
\draw[wavy, blue] (v) to (alpha);
\draw[wavy](gamma) to (delta);
\draw[wavy] (alpha) to (beta);
\draw[blue, ->] (delta) to (y);
\draw[blue, ->] (beta) to [out=130, in=170, loop] (beta);
\draw[blue, ->] (alpha) to [out=20, in=50, loop] (alpha);
\draw[blue, ->] (x) to (delta);
\draw[blue, ->] (gamma) to (y);
\draw[blue, ->] (x) to (gamma);
\end{tikzpicture}
\end{center}
\begin{equation*}\scalebox{0.9}{$
\mathcal{G}_{13, s}^{(1)} =  \left[\{\text{Id} + (t-s)\Theta_t\}S^{1/2}\right]_{ax}
G_{x\delta} G_{\gamma y}S_{uv}B_{v\alpha}(G_{\alpha\alpha}-m)S_{\alpha\beta}(G_{\beta\beta}-m)B_{u\gamma}S_{\gamma\delta}G_{x\gamma}G_{\delta y}
\left[\{\text{Id} + (t-s)\Theta_t\}S^{1/2}\right]_{yb}$}
\end{equation*}
Here, we bound the two loops at $\beta, \alpha$ and then the edges $x\textcolor{blue}{\rightarrow}\gamma$ and $\gamma\textcolor{blue}{\rightarrow}y$, getting a combined term $O(W^{\delta_\sto/5}W^{-2}\eta_s^{-1})$. Hence, after summing over the wavy lines, we get the same bound as before:
$\mathcal{G}_{13, s}^{(1)}\lesssim W^{\delta_\sto/5}W^{-2}\eta_s^{-1}\times $\begin{tikzpicture}[
    thick, % Line thickness
    wavy/.style={decorate, decoration={snake, amplitude=2pt, segment length=5pt}},
    ]
\node (a) at (0,0) {a};
\node (x) at (1, 0) {x};
\node (delta) at (2, 0) {$\delta$};
\node (y) at (3, 0) {y};
\node (b) at (4, 0) {b};

\draw[thick, double] (a) -- (x);
\draw[thick, double] (y) -- (b);
\draw[blue, ->](x) to (delta);
\draw[blue, ->] (delta) to (y);

\end{tikzpicture}

\subsubsection*{Size estimate for $\mathcal{G}_{13, s}^{(2)}$}

\begin{center}
\begin{tikzpicture}[
    thick, % Line thickness
    wavy/.style={decorate, decoration={snake, amplitude=2pt, segment length=5pt}},
    ]
\node (a) at (0,0) {a};
\node (x) at (2, 0) {x};
\node (gamma) at (4, 0) {$\gamma$};
\node (u) at (4, 1) {u};
\node (delta) at (5, 1) {$\delta$};
\node (y) at (6, 0) {y};
\node (b) at (8, 0) {b};
\node (alpha) at (4, 2) {$\alpha$};
\node (beta) at (2, 2) {$\beta$};
\node (v) at (3, 1) {v};

\draw[thick, double] (a) -- (x);
\draw[wavy] (u) -- (v);
\draw[blue, wavy] (gamma) to (u);
\draw[thick, double] (y) -- (b);
\draw[wavy, blue] (v) to (alpha);
\draw[wavy](gamma) to (delta);
\draw[blue, ->] (beta) to[out=180, in=220, loop] (beta);
\draw[wavy] (alpha) to (beta);
\draw[red, ->] (x) to [bend right=20] (y);
\draw[blue, <-] (alpha) to [bend left=30](gamma);
\draw[blue, ->] (x) to (gamma);
\draw[blue, ->] (delta) to (y);
\draw[blue, ->] (alpha) to[bend left=20] (delta);
\end{tikzpicture}
\end{center}
\begin{equation*}\scalebox{0.95}{$
\mathcal{G}_{13, s}^{(2)} =  \left[\{\text{Id} + (t-s)\Theta_t\}S^{1/2}\right]_{ax}
\overline{G_{xy}}S_{uv}B_{v\alpha}G_{\alpha\delta}G_{\gamma\alpha}S_{\alpha\beta}(G_{\beta\beta}-m)B_{u\gamma}S_{\gamma\delta}G_{x\gamma}G_{\delta y}
\left[\{\text{Id} + (t-s)\Theta_t\}S^{1/2}\right]_{yb}$}
\end{equation*}
We eliminate the loop at $\beta$ and then the edges $x\textcolor{red}{\rightarrow}\gamma$, $\gamma\textcolor{blue}{\rightarrow}\alpha$, and $\alpha\textcolor{blue}{\rightarrow}\delta$, getting back $O(W^{\delta_\sto/5}W^{-2}\eta_s^{-1})$ by estimate $(E2)$. Observe that with this we have almost the same subgraph as before with an additional wavy line $\gamma\rightsquigarrow \delta$, i.e our estimate looks like: 
$\mathcal{G}_{13, s}^{(2)}\lesssim W^{\delta_\sto/5}W^{-2}\eta_s^{-1}\times $\begin{tikzpicture}[
    thick, % Line thickness
    wavy/.style={decorate, decoration={snake, amplitude=2pt, segment length=5pt}},
    ]
\node (a) at (0,0) {a};
\node (x) at (1, 0) {x};
\node (gamma) at (2, 0) {$\gamma$};
\node (delta) at (3, 0) {$\delta$};
\node (y) at (4, 0) {y};
\node (b) at (5, 0) {b};

\draw[thick, double] (a) -- (x);
\draw[thick, double] (y) -- (b);
\draw[blue, ->](x) to (gamma);
\draw[wavy, orange] (gamma) to (delta);
\draw[blue, ->] (delta) to (y);

\end{tikzpicture}

\subsubsection*{Size estimate for $\mathcal{G}_{13, s}^{(3)}$}

\begin{center}
\begin{tikzpicture}[
    thick, % Line thickness
    wavy/.style={decorate, decoration={snake, amplitude=2pt, segment length=5pt}},
    ]
\node (a) at (0,0) {a};
\node (x) at (2, 0) {x};
\node (gamma) at (4, 0) {$\gamma$};
\node (u) at (4, 1) {u};
\node (delta) at (6, 1) {$\delta$};
\node (y) at (6, 0) {y};
\node (b) at (8, 0) {b};
\node (alpha) at (4, 2) {$\alpha$};
\node (beta) at (6, 2) {$\beta$};
\node (v) at (3, 1) {v};

\draw[thick, double] (a) -- (x);
\draw[wavy] (u) -- (v);
\draw[blue, wavy] (gamma) to (u);
\draw[thick, double] (y) -- (b);
\draw[wavy, blue] (v) to (alpha);
\draw[wavy](gamma) to (delta);
\draw[blue, ->] (alpha) to[out=200, in=160, loop] (alpha);
\draw[wavy] (alpha) to (beta);
\draw[red, ->] (x) to [bend right=20] (y);
\draw[blue, <-] (beta) to (gamma);
\draw[blue, ->] (x) to (gamma);
\draw[blue, ->] (delta) to (y);
\draw[blue, ->] (beta) to (delta);
\end{tikzpicture}
\end{center}
\begin{equation*}\scalebox{0.95}{$
\mathcal{G}_{13, s}^{(3)} =  \left[\{\text{Id} + (t-s)\Theta_t\}S^{1/2}\right]_{ax}
\overline{G_{xy}}S_{uv}B_{v\alpha}(G_{\alpha\alpha}-m)S_{\alpha\beta}G_{\beta\delta}G_{\gamma\beta}B_{u\gamma}S_{\gamma\delta}G_{x\gamma}G_{\delta y}
\left[\{\text{Id} + (t-s)\Theta_t\}S^{1/2}\right]_{yb}$}
\end{equation*}

\noindent For the last subgraph, we start with the loop $\alpha$ and then apply the second estimate for the edges $x\textcolor{red}{\rightarrow}y$, $\gamma\textcolor{blue}{\rightarrow}\beta$, and $\beta\textcolor{blue}{\rightarrow}\delta$. We get a combined term $O(W^{\delta_\sto/5}W^{-2}\eta_s^{-1})$, s.t. after summing out $\beta, \alpha, v, u, $, we get a final estimate that looks like the one before:
$\mathcal{G}_{13, s}^{(3)}\lesssim W^{\delta_\sto/5}W^{-2}\eta_s^{-1}\times $\begin{tikzpicture}[
    thick, % Line thickness
    wavy/.style={decorate, decoration={snake, amplitude=2pt, segment length=5pt}},
    ]
\node (a) at (0,0) {a};
\node (x) at (1, 0) {x};
\node (gamma) at (2, 0) {$\gamma$};
\node (delta) at (3, 0) {$\delta$};
\node (y) at (4, 0) {y};
\node (b) at (5, 0) {b};

\draw[thick, double] (a) -- (x);
\draw[thick, double] (y) -- (b);
\draw[blue, ->](x) to (gamma);
\draw[wavy, orange] (gamma) to (delta);
\draw[blue, ->] (delta) to (y);

\end{tikzpicture}. And since the partial derivative expansion is equivalent to the summation of the latter three subgraphs, we get a combined bound on $\mathcal{G}_{i3, s}$ that is $O(W^{-\frac{3}{4}}\eta_t^{-1}W^{-1}\eta_t^{-\frac{1}{2}})$ by using our earlier calculations. 
\subsubsection*{Size estimate for $\mathcal{G}_{23, s}^{(1)}$}
\begin{center}
\begin{tikzpicture}[
    thick, % Line thickness
    wavy/.style={decorate, decoration={snake, amplitude=2pt, segment length=5pt}},
    ]
\node (a) at (0,0) {a};
\node (x) at (2, 0) {x};
\node (gamma) at (4, 0) {$\gamma$};
\node (u) at (4, 1) {u};
\node (delta) at (5, 1) {$\delta$};
\node (y) at (6, 0) {y};
\node (b) at (8, 0) {b};
\node (alpha) at (4, 2) {$\alpha$};
\node (beta) at (6, 2) {$\beta$};
\node (v) at (3, 1) {v};

\draw[thick, double] (a) -- (x);
\draw[wavy] (u) -- (v);
\draw[blue, wavy] (gamma) to (u);
\draw[thick, double] (y) -- (b);
\draw[wavy, blue] (v) to (alpha);
\draw[wavy](gamma) to (delta);
\draw[blue, <-] (alpha) to[bend left=30] (beta);
\draw[wavy] (alpha) to (beta);
\draw[red, ->] (x) to [bend right=20] (gamma);
\draw[blue, ->] (x) to [bend left=20](gamma);
\draw[blue, ->] (delta) to (y);
\draw[red, ->] (beta) to (y);
\draw[red, <-] (alpha) to (delta);
\end{tikzpicture}
\end{center}
\begin{equation*}
\mathcal{G}_{23, s}^{(1)} =  \left[\{\text{Id} + (t-s)\Theta_t\}S^{1/2}\right]_{ax}
\overline{G_{x\gamma}G_{\delta \alpha}G_{\beta y}}S_{uv}B_{v\alpha}S_{\alpha\beta}G_{\beta\alpha}B_{u\gamma}S_{\gamma\delta}G_{x\gamma}G_{\delta y}
\left[\{\text{Id} + (t-s)\Theta_t\}S^{1/2}\right]_{yb}
\end{equation*}
For this subgraph, we can bound the three red edges along with $\beta\textcolor{blue}{\rightarrow}\alpha$, for which we get combined term $O( W^{\delta_\sto/5}W^{-2}\eta_s^{-1})$. After summing over the wavy edges in order $\beta, \alpha, v, u$, we get the same bound as before:
$\mathcal{G}_{23, s}^{(1)}\lesssim W^{\delta_\sto/5}W^{-2}\eta_s^{-1}\times $\begin{tikzpicture}[
    thick, % Line thickness
    wavy/.style={decorate, decoration={snake, amplitude=2pt, segment length=5pt}},
    ]
\node (a) at (0,0) {a};
\node (x) at (1, 0) {x};
\node (gamma) at (2, 0) {$\gamma$};
\node (delta) at (3, 0) {$\delta$};
\node (y) at (4, 0) {y};
\node (b) at (5, 0) {b};

\draw[thick, double] (a) -- (x);
\draw[thick, double] (y) -- (b);
\draw[blue, ->](x) to (gamma);
\draw[wavy, orange] (gamma) to (delta);
\draw[blue, ->] (delta) to (y);

\end{tikzpicture}

\subsubsection*{Size estimate for $\mathcal{G}_{23, s}^{(2)}$}

\begin{center}
\begin{tikzpicture}[
    thick, % Line thickness
    wavy/.style={decorate, decoration={snake, amplitude=2pt, segment length=5pt}},
    ]
\node (a) at (0,0) {a};
\node (x) at (2, 0) {x};
\node (gamma) at (4, 0) {$\gamma$};
\node (u) at (4, 1) {u};
\node (delta) at (6, 1) {$\delta$};
\node (y) at (6, 0) {y};
\node (b) at (8, 0) {b};
\node (alpha) at (4, 2) {$\alpha$};
\node (beta) at (6, 2) {$\beta$};
\node (v) at (3, 1) {v};

\draw[thick, double] (a) -- (x);
\draw[wavy] (u) -- (v);
\draw[blue, wavy] (gamma) to (u);
\draw[thick, double] (y) -- (b);
\draw[wavy, blue] (v) to (alpha);
\draw[wavy](gamma) to (delta);
\draw[blue, <-] (alpha) to[bend left=30] (beta);
\draw[wavy] (alpha) to (beta);
\draw[blue, ->] (x) to(gamma);
\draw[red, ->] (x) to [bend left=30](alpha);
\draw[red, ->] (delta) to[bend left=20] (y);
\draw[blue, ->] (delta) to[bend right=20] (y);
\draw[red, ->] (beta) to[bend right=10] (gamma);
\end{tikzpicture}
\end{center}
\begin{equation*}
\mathcal{G}_{23, s}^{(2)} =  \left[\{\text{Id} + (t-s)\Theta_t\}S^{1/2}\right]_{ax}
\overline{G_{x\alpha}G_{\beta \gamma}G_{\delta y}}S_{uv}B_{v\alpha}S_{\alpha\beta}G_{\beta\alpha}B_{u\gamma}S_{\gamma\delta}G_{x\gamma}G_{\delta y}
\left[\{\text{Id} + (t-s)\Theta_t\}S^{1/2}\right]_{yb}
\end{equation*}
Here, we can get bound all the red lines along with $\beta\textcolor{blue}{\rightarrow}\alpha$, s.t. after the summation over the wavy lines in order $\beta, \alpha, v, u$, we get the familiar:
$\mathcal{G}_{23, s}^{(2)}\lesssim W^{\delta_\sto/5} W^{-2}\eta_s^{-1}\times $\begin{tikzpicture}[
    thick, % Line thickness
    wavy/.style={decorate, decoration={snake, amplitude=2pt, segment length=5pt}},
    ]
\node (a) at (0,0) {a};
\node (x) at (1, 0) {x};
\node (gamma) at (2, 0) {$\gamma$};
\node (delta) at (3, 0) {$\delta$};
\node (y) at (4, 0) {y};
\node (b) at (5, 0) {b};

\draw[thick, double] (a) -- (x);
\draw[thick, double] (y) -- (b);
\draw[blue, ->](x) to (gamma);
\draw[wavy, orange] (gamma) to (delta);
\draw[blue, ->] (delta) to (y);

\end{tikzpicture}
\subsubsection*{Size estimate for $\mathcal{G}_{23, s}^{(3)}$}

\begin{center}
\begin{tikzpicture}[
    thick, % Line thickness
    wavy/.style={decorate, decoration={snake, amplitude=2pt, segment length=5pt}},
    ]
\node (a) at (0,0) {a};
\node (x) at (2, 0) {x};
\node (gamma) at (4, 0) {$\gamma$};
\node (u) at (4, 1) {u};
\node (delta) at (5, 1) {$\delta$};
\node (y) at (6, 0) {y};
\node (b) at (8, 0) {b};
\node (alpha) at (4, 2) {$\alpha$};
\node (beta) at (6, 2) {$\beta$};
\node (v) at (3, 1) {v};

\draw[thick, double] (a) -- (x);
\draw[wavy] (u) -- (v);
\draw[blue, wavy] (gamma) to (u);
\draw[thick, double] (y) -- (b);
\draw[wavy, blue] (v) to (alpha);
\draw[wavy](gamma) to (delta);
\draw[blue, <-] (alpha) to[bend left=20] (gamma);
\draw[wavy] (alpha) to (beta);
\draw[blue, ->] (x) to(gamma);
\draw[red, ->] (x) to [bend left=30](alpha);
\draw[blue, <-] (delta) to (beta);
\draw[blue, ->] (delta) to (y);
\draw[red, ->] (beta) to(y);
\end{tikzpicture}
\end{center}
\begin{equation*}
\mathcal{G}_{23, s}^{(3)} =  \left[\{\text{Id} + (t-s)\Theta_t\}S^{1/2}\right]_{ax}
\overline{G_{x\alpha}G_{\beta y}}S_{uv}B_{v\alpha}S_{\alpha\beta}G_{\beta\delta}G_{\gamma \alpha}B_{u\gamma}S_{\gamma\delta}G_{x\gamma}G_{\delta y}
\left[\{\text{Id} + (t-s)\Theta_t\}S^{1/2}\right]_{yb}
\end{equation*}

\noindent Lastly, by bounding all the red edges along with $\gamma\textcolor{blue}{\rightarrow}\alpha$ and $\beta\textcolor{blue}{\rightarrow}\delta$ and summing over the wavy edges in order $\beta, \alpha, v, u$, we get the equivalent bound: $\mathcal{G}_{23, s}^{(3)}\lesssim W^{\delta_\sto/5}W^{-2}\eta_s^{-1}\times $\begin{tikzpicture}[
    thick, % Line thickness
    wavy/.style={decorate, decoration={snake, amplitude=2pt, segment length=5pt}},
    ]
\node (a) at (0,0) {a};
\node (x) at (1, 0) {x};
\node (gamma) at (2, 0) {$\gamma$};
\node (delta) at (3, 0) {$\delta$};
\node (y) at (4, 0) {y};
\node (b) at (5, 0) {b};

\draw[thick, double] (a) -- (x);
\draw[thick, double] (y) -- (b);
\draw[blue, ->](x) to (gamma);
\draw[wavy, orange] (gamma) to (delta);
\draw[blue, ->] (delta) to (y);

\end{tikzpicture}, which makes this entire subgraph subject to the same $O(W^{-\frac{3}{4}}\eta_t^{-1}W^{-1}\eta_t^{-\frac{1}{2}})$ estimate. 

\subsubsection*{Size estimate for $\mathcal{G}_{33, s}^{(1)}$}

\begin{center}
\begin{tikzpicture}[
    thick, % Line thickness
    wavy/.style={decorate, decoration={snake, amplitude=2pt, segment length=5pt}},
    ]
\node (a) at (0,0) {a};
\node (x) at (2, 0) {x};
\node (gamma) at (4, 0) {$\gamma$};
\node (u) at (4, 1) {u};
\node (delta) at (6, 1) {$\delta$};
\node (y) at (6, 0) {y};
\node (b) at (8, 0) {b};
\node (alpha) at (4, 2) {$\alpha$};
\node (beta) at (2, 2) {$\beta$};
\node (v) at (3, 1) {v};

\draw[thick, double] (a) -- (x);
\draw[wavy] (u) -- (v);
\draw[blue, wavy] (gamma) to (u);
\draw[thick, double] (y) -- (b);
\draw[wavy, blue] (v) to (alpha);
\draw[wavy](gamma) to (delta);
\draw[blue, ->] (alpha) to[bend left=30] (gamma);
\draw[wavy] (alpha) to (beta);
\draw[red, ->] (x) to(gamma);
\draw[blue, ->] (x) to (beta);
\draw[blue, ->] (delta) to[bend left=20] (y);
\draw[red, ->] (delta) to[bend right=20](y);
\draw[blue, <-] (alpha) to[bend right=30] (beta);
\end{tikzpicture}
\end{center}
\begin{equation*}
\mathcal{G}_{33, s}^{(1)} =  \left[\{\text{Id} + (t-s)\Theta_t\}S^{1/2}\right]_{ax}
\overline{G_{x\gamma}G_{\delta y}}G_{x\beta}S_{uv} B_{v\alpha}S_{\alpha \beta} G_{\beta \alpha} B_{u\gamma} S_{\gamma\delta} G_{\alpha \gamma}G_{\delta y}
\left[\{\text{Id} + (t-s)\Theta_t\}S^{1/2}\right]_{yb}
\end{equation*}

\noindent By bounding all the blue edges we get a combined $O(W^{\delta_\sto/5}W^{-2}\eta_s^{-1})$ term, s.t. after summing over the wavy edges in order $\beta, \alpha, v, u$, we get a  bound:
$\mathcal{G}_{33, s}^{(1)}\lesssim W^{\delta_\sto/5}W^{-2}\eta_s^{-1}\times $\begin{tikzpicture}[
    thick, % Line thickness
    wavy/.style={decorate, decoration={snake, amplitude=2pt, segment length=5pt}},
    ]
\node (a) at (0,0) {a};
\node (x) at (1, 0) {x};
\node (gamma) at (2, 0) {$\gamma$};
\node (delta) at (3, 0) {$\delta$};
\node (y) at (4, 0) {y};
\node (b) at (5, 0) {b};

\draw[thick, double] (a) -- (x);
\draw[thick, double] (y) -- (b);
\draw[red, ->](x) to (gamma);
\draw[wavy, orange] (gamma) to (delta);
\draw[red, ->] (delta) to (y);

\end{tikzpicture}

\subsubsection*{Size estimate for $\mathcal{G}_{33, s}^{(2)}$}

\begin{center}
\begin{tikzpicture}[
    thick, % Line thickness
    wavy/.style={decorate, decoration={snake, amplitude=2pt, segment length=5pt}},
    ]
\node (a) at (0,0) {a};
\node (x) at (2, 0) {x};
\node (delta) at (4, 0) {$\delta$};
\node (y) at (6, 0) {y};
\node (b) at (8, 0) {b};
\node (gamma) at (4, 1) {$\gamma$};
\node (u) at (5, 1) {u};
\node (v) at (6, 1) {v};
\node (alpha) at (6, 2) {$\alpha$};
\node (beta) at (4, 2) {$\beta$};


\draw[thick, double] (a) -- (x);
\draw[wavy] (u) -- (v);
\draw[blue, wavy] (gamma) to (u);
\draw[thick, double] (y) -- (b);
\draw[wavy, blue] (v) to (alpha);
\draw[wavy](gamma) to (delta);
\draw[wavy] (alpha) to (beta);
\draw[blue, ->] (delta) to(y);
\draw[blue, <-] (alpha) to [bend right=30] (beta);
\draw[blue, ->] (x) to (delta);
\draw[blue, ->] (gamma) to (beta);
\draw[red, ->] (x) to [bend right=25](y);
\draw[blue, ->] (alpha) to (gamma);
\end{tikzpicture}
\end{center}
\begin{equation*}
\mathcal{G}_{33, s}^{(2)} =  \left[\{\text{Id} + (t-s)\Theta_t\}S^{1/2}\right]_{ax}
\overline{G_{x y}}G_{x\delta}G_{\gamma \beta}S_{uv} B_{v\alpha}S_{\alpha \beta} G_{\beta \alpha} B_{u\gamma} S_{\gamma\delta} G_{\alpha \gamma}G_{\delta y}
\left[\{\text{Id} + (t-s)\Theta_t\}S^{1/2}\right]_{yb}
\end{equation*}
For this graph, let us bound the edges $x\textcolor{red}{\rightarrow}y$, $\gamma\textcolor{blue}{\rightarrow}\beta$, $\beta\textcolor{blue}{\rightarrow}\alpha$ and $\alpha\textcolor{blue}{\rightarrow}\gamma$, s.t. after summing over the wavy lines in order of the indices $\beta, \alpha, v, u, \gamma$, we get a bound:$\mathcal{G}_{33, s}^{(2)}\lesssim W^{\delta_\sto/5}W^{-2}\eta_s^{-1}\times $\begin{tikzpicture}[
    thick, % Line thickness
    wavy/.style={decorate, decoration={snake, amplitude=2pt, segment length=5pt}},
    ]
\node (a) at (0,0) {a};
\node (x) at (1, 0) {x};
\node (delta) at (2, 0) {$\delta$};
\node (y) at (3, 0) {y};
\node (b) at (4, 0) {b};

\draw[thick, double] (a) -- (x);
\draw[thick, double] (y) -- (b);
\draw[blue, ->](x) to (delta);
\draw[blue, ->] (delta) to (y);

\end{tikzpicture}

\subsubsection*{Size estimate for $\mathcal{G}_{33, s}^{(3)}$}


\begin{center}
\begin{tikzpicture}[
    thick, % Line thickness
    wavy/.style={decorate, decoration={snake, amplitude=2pt, segment length=5pt}},
    ]
\node (a) at (0,0) {a};
\node (x) at (2, 0) {x};
\node (delta) at (4, 0) {$\delta$};
\node (y) at (6, 0) {y};
\node (b) at (8, 0) {b};
\node (gamma) at (4, 1) {$\gamma$};
\node (u) at (5, 1) {u};
\node (v) at (6, 1) {v};
\node (alpha) at (5, 2) {$\alpha$};
\node (beta) at (3, 2) {$\beta$};


\draw[thick, double] (a) -- (x);
\draw[wavy] (u) -- (v);
\draw[blue, wavy] (gamma) to (u);
\draw[thick, double] (y) -- (b);
\draw[wavy, blue] (v) to (alpha);
\draw[wavy](gamma) to (delta);
\draw[wavy] (alpha) to (beta);
\draw[blue, ->] (delta) to(y);
\draw[blue, ->] (x) to (beta);
\draw[blue, <-] (delta) to (beta);
\draw[red, ->] (x) to [bend right=25](y);
\draw[blue, ->] (alpha) to [bend left=15](gamma);
\draw[red, ->] (gamma) to [bend left=15] (alpha);
\end{tikzpicture}
\end{center}
\begin{equation*}
\mathcal{G}_{33, s}^{(3)} =  \left[\{\text{Id} + (t-s)\Theta_t\}S^{1/2}\right]_{ax}
\overline{G_{x y}}G_{x\beta}S_{uv} B_{v\alpha}S_{\alpha \beta} G_{\beta \delta}G_{\gamma \alpha} B_{u\gamma} S_{\gamma\delta} G_{\alpha \gamma}G_{\delta y}
\left[\{\text{Id} + (t-s)\Theta_t\}S^{1/2}\right]_{yb}
\end{equation*}
Lastly, by bounding all of the red edges along with $\beta\textcolor{blue}{\rightarrow}\delta$, $\alpha\textcolor{blue}{\rightarrow}\gamma$, we get an estimate $O(W^{\delta_\sto/5}W^{-2}\eta_s^{-1})$. This leaves us with a bound that looks like 
$\mathcal{G}_{33, s}^{(3)}\lesssim W^{\delta_\sto/5}W^{-2}\eta_s^{-1}\times $
\begin{tikzpicture}[
    thick, % Line thickness
    wavy/.style={decorate, decoration={snake, amplitude=2pt, segment length=5pt}},
    ]
\node (a) at (0,0) {a};
\node (x) at (1, 0) {x};
\node (beta) at (2, 0) {$\beta$};
\node (delta) at (3, 0){$\delta$};
\node (y) at (4, 0) {y};
\node (b) at (5, 0) {b};

\draw[thick, double] (a) -- (x);
\draw[thick, double] (y) -- (b);
\draw[blue, ->](x) to (beta);
\draw [wavy, orange] (beta) to (delta);
\draw[blue, ->] (delta) to (y);

\end{tikzpicture}, which clearly gives for this set of subgraphs the same bound $O(W^{-\frac{3}{4}}\eta_t^{-1}W^{-1}\eta_t^{-\frac{1}{2}})$ as the ones before. 
 


\subsubsection*{Size estimate for $\mathcal{G}_{43, s}^{(1)}$}
\begin{center}
\begin{tikzpicture}[
    thick, % Line thickness
    wavy/.style={decorate, decoration={snake, amplitude=2pt, segment length=5pt}},
    ]
\node (a) at (0,0) {a};
\node (x) at (2, 0) {x};
\node (delta) at (4, 0) {$\delta$};
\node (y) at (6, 0) {y};
\node (b) at (8, 0) {b};
\node (gamma) at (4, 1) {$\gamma$};
\node (u) at (2, 1) {u};
\node (v) at (2, 2) {v};
\node (alpha) at (4, 2) {$\alpha$};
\node (beta) at (6, 2) {$\beta$};


\draw[thick, double] (a) -- (x);
\draw[wavy] (u) -- (v);
\draw[blue, wavy] (gamma) to (u);
\draw[thick, double] (y) -- (b);
\draw[wavy, blue] (v) to (alpha);
\draw[wavy](gamma) to (delta);
\draw[wavy] (alpha) to (beta);
\draw[red, ->] (delta) to(y);
\draw[blue, <-] (alpha) to [bend left=30] (beta);
\draw[blue, ->] (x) to (gamma);
\draw[blue, ->] (delta) to (beta);
\draw[red, ->] (x) to [bend right=25](gamma);
\draw[blue, ->] (alpha) to (y);
\end{tikzpicture}
\end{center}
\begin{equation*}
\mathcal{G}_{43, s}^{(1)} =  \left[\{\text{Id} + (t-s)\Theta_t\}S^{1/2}\right]_{ax}
\overline{G_{x \gamma}G_{\delta y}}S_{uv} B_{v\alpha}S_{\alpha \beta} G_{\beta \alpha}G_{\alpha y} B_{u\gamma} S_{\gamma\delta} G_{x\gamma}G_{\delta \beta}
\left[\{\text{Id} + (t-s)\Theta_t\}S^{1/2}\right]_{yb}
\end{equation*}
Bounding all the blue edges for $O(W^{\delta_\sto/20}W^{-1/2}\eta_s^{-1/4})$ and then summing over the wavy lines in order $\beta,\alpha, v, u, \gamma$, results in estimate:
$\mathcal{G}_{43, s}^{(1)}\lesssim W^{\delta_\sto/5}W^{-2}\eta_s^{-1}\times $\begin{tikzpicture}[
    thick, % Line thickness
    wavy/.style={decorate, decoration={snake, amplitude=2pt, segment length=5pt}},
    ]
\node (a) at (0,0) {a};
\node (x) at (1, 0) {x};
\node (gamma) at (2, 0) {$\gamma$};
\node (delta) at (3, 0) {$\delta$};
\node (y) at (4, 0) {y};
\node (b) at (5, 0) {b};

\draw[thick, double] (a) -- (x);
\draw[thick, double] (y) -- (b);
\draw[red, ->](x) to (gamma);
\draw[wavy, orange] (gamma) to (delta);
\draw[red, ->] (delta) to (y);

\end{tikzpicture}



\subsubsection*{Size estimate for $\mathcal{G}_{43, s}^{(2)}$}


\begin{center}
\begin{tikzpicture}[
    thick, % Line thickness
    wavy/.style={decorate, decoration={snake, amplitude=2pt, segment length=5pt}},
    ]
\node (a) at (0,0) {a};
\node (x) at (2, 0) {x};
\node (delta) at (4, 0) {$\delta$};
\node (y) at (6, 0) {y};
\node (b) at (8, 0) {b};
\node (gamma) at (4, 1) {$\gamma$};
\node (u) at (2, 1) {u};
\node (v) at (2, 2) {v};
\node (alpha) at (4, 2) {$\alpha$};
\node (beta) at (6, 2) {$\beta$};


\draw[thick, double] (a) -- (x);
\draw[wavy] (u) -- (v);
\draw[blue, wavy] (gamma) to (u);
\draw[thick, double] (y) -- (b);
\draw[wavy, blue] (v) to (alpha);
\draw[wavy](gamma) to (delta);
\draw[wavy] (alpha) to (beta);
\draw[blue, ->] (gamma) to(alpha);
\draw[blue, ->] (x) to (gamma);
\draw[blue, ->] (delta) to [bend left=15](beta);
\draw[blue, <-] (delta) to [bend right=15](beta);
\draw[red, ->] (x) to [bend right=25](y);
\draw[blue, ->] (alpha) to (y);
\end{tikzpicture}
\end{center}
\begin{equation*}
\mathcal{G}_{43, s}^{(2)} =  \left[\{\text{Id} + (t-s)\Theta_t\}S^{1/2}\right]_{ax}
\overline{G_{x y}}S_{uv} B_{v\alpha}S_{\alpha \beta} G_{\beta \delta}G_{\gamma \alpha}G_{\alpha y} B_{u\gamma} S_{\gamma\delta} G_{x\gamma}G_{\delta \beta}
\left[\{\text{Id} + (t-s)\Theta_t\}S^{1/2}\right]_{yb}
\end{equation*}

\noindent Bounding $\gamma\textcolor{blue}{\rightarrow}\alpha$, $\delta\textcolor{blue}{\rightarrow}\beta$, $\beta\textcolor{blue}{\rightarrow}\delta$ and $x\textcolor{red}{\rightarrow}y$ for a combined $O(W^{\delta_\sto/5}W^{-2}\eta_s^{-1})$ term, and summing over $\beta, \alpha, v, u$, yields the all-too familiar bound:
$\mathcal{G}_{43, s}^{(2)}\lesssim W^{\delta_\sto/5} W^{-2}\eta_s^{-1}\times $
\begin{tikzpicture}[
    thick, % Line thickness
    wavy/.style={decorate, decoration={snake, amplitude=2pt, segment length=5pt}},
    ]
\node (a) at (0,0) {a};
\node (x) at (1, 0) {x};
\node (gamma) at (2, 0) {$\gamma$};
\node (alpha) at (3, 0){$\alpha$};
\node (y) at (4, 0) {y};
\node (b) at (5, 0) {b};

\draw[thick, double] (a) -- (x);
\draw[thick, double] (y) -- (b);
\draw[blue, ->](x) to (gamma);
\draw [wavy, orange] (gamma) to (alpha);
\draw[blue, ->] (alpha) to (y);

\end{tikzpicture}


\subsubsection*{Size estimate for $\mathcal{G}_{43, s}^{(3)}$}


\begin{center}
\begin{tikzpicture}[
    thick, % Line thickness
    wavy/.style={decorate, decoration={snake, amplitude=2pt, segment length=5pt}},
    ]
\node (a) at (0,0) {a};
\node (x) at (2, 0) {x};
\node (gamma) at (4, 0) {$\gamma$};
\node (u) at (4, 1) {u};
\node (delta) at (5, 1) {$\delta$};
\node (y) at (6, 0) {y};
\node (b) at (8, 0) {b};
\node (alpha) at (4, 2) {$\alpha$};
\node (beta) at (6, 2) {$\beta$};
\node (v) at (3, 1) {v};

\draw[thick, double] (a) -- (x);
\draw[wavy] (u) -- (v);
\draw[blue, wavy] (gamma) to (u);
\draw[thick, double] (y) -- (b);
\draw[wavy, blue] (v) to (alpha);
\draw[wavy](gamma) to (delta);
\draw[wavy] (alpha) to (beta);
\draw[blue, ->] (x) to(gamma);
\draw[red, ->] (x) to [bend right=30](y);
\draw[blue, <-] (delta) to (alpha);
\draw[blue, ->] (delta) to (beta);
\draw[blue, ->] (gamma) to (y);
\draw[blue, <-] (alpha) to [bend left=30] (beta);
\end{tikzpicture}
\end{center}
\begin{equation*}
\mathcal{G}_{43, s}^{(3)} =   \left[\{\text{Id} + (t-s)\Theta_t\}S^{1/2}\right]_{ax}
\overline{G_{x y}}S_{uv} B_{v\alpha}S_{\alpha \beta} G_{\beta \alpha}G_{\alpha \delta}G_{\gamma y} B_{u\gamma} S_{\gamma\delta} G_{x\gamma}G_{\delta \beta}
\left[\{\text{Id} + (t-s)\Theta_t\}S^{1/2}\right]_{yb}
\end{equation*}
Here, we bound the edges $\alpha\textcolor{blue}{\rightarrow}\delta$, $x\textcolor{red}{\rightarrow}y$, $\beta\textcolor{blue}{\rightarrow}\alpha$,  and $\delta\textcolor{blue}{\rightarrow}\beta$ for a combined term $O(W^{\delta_\sto/5}W^{-2}\eta_s^{-1})$. After summing out $\beta, \alpha, v, u, \delta$ in that order, we get 
$\mathcal{G}_{43, s}^{(3)}\lesssim W^{\delta_\sto/5}W^{-2}\eta_s^{-1}\times $\begin{tikzpicture}[
    thick, % Line thickness
    wavy/.style={decorate, decoration={snake, amplitude=2pt, segment length=5pt}},
    ]
\node (a) at (0,0) {a};
\node (x) at (1, 0) {x};
\node (gamma) at (2, 0) {$\gamma$};
\node (y) at (3, 0) {y};
\node (b) at (4, 0) {b};

\draw[thick, double] (a) -- (x);
\draw[thick, double] (y) -- (b);
\draw[blue, ->](x) to (gamma);
\draw[blue, ->] (gamma) to (y);

\end{tikzpicture}. With this, we have verified that all the partial derivative terms satisfy the estimate $O(W^{-\frac{3}{4}}\eta_t^{-1}W^{-1}\eta_t^{-\frac{1}{2}})$. 
\subsubsection*{Combined bound}
All of the latter graph expansions allowed us to confirm the following result: $$\max_{a,b}\max_{i=1, \cdots, 4}\max_{j=0, 1, 2, 3} \left\vert \int_0^{t\wedge \tau_\sto}\mathcal{G}_{ij, s}(z)_{ab}ds\right\vert\lesssim W^{-\frac{3}{4}}\eta_t^{-1}W^{-1}\eta_t^{-\frac{1}{2}}$$
It is important to note that so far we have been working with the loop and regular vertex expansions without the regularization terms. They require more technical tools that work with isotropic bounds on graph modifications, and as such are outside the scope of this thesis. We simply state that the implication of their result (Lemma 22 \cite{bandSDE}) is within the stochastic bound we established above. As such, we conclude that when we integrate the graph \hyperref[graph]{$(6)$} over $s\in [0, t\wedge \tau_\sto]$, we get $\prec W^{-\frac{3}{4}}\eta_t^{-1}W^{-1}\eta_t^{-\frac{1}{2}}$, since all of the component terms in its sum expansion are of that order. Observe that everything we did applies WLOG to its conjguate, hence: $$\max_{a, b}\left\vert \int_0^{\tau_\sto\wedge t} \{\text{Id} + (t-s)\Theta_t\} S^{1/2}\Omega_s(z)S^{1/2} \{\text{Id} + (t-s)\Theta_t\}ds\right\vert \prec W^{-\frac{3}{4}}\eta_t^{-1}W^{-1}\eta_t^{-\frac{1}{2}}$$
The authors then argue that by virtue of "naive polynomial-in-W bounds" for the LHS and a "standard net argument" \cite{bandSDE} one can conclude:\begin{boxtheorem}{Theorem 6.1}
$$\max_{t\in [0, 1]}\max_{a, b}\frac{\left\vert \int_0^{\tau_\sto\wedge t} \{\text{Id} + (t-s)\Theta_t\} S^{1/2}\Omega_s(z)S^{1/2} \{\text{Id} + (t-s)\Theta_t\}ds\right\vert }{W^{-\frac{3}{4}}\eta_t^{-1}W^{-1}\eta_t^{-\frac{1}{2}}}\prec 1$$\end{boxtheorem}
While the technical details are again outside the scope of this thesis, we can elaborate on the meaning of both those approaches. The first one refers to straightforward power counting of each matrix component in the integrand. Specifically, $[\text{Id} + (t - s)\Theta_t]$ contributes at most $O(1)$ since $t,s \in [0,1]$ and $\Theta_t$ has bounded entries. The matrices $S_t^{\frac{1}{2}}$, $\Omega_s(z)$, and $S_s^{\frac{1}{2}}$ each contribute additional powers of $W$, but their product maintains an overall scaling of $O(W^C)$ where $C\lesssim 1$ does not depend on $W$. The net argument, on the other hand, provides control over the supremum by discretizing the integration interval $[0, \tau_{\text{stop}} \wedge t]$ into an $\varepsilon$-net with spacing $\delta = W^{-D}$ for some appropriate constant $D$. As such, we get the aforementioned $O(W^C)$ bound, and by continuity of the integrand (with Lipschitz constant also polynomially bounded in $W$), these bounds extend to the entire interval with negligible error. When normalized by the denominator $W^{-\frac{3}{4}}\eta_t^{-1} \cdot W^{-1}\eta_t^{-\frac{1}{2}}$, the ratio remains bounded by a constant, so $\prec 1$. 




\subsection{Estimate for $\E_t^{M, \text{stop}}(z)$}
Recall that our martingale term was the sum of two quadratic variations, namely:
$$\E_t^{M, \text{stop}}(z)=-\int_0^{t\wedge \tau_\sto} \{\text{Id} + (t-s)\Theta_t\} S^{1/2}dM_s(z)S^{1/2} \{\text{Id} + (t-s)\Theta_t\}= $$$$=-\int_0^{t\wedge \tau_\sto} \{\text{Id} + (t-s)\Theta_t\}S^{\frac{1}{2}}\sum_{u, v}\left[\overline{G_{xy}}G_{xu}dH_{uv}(t)G_{vy}\right]S^{\frac{1}{2}}\{\text{Id} + (t-s)\Theta_t\}-$$
$$- \int_0^{t\wedge \tau_\sto} \{\text{Id} + (t-s)\Theta_t\}S^{\frac{1}{2}}\sum_{u, v}\left[G_{xy}\overline{G_{xu}dH_{uv}(t)}\,\overline{G_{vy}}\right]S^{\frac{1}{2}}\{\text{Id} + (t-s)\Theta_t\}=$$
$$=\E_t^{M, 1}(z)+\E_t^{M, 2}(z),$$

\noindent which can in turn be written diagrammatically using the same notation as before (\hyperref[def2]{Def. 2}):
\begin{center}
\begin{tikzpicture}[
    thick, % Line thickness
    wavy/.style={decorate, decoration={snake, amplitude=2pt, segment length=5pt}},
    ]
\node (a) at (0,0) {a};
\node (x) at (1, 0) {x};
\node (y) at (3, 0) {y};
\node (b) at (4, 0) {b};
\node (u) at (1.5, 1) {u};
\node (v) at (2.5, 1){v};

\draw[thick, double] (a) -- (x);
\draw[thick, double] (y) -- (b);
\draw[red, ->](x) to (y);
\draw[blue, ->] (x) to (u);
\draw[blue, ->] (v) to (y);
\draw[wavy, purple] (u) to (v);

\end{tikzpicture} \begin{tikzpicture}[
    thick, % Line thickness
    wavy/.style={decorate, decoration={snake, amplitude=2pt, segment length=5pt}},
    ]
\node (a) at (0,0) {a};
\node (x) at (1, 0) {x};
\node (y) at (3, 0) {y};
\node (b) at (4, 0) {b};
\node (u) at (1.5, 1) {u};
\node (v) at (2.5, 1){v};

\draw[thick, double] (a) -- (x);
\draw[thick, double] (y) -- (b);
\draw[blue, ->](x) to (y);
\draw[red, ->] (x) to (u);
\draw[red, ->] (v) to (y);
\draw[wavy, red] (u) to (v);

\end{tikzpicture},
\end{center}
where the purple wavy line represents $dH_{uv}$ and the red wavy line is its complement. In order to estimate the latter term, we will apply the Burkholder-Davis-Gundy (BGD) inequality \cite{bdg}, defined below:
\begin{boxtheorem}{\centering BDG inequality}
$$\mathbb{E}\left[\sup_{0\leq s\leq t}\left\vert M_{xy}(s)\right\vert^{2p}\right]\leq C_p\mathbb{E}\left[ M_{xy} \right]_t^p$$
\end{boxtheorem}
In order to express the Quadratic Variation (QV) on the RHS, we will need to rewrite our separate graphs into one combined QV graph. To do this, we will use that fact that the Brownian increments $dB_{\alpha\beta}(t)$ are mutually uncorrelated, s.t. the cross-terms in the quadratic variation vanish unless $(u, v) = (u', v')$. This means that we can identify the latter pairs, getting a combined graph:


\begin{center}
\begin{tikzpicture}[
    thick, % Line thickness
    wavy/.style={decorate, decoration={snake, amplitude=2pt, segment length=5pt}},
    ]
\node (a) at (0,0) {a};
\node (xp) at (1, -1) {x'};
\node (yp) at (3, -1) {y'};
\node (b) at (4, 0) {b};
\node (u) at (1.5, 0) {u};
\node (v) at (2.5, 0){v};
\node (x) at (1, 1){x};
\node (y) at (3, 1){y};

\draw[thick, double] (a) -- (xp);
\draw[thick, double] (yp) -- (b);
\draw[blue, ->](xp) to (yp);
\draw[red, ->] (xp) to (u);
\draw[red, ->] (v) to (yp);
\draw[wavy, red] (u) to (v);
\draw[thick, double] (a) -- (x);
\draw[thick, double] (y) -- (b);
\draw[red, ->](x) to (y);
\draw[blue, ->] (x) to (u);
\draw[blue, ->] (v) to (y);

\end{tikzpicture}
\end{center}
If we are to approach the bounding in the same fashion as before, we will not get appropriate bounds. To see this, let us start by eliminating the wavy line for a term  $\sqrt{S_{uv}}\lesssim W^{-1/2}$ per estimate $(E4)$, along with $O(\sqrt{dt})$ per the Brownian variance $\mathbb{E}\left\vert dB_{uv}(t)\right\vert^2=dt$. We want to keep a connected component similar to before, so let WLOG this be the path $x'=\!= a=\!= x\textcolor{blue}{\longrightarrow }y=\!= b=\!= y'$. By applying estimate $(E2)$ to all straight lines connected to $u$ and $v$, we get a combined term $O(W^{\delta_\sto/5}W^{-2}\eta_s^{-1})$. To get to the the aforementioned path, we need to bound the $x'\textcolor{blue}{\rightarrow}y'$ using estimate $(1)$ (since the edge is not "free") with which we get a term $O(W^{\delta_\sto/20}W^{1/2}\eta_s^{-3/4})$. For what remains we will have to split three double lines into their respective $\text{Id}$ and $(t-s)\Theta_t$ components, yielding $2^3$ separate terms, each of the form $O\left(W^{\delta_{ab}}W^{-\frac{3}{2}}\eta_t^{-\frac{3a+2b}{2}}\eta_{s}^{\frac{5a+3b}{4}}\right)$, where $a, b\in \{0, 1, 2, 3\}$,  $a+b=3$, and $\delta_{ab}>0$ is $\lesssim\delta_\sto$. We can easily observe that this gives us at most $W^{-\frac{3}{2}}$, which is not enough for at least a coefficient of order $W^{-\frac{7}{4}}$ in the bound, necessary for our proof. As such, we will have to approach $\E_t^{M, \sto}$ differently:


\subsubsection*{Alternative bound}
The latter calculations demonstrate that we need to be more thoughtful with our estimates. Let us come back to our original form of the $S^{1/2}dM_s(z)S^{1/2}$ term. Observe that we can split the components $\E_t^{M, 1}$ and $\E_t^{M, 2}$ into 4 separate parts each by expanding the square brackets. We will this explicitly for $\mathcal{E}_{t}^{M, 1}$, whereas the second term is equivalent up to complex conjugation:

$$\mathcal{E}_t^{M, 1} = -\int_0^{t\wedge \tau_\sto}(t-s)^2\Theta_tS^{1/2}\left\{\sum_{u, v}\left[\overline{G_{xy}}G_{xu}dH_{uv}(t)G_{vy}\right]\right\}S^{1/2}\Theta_t+$$
$$-\int_0^{t\wedge \tau_\sto}(t-s)\Theta_tS^{1/2}\left\{\sum_{u, v}\left[\overline{G_{xy}}G_{xu}dH_{uv}(t)G_{vy}\right]\right\}S^{1/2}$$
$$-\int_0^{t\wedge \tau_\sto}(t-s)S^{1/2}\left\{\sum_{u, v}\left[\overline{G_{xy}}G_{xu}dH_{uv}(t)G_{vy}\right]\right\}S^{1/2}\Theta_t+$$
$$-\int_0^{t\wedge \tau_\sto}S^{1/2}\left\{\sum_{u, v}\left[\overline{G_{xy}}G_{xu}dH_{uv}(t)G_{vy}\right]\right\}S^{1/2} =$$
$$= \mathcal{E}_t^{M, 11}+\mathcal{E}_t^{M, 12}+\mathcal{E}_t^{M, 13}+\mathcal{E}_t^{M, 14}$$
Now, we will apply the BDG inequality to the individual $\mathcal{E}_t^{M, 1i}$ terms for $j=1, 2, 3, 4$:

\begin{equation*}\mathbb{E}\left\vert \mathcal{E}^{M, 1i}_t(z)_{ab}\right\vert^{2p}\leq C_p\mathbb{E}\left[ \mathcal{E}^{M, 1i}_t(z)_{ab} \right]^p\tag{6.3}\end{equation*}


\noindent To calculate the quadratic variation (QV) on the RHS, recall from before that the Brownian increments $dB_{\alpha\beta}(t)$ are mutually uncorrelated, i.e our term $dH_s = \sum_{u, v} \sqrt{S_{u, v}}dB_{u, v}(s)$ has covariance structure: $$\mathbb{E}\left[dH_s\overline{dH_s}\right] = \mathbb{E}\left[\sum_{u, u', v, v'}\sqrt{S_{uv}}dB_{uv}(s)\sqrt{S_{u'v'}}dB_{u'v'}(s)\right]_{ab}=S_{ab}ds$$
Hence the QV of $\left[\mathcal{E}^{M, 11}\right]_{ab}$ and the other terms is:
$$ \int_0^{t\wedge \tau_\sto}(t-s)^4\sum_{x, x', y, y', u, v}(\Theta_tS^{1/2})_{ax}(\Theta_tS^{1/2})_{ax'}\overline{G_{xy}}G_{x'y'}G_{xu}\overline{G}_{x'u}S_{uv}G_{vy}\overline{G_{vy'}}(S^{1/2}\Theta_t)_{yb}(S^{1/2}\Theta_t)_{yb'}ds$$

$$\left[\mathcal{E}^{M, 12}\right]_{ab} = \int_0^{t\wedge \tau_\sto}(t-s)^2\sum_{x, x', y, y', u, v}(\Theta_tS^{1/2})_{ax}(\Theta_tS^{1/2})_{ax'}\overline{G_{xy}}G_{x'y'}G_{xu}\overline{G}_{x'u}S_{uv}G_{vy}\overline{G_{vy'}}S^{1/2}_{yb}S^{1/2}_{yb'}ds$$

$$\left[\mathcal{E}^{M, 13}\right]_{ab} = \int_0^{t\wedge \tau_\sto}(t-s)^2\sum_{x, x', y, y', u, v}S^{1/2}_{ax}S^{1/2}_{ax'}\overline{G_{xy}}G_{x'y'}G_{xu}\overline{G}_{x'u}S_{uv}G_{vy}\overline{G_{vy'}}(S^{1/2}\Theta_t)_{yb}(S^{1/2}\Theta_t)_{yb'}ds$$

$$\left[\mathcal{E}^{M, 14}\right]_{ab} = \int_0^{t\wedge \tau_\sto}\sum_{x, x', y, y', u, v}S^{1/2}_{ax}S^{1/2}_{ax'}\overline{G_{xy}}G_{x'y'}G_{xu}\overline{G}_{x'u}S_{uv}G_{vy}\overline{G_{vy'}}S^{1/2}_{yb}S^{1/2}_{yb'}ds$$

\noindent Observe that we have the same primary components, differing up to the utmost left and right terms. Hence, we can abstract away the following four separate matrices: 

$$\Upsilon_{yy}^u := \sum_v S_{uv} G_{vy} \overline{G}_{vy'} (\Theta_t S^{\frac{1}{2}})_{yb} (\Theta_t S^{\frac{1}{2}})_{y'b}$$
$$\Omega_{yy}^u := \sum_{x,x'} \overline{G}_{xy} G_{x'y'} G_{xu} \overline{G}_{x'u} (\Theta_t S^{\frac{1}{2}})_{ax} (\Theta_t S^{\frac{1}{2}})_{ax'}$$
$$\Gamma_{yy}^u := \sum_v S_{uv} G_{vy} \overline{G}_{vy'} S^{\frac{1}{2}}_{yb} S^{\frac{1}{2}}_{y'b}$$
$$\Xi_{yy}^u := \sum_{x,x'} \overline{G}_{xy} G_{x'y'} G_{xu} \overline{G}_{x'u} S^{\frac{1}{2}}_{ax} S^{\frac{1}{2}}_{ax'}$$
Now, observe that our terms become expressed simply as the two-by-two products:
\begin{equation*}
\begin{array}{ll}
 \left[\mathcal{E}_{t}^{M, 11}(z)_{ab}\right] = \int_0^{t\wedge \tau_\sto} (t-s)^4\sum_{u, y} (\Omega^u\Upsilon^{u, *})_{yy'}ds &  \left[\mathcal{E}_{t}^{M, 12}(z)_{ab}\right] = \int_0^{t\wedge \tau_\sto} (t-s)^2\sum_{u, y} (\Omega^u\Gamma^{u, *})_{yy'}ds \\
 \left[\mathcal{E}_{t}^{M, 13}(z)_{ab}\right] = \int_0^{t\wedge \tau_\sto} \sum_{u, y} (t-s)^2(\Xi^u\Upsilon^{u, *})_{yy'}ds  & \left[\mathcal{E}_{t}^{M, 14}(z)_{ab}\right] = \int_0^{t\wedge \tau_\sto} \sum_{u, y} (\Xi^u\Gamma^{u, *})_{yy'}ds
\end{array}\end{equation*}
Observe that both $\Omega^u$ and $\Xi^u$ are positive-semidefinite, since we can rewrite them as:
$$\Omega^u = \left(\sum_x (\Theta_t S^{1/2})_{ax}G_{xu}\overline{G}_{xy}\right)\left(\sum_{x'} (\Theta_t S^{1/2})_{ax'}G_{x'u}\overline{G}_{x'y'}\right)^*$$
$$\Xi^u = \left(\sum_x  S^{1/2}_{ax}G_{xu}\overline{G}_{xy}\right)\left(\sum_{x'} S^{1/2}_{ax'}G_{x'u}\overline{G}_{x'y'}\right)^*$$
This means that we can apply the von Neumann trace inequality, which states that the trace is bounded above by singular values $\left\vert\text{Tr}(AB)\right\vert\leq \sum_{i=1}^N \alpha_i\beta_i$, hence by the positive-semidefiniteness: 
$$\left\vert\sum_y\left(\Omega^u\Upsilon^{u, *}\right)_{yy}\right\vert = \left\vert\text{Tr}\left(\Omega^u\Upsilon^{u, *}\right)\right\vert\leq \left\|\Omega^u\right\|_{op}\sum_{y}\Omega_{yy'}^uds.$$
This means that we can bound each of the earlier terms as follows:

$$\left[\mathcal{E}_{t}^{M, 11}(z)_{ab}\right] = \int_0^{t\wedge \tau_\sto} (t-s)^4\sum_{u, y} (\Omega^u\Upsilon^{u, *})_{yy}ds\leq \int_0^{t\wedge \tau_\sto} (t-s)^4\sum_u \|\Upsilon^u\|_{op}\sum_y \Omega_{yy'}^uds$$

$$\left[\mathcal{E}_{t}^{M, 12}(z)_{ab}\right] = \int_0^{t\wedge \tau_\sto} (t-s)^2\sum_{u, y} (\Omega^u\Gamma^{u, *})_{yy'}ds\leq \int_0^{t\wedge \tau_\sto} (t-s)^2\sum_u \|\Gamma^u\|_{op}\sum_y \Omega_{yy'}^uds$$

$$\left[\mathcal{E}_{t}^{M, 13}(z)_{ab}\right] = \int_0^{t\wedge \tau_\sto} \sum_{u, y} (t-s)^2(\Xi^u\Upsilon^{u, *})_{yy'}ds \leq  \int_0^{t\wedge \tau_\sto} (t-s)^2\sum_u \|\Upsilon^u\|_{op}\sum_y \Xi_{yy'}^uds$$

$$\left[\mathcal{E}_{t}^{M, 14}(z)_{ab}\right] = \int_0^{t\wedge \tau_\sto} \sum_{u, y} (\Xi^u\Gamma^{u, *})_{yy'}ds\leq \int_0^{t\wedge \tau_\sto} \sum_u \|\Gamma^u\|_{op}\sum_y \Xi_{yy'}^uds$$
We can estimate the latter bounds by using the fact that $\vert G\vert^2 = GG^*$ and $|G|^4=|G|^2|G|^2$:
$$\sum_{u, y}\Omega_{yy}^u=\sum_{u,y}\sum_{x,x'} \overline{G}_{xy} G_{x'y'} G_{xu} \overline{G}_{x'u} (\Theta_t S^{\frac{1}{2}})_{ax} (\Theta_t S^{\frac{1}{2}})_{ax'} =$$
$$= \sum_{x, x'}|G|^2_{xx'}|{G}|^2_{x'x} (\Theta_t S^{\frac{1}{2}})_{ax} (\Theta_t S^{\frac{1}{2}})_{ax'}\leq 4 \sum_{x}|G|_{xx}^2 |(\Theta_t S^{\frac{1}{2}})_{ax}|^2,$$
following from the Schwarz inequality, since $|G|_{xx'}^2|G|^2_{x'x}=||G|_{xx'}^2|^2$ by virtue of $|G|^2$ being Hermitian. On the other hand, by the \hyperref[ward]{Ward identity}, $\max_{a, b}|(\Theta_t)_{ab}|\lesssim W^{-1}\eta_t^{-1/2}$ and $\sum_\beta |(\Theta_t S^{1/2})_{\alpha\beta}|\lesssim \eta_t^{-1}$ (\hyperref[lemma-a2]{Lemma A.2}), the stopping time $s\leq \tau_\sto$ and $(t-s)\lesssim \eta_s$ (\hyperref[lemma-a21]{Lemma A.1.2}) , we get: $$4 \sum_{x}|G|_{xx}^4 |(\Theta_t S^{\frac{1}{2}})_{ax}|^2=4 \eta_s^{-2}\sum_x|\text{Im}G_{xx}|^2|(\Theta_t S^{\frac{1}{2}})_{ax}|^2\leq 4 \eta_s^{-2}|\max_{k}\text{Im}G_{kk}|^2\sum_x|(\Theta_t S^{\frac{1}{2}})_{ax}|^2$$
$$\lesssim\eta_s^{-3}\max_k \text{Im}G_{kk}\sum_{x} |(\Theta_t S^{\frac{1}{2}})_{ax}|^2\leq W^{\delta_\sto}\eta_s^{-3}W^{-1}\eta_t^{-3/2},$$

\noindent By the same logic as for $\Omega^u_{yy}$, along with the fact that $S_{ab}^{1/2}\lesssim W^{-1}\Rightarrow \sum_b S_{ab}^{1/2} \lesssim 1$, we have that: 

$$\sum_{u, y}\Xi_{yy}^u=\sum_{u,y}\sum_{x,x'} \overline{G}_{xy} G_{x'y'} G_{xu} \overline{G}_{x'u} S^{\frac{1}{2}}_{ax} S^{\frac{1}{2}}_{ax'} = \sum_{x, x'}|G|^2_{xx'}|{G}|^2_{x'x} S^{\frac{1}{2}}_{ax} S^{\frac{1}{2}}_{ax'}\leq $$
$$\leq 4 \sum_{x}|G|_{xx}^4 |S^{\frac{1}{2}}_{ax}|^2\lesssim \eta_s^{-3}\max_k\text{Im}G_{kk}\sum_x|S_{ax}^{1/2}|^2\lesssim W^{\frac{\delta_\sto}{10}}W^{-1}\eta_s^{-3}$$

\noindent We need to bound the other two matrices,. We can do this by using the fact that $S^u$ is the diagonal matrix $S^u_{ij} = \delta_{ij}S_{ui} = O(W^{-1})$, along with estimate $(E4)$ and the fact that $\|G\|_{op} = \sup_{\lambda_i}\frac{1}{|\lambda -z|}  = \eta_s^{-1}$,
hence we get the following bounds:
$$\|\Upsilon^u\|_{op}\leq \max_{\alpha, \beta}|(\Theta_t S^{1/2})|^2\|G^* S^u G\|_{op}\leq \max_{\alpha, \beta}|(\Theta_t S^{1/2})|^2 W^{-1}\eta_s^{-2}$$
$$\|\Gamma^u\|_{op}\leq\max_{\alpha,\beta}|S_{\alpha\beta}^{1/2}|^2\|G^* S^u G\|_{op}\lesssim \max_{\alpha,\beta}|S_{\alpha\beta}^{1/2}|^2 W^{-1}\eta_s^{-2}\lesssim W^{-3}\eta_s^{-2}$$

\noindent With this we have all the necessary pieces to complete the estimation of all the QV terms $\mathcal{E}^{M, 1i}_t$, by changing the variable of integration w.r.t $\sigma = |\text{Im}w_s|$. As such, for any $\delta>0$:
$$\left[\mathcal{E}_{t}^{M, 11}(z)_{ab}\right]\lesssim \int_0^{\tau_\sto\wedge t} \eta_s^{-1} W^{\delta_\sto}W^{-2}\eta_t^{-\frac{3}{2}}\max_{\alpha, \beta}|(\Theta_t S^{1/2})_{\alpha, \beta}|^2ds$$
$$\left[\mathcal{E}_{t}^{M, 11}(z)_{ab}\right]\lesssim W^{\delta_\sto}W^{-2}\eta_t^{-\frac{3}{2}} \max_{\alpha, \beta}|(\Theta_t S^{1/2})_{\alpha, \beta}|^2 \int_{\eta_t}^{\eta_0} \sigma^{-1}ds\lesssim W^{-4+\delta+\delta_\sto}\eta_t^{-\frac{5}{2}},$$
since $\max_{\alpha, \beta}|(\Theta_t S^{1/2})_{\alpha, \beta}|^2 \lesssim W^{-2}\eta_t^{-1}$ (\hyperref[lemma-a2]{Lemma A.2}) and $\log \eta_t^{-1}\lesssim \log \eta^{-1}\lesssim W^{\delta}$. Similarly, the second and third terms are: 

$$\left[\mathcal{E}_{t}^{M, 12}(z)_{ab}\right] \lesssim  W^{\delta_\sto} W^{-3}W^{-1}\eta_t^{-3/2}\int_0^{t\wedge \tau_\sto}(t-s)^2\eta_s^{-5}ds$$
$$\lesssim W^{\delta_\sto}W^{-4}\eta_t^{-3/2}\int_{\eta_t}^{\eta_0}\sigma^{-3}ds\lesssim W^{\delta_\sto} W^{-4}\eta_t^{-\frac{7}{2}}$$

$$\left[\mathcal{E}_{t}^{M, 13}(z)_{ab}\right] \lesssim W^{\frac{\delta_\sto}{10}}W^{-2} \max_{\alpha, \beta}|(\Theta_t S^{1/2})_{\alpha, \beta}|^2 \int_0^{\tau_\sto\wedge t} (t-s)^2 \eta_s^{-5}ds\lesssim $$
$$\lesssim W^{\frac{\delta_\sto}{10}}W^{-2} \max_{\alpha, \beta}|(\Theta_t S^{1/2})_{\alpha, \beta}|^2 \int_{\eta_t}^{\eta_0}\sigma^{-3}ds\lesssim W^{\frac{\delta_\sto}{10}}W^{-4} \eta_t^{-3}$$
As for the last term: $$\left[\mathcal{E}_{t}^{M, 14}(z)_{ab}\right] \leq \int_0^{t\wedge \tau_\sto} \sum_u \|\Gamma^u\|_{op}\sum_y \Xi_{yy'}^uds\lesssim \int_0^{t\wedge \tau_\sto} W^{-1}\eta_s^{-3}W^{\delta_\sto/10}W^{-3}\eta_s^{-2}$$
$$\lesssim W^{\frac{\delta_\sto}{10}-4}\int_{\eta_t}^{\eta_0}\sigma^{-5}\lesssim W^{\delta_\sto/10}W^{-4}\eta_t^{-4},$$
we can observe that a naive bound of the sort we did for the last three terms is not sufficient ($\eta_s^{-4}$ instead of $\eta_s^{-3}$) , given that we do not have $(t-s)$ multiples. For this reason, the authors instead approach this term with an interpolation with a second bound \cite{bandSDE}. First, recall that:
$$\E_t^{M, 14}(z)= \int_0^{t\wedge \tau_\sto}S^{1/2}\left\{\sum_{u, v}\left[\overline{G_{xy}}G_{xu}dH_{uv}(t)G_{vy}\right]\right\}S^{1/2}$$
By the fact that $S^{1/2}$ is an averaging operator (\hyperref[diffusion]{Sec. 4.1}), one can estimate: $$\left[\E_t^{M, 14}(z)\right] \leq \max_{\alpha, \beta}\left[\left\{\int_0^{\tau_\sto\wedge}S^{1/2} \sum_{u, v}\left[\overline{G_{xy}}G_{xu}dH_{uv}(t)G_{vy}\right]\right\}_{\alpha\beta}\right]=$$

\begin{equation*}=\max_{\alpha\beta}\int_0^{\tau_\sto\wedge}\sum_{x, x', u, v}S^{\frac{1}{2}}_{\alpha x}S^{\frac{1}{2}}_{\alpha x'}\overline{G_{x\beta}}G_{x'\beta}G_{xu}\overline{G_{x'u}}S_{uv}|G_{v\beta}|^2ds=\tag{6.4}\end{equation*}
\begin{equation*}=\max_{\alpha\beta}\int_0^{\tau_\sto\wedge}\sum_{x, x', u, v}\left(G^*S^{\frac{1}{2}, \alpha}GG^*S^{\frac{1}{2},\alpha}G\right)_{uu}S_{uv}|G_{v\beta}|^2ds\tag{6.4.1}\end{equation*}
By first applying the \hyperref[ward]{Ward identity} and then the operator norm bound $\|G\|_{op}\leq \eta_s^{-1}$, we have: 
$$\left(G^*S^{\frac{1}{2}, \alpha}GG^*S^{\frac{1}{2},\alpha}G\right)_{uu} = \eta_s^{-1}\left(G^*S^{\frac{1}{2}, \alpha}\text{Im}GS^{\frac{1}{2},\alpha}G\right)_{uu}\lesssim$$
$$\lesssim \eta_s^{-1}\left\|\sqrt{S^{\frac{1}{2}, \alpha}}\im G\sqrt{S^{\frac{1}{2}, \alpha}}\right\|_{\mathrm{op}}(G^*S^{\frac{1}{2}, \alpha}G)_{uu}\lesssim$$
\begin{equation*}\lesssim W^{-1}\eta_s^{-2}(G^*S^{\frac{1}{2}, \alpha}G)_{uu} = W^{-1}\eta_s^{-2}\sum_\gamma S^{\frac{1}{2}}_{\alpha\gamma}|G_{\gamma\beta}|^2\tag{6.4.2}\end{equation*}
Now, in both the display after $(6.4)$ and the one above we have terms of the form $S_{xy}|G_{yz}|^2$. To estimate this, we will use the a time $s\leq \tau_\sto$ before the stopping time, for which we know \cite{bandSDE}: $$|T_s(z)_{zy}|\lesssim |(\Theta_s)_{zy}|+|\E_s(z)_{xy}|\lesssim W^{-1}\eta_s^{-\frac{1}{2}}$$
$$\Rightarrow \sum_v S_{uv}|G_{v\beta}|^2\lesssim W^{\delta_\sto/10}S_{u\beta}+W^{\delta_\sto/10}\left[S^{1/2}T_s(z)S^{1/2}\right]_{u\beta}\lesssim$$ 
$$\lesssim W^{\delta_\sto}S_{u\beta} + W^{\delta_\sto/10}W^{-1}\eta_s^{-1/2}\lesssim W^{\delta_\sto/10}W^{-1}\eta_s^{-1/2},$$
since $\eta_s = O(1)$. This means that we can combine $(6.4.1)$ and $(6.4.2)$ along with the result above in order to integrate with a change of variable as before: \newpage
$$\left[\E_t^{M, 14}(z)_{ab}\right]\lesssim \int_0^{t\wedge \tau_\sto} W^{\delta_\sto/5}W^{-3}\eta_s^{-3}ds\lesssim$$
$$\lesssim W^{\delta_\sto/5}W^{-3}\int_{\eta_t}^{\eta_0}\sigma^{-3}d\sigma\lesssim W^{\delta_\sto/5}W^{-3}\eta_t^{-2}$$

\noindent By interpolating with the first bound we found, we get: 
$$\left[\E_t^{M, 14}(z)_{ab}\right] \lesssim \sqrt{ W^{\delta_\sto/10}W^{-4}\eta_t^{-4}}\sqrt{W^{\delta_\sto/5}W^{-3}\eta_t^{-2}}\lesssim W^{\frac{3\delta_\sto}{20}}W^{-\frac{7}{2}}\eta_t^{-3}\lesssim W^{\frac{\delta_\sto}{5}}W^{-\frac{7}{2}}\eta_t^{-3}$$
As such, we have in fact showed that for all $i=1, \cdots 4$ $$[\E_t^{M, 1j}(z)_{ab}]\lesssim W^{\frac{\delta_\sto}{5}}\lesssim W^{\delta_\sto/5}W^{-\frac{3}{2}}\eta_t^{-2}W^{-2}\eta_t^{-1}$$
By applying Chebyshev to the BDG inequality $(6.3)$, we have that for any $\delta>0$ and $p\geq 1$: 
\begin{equation*}\mathbb{P}\left(\left\vert \E_t^{M, 1j}(z)_{ab}\right\vert\geq N^\delta [\E_t^{M, 1j}(z)_{ab}]^{1/2}\right)\leq C_p N^{-2p\delta}\tag{6.5}\end{equation*}
Using Boole's inequality and a "standard Hölder continuity argument"  are sufficient to show \cite{bandSDE}:
\begin{boxtheorem}{Theorem 6.2}
$$\frac{\left\vert \E_t^{M_1, j}(z)_{ab}\right\vert}{W^{-\frac{3}{4}}\eta_t^{-1}W^{-1}\eta_t^{-\frac{1}{2}}}\prec W^{\frac{\delta_\sto}{5}}$$
\end{boxtheorem}



\subsection{Estimate for $\E_t^{S, \text{stop}}(z)$}

Recall \hyperref[cor-a3]{Corollary A.3} of \hyperref[lemma-a2]{Lemma A.2}: $$\sup_x \sum_y \{\text{Id}+(t-s)\Theta_t\}_{xy}+\sup_y \sum_x \{\text{Id}+(t-s)\Theta_t\}_{xy} = 1+O(\eta_t^{-1}\eta_s)$$
Then by Hölder's inequality and the bound above, we have: $$|\E_t^{S}(z)_{ab}|=\left\vert\int_0^{t\wedge \tau_\sto} \{\text{Id} + (t-s)\Theta_t\}\E_s^2(z) \{\text{Id} + (t-s)\Theta_t\}ds\right\vert \leq$$
$$\leq \int \left[1+O(\eta_t^{-1}\eta_s)\right]^2 \max_{a, b}|\E_s(z)_{ab}^2|ds \lesssim \int_0^{\tau_\sto \wedge t} \max_{a, b}|\E_s(z)_{ab}^2|ds + \int_0^{\tau_\sto \wedge t}\eta_t^{-2}\eta_s^{2} \max_{a, b}|\E_s(z)^2_{ab}|ds$$
We can bound $\max_{a, b}|\E_s(z)_{ab}^2|ds$ by using the definition of our stopping time. Since $s\leq \tau_\sto = \tau_{\sto, 1}\wedge \tau_{\sto, 2}$, by property of $\tau_{\sto, 1}$ and given a bound $N\leq W^{\frac{11}{8}-\nu}$ we have: $$\max_{a, b}|\E_s(z)_{ab}|<W^{\delta_\sto} W^{-\frac{3}{4}} \eta_s^{-1}\cdot W^{-1}\eta_s^{-\frac{1}{2}}$$
$$\max_{a, b}|\E_s(z)_{ab}|^2<W^{2\delta_\sto} W^{-\frac{3}{2}} \eta_s^{-2}\cdot W^{-2}\eta_s^{-1}=W^{2\delta_{\sto}} W^{-\frac{7}{2}}\eta_s^{-3} $$
Now, observe that by matrix multiplication, for any $a, b\in \Gamma$, we have $\E_s(z)^2_{ab} = \sum_{j=1}^N \E_s(z)_{aj}\E_s(z)_{jb}$:
$$\Rightarrow |\E_s(z)_{ab}^2|\leq \sum_{j=1}^N |\E_s(z)_{aj}||\E_s(z)_{jb}|\leq \sum_{j=1}^N \max_{x,y}|\E_s(z)_{xy}|^2 = N\max_{x,y}|\E_s(z)_{xy}|^2$$
$$\Rightarrow \max_{a,b}|\E_s(z)^2_{ab}|\leq W^{2\delta_\sto}W^{-\frac{28-11}{8}-\nu}\eta_s^{-2} = W^{2\delta_\sto}W^{-\frac{17}{8}-\nu}\eta_s^{-2}$$
Putting the last two results together gives us the following bound on $|\E_t^{S, \mathrm{stop}}(z)_{ab}|$:
$$|\E_t^{S, \mathrm{stop}}(z)_{ab}|\lesssim \int_0^t W^{2\delta_\sto}W^{-\frac{17}{8}-\nu}\eta_s^{-2}ds + \int_0^t \eta_t^{-2}W^{2\delta_\sto}W^{-\frac{17}{8}-\nu}\eta_s^{-1}ds$$
\begin{equation*}
|\E_t^{S, \sto}(z)_{ab}|\lesssim W^{2\delta_\sto} W^{-\frac{17}{8}-\nu}\int_{\eta_t}^{\im w_0}\sigma^{-3}d\sigma + \eta_t^{-2} W^{2\delta_\sto} W^{-\frac{17}{8}-\nu}\int_{\eta_t}^{\im w_0}\sigma^{-1}d\sigma
\end{equation*}
\begin{equation*}
\lesssim W^{2\delta_\sto} W^{-\frac{17}{8}-\nu} \left(\frac{\im w_0^{-2}-\eta_t^{-2}}{-2}\right)+W^{2\delta_\sto}W^{-\frac{17}{8}-\nu}\eta_t^{-2}\left(\ln \im w_0-\ln \eta_t\right)
\end{equation*}
By \hyperref[lemma-a12]{Lemma A.1.2} for any $s\in [0, t]$, we have $\im w_s = \im w_t+(t-s)\im m(z)\geq \im w_t$
and \hyperref[2.1]{$\im m(z)\asymp 1$}, which imply $\eta\leq \im w_s\lesssim 1$. And since $|\log \eta|\prec 1$ (follows by $\eta\geq W^{-3/4}$), we have that $|\log \im w_s|\prec 1$. Hence, we can drop the negative term w.r.t. $\im w_0^{-2}$ and bound above :
\begin{equation*}
|\E_t^{S, \sto}(z)_{ab}|\lesssim W^{2\delta_\sto} W^{-\frac{17}{8}-\nu} \eta_t^{-2}+W^{2\delta_\sto}W^{-\frac{17}{8}-\nu}\eta_t^{-2}\log \eta_t^{-1} \lesssim W^{2\delta_\sto}W^{-\frac{17}{8}-\nu}\eta_t^{-2}\log \eta_t^{-1}
\end{equation*}
Observe that the latter expression can be split in the following terms: $$|\E_t^{S, \sto}(z)_{ab}|\lesssim W^{2\delta_\sto} W^{-\frac{3}{8}-\nu} \eta_t^{-\frac{1}{2}}\log \eta_t^{-1}\cdot W^{-\frac{3}{4}}\eta_t^{-1}\cdot W^{-1}\eta_t^{-\frac{1}{2}}$$
In this context, since $\eta_t \geq \eta\geq W^{-3/4}$, the first term is: $$W^{2\delta_\sto} W^{-\frac{3}{8}-\eta}\eta_t^{-\frac{1}{2}}\log \eta_t^{-1}\lesssim W^{2\delta_\sto-\nu}\log W,$$
But since $\eta>0$ is fixed, we can pick $\delta_\sto$ small, s.t. the latter term is $\lesssim 1$, meaning that our result above can be stated as: 
\begin{boxtheorem}{Theorem 6.3. }There $\exists \delta>0$, s.t.:
$$|\E_t^{S, \sto}(z)_{ab}|\lesssim W^{-\delta}\cdot W^{-\frac{3}{4}}\eta_t^{-1}\cdot W^{-1}\eta_t^{-\frac{1}{2}}$$
\end{boxtheorem}

\newpage

\section{Proof of Theorems}

\subsection*{7.1 Theorem 1 - Stopping time}
\label{proof-stop}

\begin{boxtheorem}{Theorem 1: (Stopping time).} $\mathbb{P}[\tau_{\text{stop}, i} \neq  1] \lesssim_D N^{-D}$ for any $i=1, 2$ and $D > 0$.
\end{boxtheorem}
\textit{Proof} - By Boole's inequality: $$\mathbb{P}\left[\tau_{\sto, 1}\neq 1\right]\leq \mathbb{P}[\left\{\tau_{\sto, 2}\neq 1\}\cap \{\tau_{\sto, 1=1}\}\right]+\mathbb{P}[\tau_\sto=\tau_{\sto, 1}<1]$$
For any $\mathcal{A}\in \{\tau_\sto = \tau_{\sto, 1}<1\}$, we have $|\E_{\tau_\sto}(z)_{ab}|\geq W^{\delta_\sto}W^{-3/4}\eta_t^{-1}W^{-1}\eta_t^{-1/2}$, since $\E_t(z)$ is almost surely continuous in $t$. By definition of the stopped SDE, we have $\E_{\tau_\sto}(z) = \E_{\tau_\sto}^{\sto}(z)$, i.e: $$\mathbb{P}\left[\tau_\sto = \tau_{\sto, 1}<1\right]\leq \mathbb{P}\left[\max_{t\in [0, 1]}\max_{a, b}\frac{|\E_t^{M, \sto}(z)_{ab}|}{W^{-\frac{3}{4}}\eta_t^{-\frac{3}{4}}W^{-1}\eta_t^{-\frac{1}{2}}}\geq W^{\delta_\sto}\right]$$
By applying our theorems $6.1$ through $6.3$ to the RHS, we know from the definition of stochastic domination that for any $D>0$: 
$$\mathbb{P}\left[ \tau_\sto = \tau_{\sto, 1}<1\right]\leq C_DN^{-D}$$
This means that we now need to establish an equivalent bound for the first term on the RHS. Assume $\exists \mathcal{B}$ an event for which the following is true:
$$\max_{t\in [0, 1]}\max_{a, b}\frac{|\E_t(z)_{ab}|}{W^{\delta_\sto}W^{-\frac{3}{4}}\eta_t^{-\frac{3}{4}}W^{-1}\eta_t^{-\frac{1}{2}}}\lesssim 1$$
Since $T_t(z) = \Theta_t+\E_t(z)$ and $|(\Theta_t)_{ab}|\lesssim W^{-1}\eta_t^{-1/2}$ (Lemma A.2), this implies that on $\mathcal{B}$, we have $T_t(z)_{ab}\lesssim W^{-1}\eta_t^{-\frac{1}{2}}$. By the behavior of $S^{1/2}$ as an averaging operator, we have for all $a, b$ and $t\in [0, 1]$: $$(S^{1/2}T_t(z)S^{1/2})_{ab}\lesssim W^{-1}\eta_t^{-1/2}.$$ 
This means that by \hyperref[corollary-b5]{Corollary B.5} on the event $\mathcal{B}$, we get precisely: $$\frac{|G_t(z)_{ab}-m(z)\delta_{ab}|}{(S^{1/2}T_t(z)S^{1/2})_{ab}+S^{1/2}_{ab}+W^{-C}}\prec 1$$
Hence, $\mathbb{P}[\left\{\tau_{\sto, 2}\neq 1\}\cap \{\tau_{\sto, 1=1}\}\right]\rightarrow 0$ as $N\rightarrow \infty$, completing the theorem. 

\subsection*{7.2 Theorem 2 - Quantum Diffusion}
\begin{boxtheorem}{Theorem 2: (Quantum Diffusion)}
For $|E|<2$ fixed, assume $\exists \nu>0$, s.t. $\eta \asymp W^2 N^{-2}$ and $W\geq W^{8/11+\nu}$. Then: $$\max_{x, y}|T_{xy}-\Theta_{xy}|\prec W^{-\frac{7}{4}}\eta^{-\frac{3}{2}}$$
\end{boxtheorem}
\textit{Proof} - By the previous theorem and the definition of stopping time, we know: $$\max_{a,b}|\E_1(z)_{ab}|\prec W^{-\frac{7}{4}} \eta_1^{-\frac{3}{2}}.$$
But recall our flow construction that $\eta_1=\eta$, s.t $T_1=T$ and $\Theta_1=\Theta$. This means that: $$\max_{a,b}|\E_1(z)_{ab}| =\max_{x, y}|T_{xy}-\Theta_{xy}|\prec W^{-\frac{7}{4}}\eta^{-\frac{3}{2}}$$

\subsection*{7.3 Theorem 3}
\label{proof-deloc}
\begin{boxtheorem}{Theorem 3}
\noindent Assume $|E|<2$ is fixed and that $\exists \nu>0$, s.t. $\eta \asymp W^2 N^{-2}$ and $W\geq N^{8/11+\nu}$. Then:
\begin{center}$\max_{x, y}|G_{xy} - \delta_{xy}m(z)|^2\prec W^{-1}\eta^{-\frac{1}{2}}$\end{center}
\end{boxtheorem}
\textit{Proof} - By Theorem 1 and the definition $(3.4.2)$ of $\tau_\sto$, we have that: $$\max_{x, y}|G_{xy} - \delta_{xy}m(z)|^2\prec \max_{a, b}(S^{1/2}T_1(z)S^{1/2})_{ab}+\max_{a, b}S_{ab}^{1/2}$$
By Theorem 2 we showed above, \hyperref[lemma-a2]{Lemma A2}, and the behavior of $S^{1/2}$ as the averaging operator: $$\max_{a, b}T_1(z)_{ab}\prec W^{_1}\eta_1^{-1/2} = W^{-1}\eta^{-1/2}\Rightarrow \max_{ab}\left(S^{1/2}T_1(z)S^{1/2}\right)_{ab}\prec W^{-1}\eta^{-1/2}$$
$$\Rightarrow \max_{x, y}|G_{xy} - \delta_{xy}m(z)|^2\prec  W^{-1}\eta^{-1/2}+W^{-1}\prec  W^{-1}\eta^{-1/2}$$
\newpage
\subsection*{7.4 Theorem 4 - Delocalization}
\begin{boxtheorem}{Theorem 4: (Delocalization).}
For any $\ell\ll N$ and fixed $\epsilon, \kappa, c>0$: $$\frac{|\mathcal{A}_{\epsilon, \ell, \kappa}|}{N}\lesssim \sqrt{\epsilon}+\mathcal{O}(N^{-c})$$
\end{boxtheorem}
\textit{Proof} - The following proof is entirely due to Erd\H{o}s, et. al. \cite{21}, where it appears as Proposition 7.1. We present it here in order to better understand the nature of delocalization. The authors assume that $\Lambda(z) = \max_{x, y}|G_{xy}(z)-\delta_{xy}m(z)|\prec \Psi$, where $\Psi$ is a deterministic admissible parameter if $M^{-1/2}\leq \Psi^{(N)}(z)\leq M^{-\gamma/2}$ for all $N$. Note that $E\in [-2+\kappa, 2-\kappa]$ and $\eta\in M^{-1+\gamma}\leq \eta\leq 10$, where $M = \frac{1}{\max_{i, j}S_{ij}}$ and $\kappa>0$ is globally fixed. In the context of the work by Dubova and Yang we have presented so far, these assumptions are satisfied, given that Theorem 4 implies $|G_{xy}(z)-\delta_{xy}m(z)|^2\prec W^{-1}\eta^{-1/2}\lesssim N^{-1}\eta^{-1}$ by using the global diffusion scale $\eta\asymp W^2 N^{-2}$. Furthermore, by the other assumption $N\ll W^{11/8}$, one has $N^{-1}\eta^{-1}\asymp NW^{-1}\ll W^{-5/8}$\cite{bandSDE}. By Ward and the latter admissible parameter, we have for all $x$: \begin{equation*}\frac{\eta}{\im m}\sum_y|G_{yx}|^2 = \frac{\im G_{xx}}{\im m}\lesssim 1+O(\Psi)\tag{7.4.1}\end{equation*}
uniformly in $E\in [-2+\kappa, 2-\kappa]$. Hence, the map $y\mapsto \frac{\eta}{\im m(z)}|G_{yx}|^2$ is approximately a PDF on $\Gamma$:
$$\frac{\eta}{\im m(z)}\|P_{x, \ell}G \mathbf{e}_x\|^2 = \frac{\eta}{\im m(z)}\sum_y \mathbf{1}\left(|y-x|\geq \ell\right)|G_{yx}|^2=$$
\begin{equation*}=\frac{\eta}{\im m(z)}\sum_y|G_{yx}|^2 - \frac{\eta}{\im m(z)}\sum_y \mathbf{1}\left(|y-x|\geq \ell\right)|G_{yx}|^2=1+O(N^{-c}),\tag{7.4.2}\end{equation*}
for some $c>0$ by using Theorem 4, Lemma B.1 and $(7.4.1)$. Using the spectral decomposition $(1.1)$: 
\begin{equation*}\frac{\eta}{\im m(z)} \|P_{x, \ell} G\mathbf{e}_x\|^2 = \frac{\eta}{\im m(z)}\left\|\sum_\alpha \frac{1}{\lambda_\alpha-z}\overline{u}_\alpha(x)P_{x, \ell}\mathbf{u}_\alpha\right\|^2\leq \mathcal{L}+\mathcal{R}=\tag{7.4.3}\end{equation*}
$$= \frac{\eta}{\im m(z)}\left(1+\frac{1}{\xi}\right)\left\|\sum_{\alpha\in \mathcal{A}_{\epsilon, \ell}}\frac{1}{\lambda_\alpha-z}\overline{u}_\alpha(x)P_{x, \ell}\mathbf{u}_\alpha\right\|^2 +  \frac{\eta}{\im m(z)}\left(1+\frac{1}{\xi}\right)\left\|\sum_{\alpha\in \mathcal{A}^C_{\epsilon, \ell}}\frac{1}{\lambda_\alpha-z}\overline{u}_\alpha(x)P_{x, \ell}\mathbf{u}_\alpha\right\|^2,$$
where $\xi>0$ is arbitrary and $\mathcal{A}^C_{\epsilon, \ell} =\{1, \cdots, N\}\backslash A_{\epsilon, \ell}$ is the complement set. We can bound the left $\mathcal{L}$ term, using $(2.1)$ and Lemma A.1:
$$\mathcal{L}\leq \frac{\eta}{\im m(z)}\left\|\sum_{\alpha\in \mathcal{A}_{\epsilon, \ell}}\frac{1}{\lambda_\alpha-z}\overline{u}_\alpha(x)\mathbf{u}_\alpha\right\|^2\leq \frac{\eta}{\im m(z)}\sum_\alpha\frac{|u_\alpha(x)|^2}{|\lambda_\alpha-z|^2} = \frac{\im G_{xx}}{\im m(z)}\lesssim 1 + O(\Psi)$$
$$\Rightarrow \mathcal{L}\leq \frac{\eta}{\im m(z)}\left\|\sum_{\alpha\in \mathcal{A}_{\epsilon, \ell}}\frac{1}{\lambda_\alpha -z}\overline{u}_\alpha(x)\mathbf{u}_\alpha\right\|^2\leq (1+O(\Psi))\sqrt{\frac{\eta}{\im m(z)}}\left|\sum_{\alpha\in \mathcal{A}_{\epsilon, \ell}}\frac{1}{\lambda_\alpha-z}\overline{u}_\alpha(x)P_{x, \ell}\mathbf{u}_\alpha\right\|\leq$$
\begin{equation*}\leq (C+O(\Psi))\sum_{\alpha\in \mathcal{A}_{\epsilon, \ell}}\sqrt{\frac{\eta}{(E-\lambda_\alpha)^2+\eta^2}}|u_\alpha(x)|\|P_{x, \ell}\mathbf{u}_\alpha\|\tag{7.4.4}\end{equation*}
By the same token, we get the following bound on $\mathcal{R}$: 

\begin{equation*} \tag{7.4.5}\mathcal{R}\leq \left\|\sum_{\alpha\in \mathcal{A}_{\epsilon, \ell}^C}\frac{1}{\lambda_\alpha-z}\overline{u}_\alpha(x)\mathbf{u}_\alpha\right\|^2 = \sum_{\alpha\in \mathcal{A}_{\epsilon, \ell})}\frac{|u_\alpha(x)|^2}{|\lambda_\alpha-z|^2}\end{equation*}
Putting all the results we got, namely $(7.4.3), (7.4.3), (7.4.4)$ and $(7.4.5)$, we get uniformly on $E\in [-2+\kappa, 2-\kappa]$: 
$$1+O(N^{-c})\leq (C+O(N^{-c}))\left(1+\frac{1}{\xi}\right)\sum_{\alpha\in \mathcal{A}_{\epsilon, \ell}}\sqrt{\frac{\eta}{(E-\lambda_\alpha)^2+\eta^2}}|u_\alpha(x)|\|P_{x, \ell}\mathbf{u}_\alpha\|+$$
$$+(1+\xi)\sum_{\alpha\in \mathcal{A}_{\epsilon, \ell}^c}\frac{\eta|u_\alpha(x)|^2}{\im m(z)|\lambda_\alpha-z|^2}$$
Using the fact that$\int_{[-2+\kappa, 2-\kappa]}dR\frac{\eta}{(E-\lambda_\alpha)^2+\eta^2}\geq \frac{\pi}{2}$, we can use Jensen's inequality:
$$1+O(N^{-c})\leq (C+O(N^{-c}))\left(1+\frac{1}{\xi}\right)\sum_{\alpha\in \mathcal{A}_{\epsilon, \ell}}|u_\alpha(x)|\|P_{x, \ell}\mathbf{u}_\alpha\|+$$
\begin{equation*}+(1+\xi)\frac{1}{4-2\kappa}\int_{[-2+\kappa, 2-\kappa]}dE\sum_{\alpha\in \mathcal{A}_{\epsilon, \ell}^c}\frac{\eta|u_\alpha(x)|^2}{\im m(z)|\lambda_\alpha-z|^2}\tag{7.4.6}\end{equation*}
By the trace inequality: 
$$\frac{1}{4-2\kappa}\int_{[-2+\kappa, 2-\kappa]}\frac{dE}{N}\sum_{\alpha\in \mathcal{A}_{\epsilon, \ell}^c}\frac{\eta|u_\alpha(x)|^2}{\im m(z)|\lambda_\alpha-z|^2}\leq \frac{1}{4-2\kappa} \int dE \frac{1}{N}\frac{\im \text{Tr }G}{\im m(z)} = 1+O(\Psi)$$
Using this, we can average out $(7.4.6)$ with $N^{-1}\sum_x$:
$$\Rightarrow \frac{1}{4-2\kappa}\int dE \frac{1}{N}\sum_{\alpha\in \mathcal{A}_{\epsilon, \ell}^C}\frac{\eta}{\im m(z)|\lambda_\alpha-E-i\eta|^2}\geq 1-(C+O(N^{-c}))\sqrt{\epsilon})\sqrt{\epsilon}-O(N^{-c})\geq 1-C\sqrt{\epsilon}-O(N^{-c})$$
\begin{equation*}\Rightarrow \frac{1}{4-2\kappa}\int dE \frac{1}{N}\sum_{\alpha\in \mathcal{A}_{\epsilon, \ell}^C}\frac{\eta}{\im m(z)|\lambda_\alpha-E-i\eta|^2}\leq C\sqrt{\epsilon}+O(N^{-c})\tag{$\square$},\end{equation*}
where we set $\xi =\sqrt{\epsilon}$. 



\newpage

\section{Non-Gaussian case}
\label{sec:non-gaussian}
Up to this point, the work presented on band matrices has been focused exclusively on the case with Gaussian entries. Given that the delocalization behavior is conjectured to extend beyond this restriction, the goal for the last section of this thesis is to do exactly that, proving a \textit{Five Moment Theorem} for RBMs with the help of the Lindeberg exchange strategy:


\subsection{Lindeberg exchange strategy}
This method was developed by Lindeberg in his seminal proof for the generalized CLT \cite{lindeberg} and later extended by Chatterjee \cite{chatterjee}. The argument proceeds as follows \cite{tao}: Suppose that $X_1, \cdots, X_n\sim [\mu, 1]$ are i.i.d, and let $Y_1, \cdots, Y_n\sim[\mu, 1]$ be another such set. One would like to show that for any smooth, compactly supported function $F$: $$\mathbb{E}F\left(\frac{X_1+\cdots+X_n}{\sqrt{n}}\right) = \mathbb{E}F\left(\frac{Y_1+\cdots+Y_n}{\sqrt{n}}\right)+o(1),$$
by swapping the entries one at a time, with the latter being the key step in this argument.  Define
$S= \frac{X_1+\cdots+X_{n-1}}{\sqrt{n}}$, where clearly $ \frac{X_1+\cdots+X_n}{\sqrt{n}}=S+n^{-1/2}X_n $. Using the smoothness and compact support of $F$, we can apply the Taylor expansion around $S$: $$F\left(\frac{X_1+\cdots+X_n}{\sqrt{n}}\right)=F(S)+n^{-1/2}X_nF'(S)+\frac{1}{2}n^{-1}X_n^2 F''(S) + O(n^{-3/2}|X_n|^3)$$
By taking the expectation, we get: 
$$ \mathbb{E}F\left(\frac{X_1+\cdots+X_n}{\sqrt{n}}\right)=\mathbb{E}F(S)+n^{-1/2}(\mathbb{E}X_n)\mathbb{E}F'(S)+\frac{1}{2}n^{-1}(\mathbb{E}X_n^2)\mathbb{E}F''(S) + O(n^{-3/2})$$
The following step is to swap the last element $X_n$ with an independent $Y_n$ from the other set $S' = \frac{X_1+\cdots X_{n-1}+Y_n}{\sqrt{n}}$. By the same logic as above, we have that $\mathbb{E}F(S')$ is: 
$$ \mathbb{E}F\left(\frac{X_1+\cdots+X_{n-1}+Y_n}{\sqrt{n}}\right)=\mathbb{E}F(S)+n^{-1/2}(\mathbb{E}Y_n)\mathbb{E}F'(S)+\frac{1}{2}n^{-1}(\mathbb{E}Y_n^2)\mathbb{E}F''(S) +O (n^{-3/2})$$
By taking their difference and using the matching moments of second order, we get: 
$$ \mathbb{E}F\left(\frac{X_1+\cdots+X_n}{\sqrt{n}}\right) = \mathbb{E}F\left(\frac{X_1+\cdots+X_{n-1}+Y_n}{\sqrt{n}}\right)+o\left(\frac{1}{n}\right)$$
If we have higher matching moments, we could simply continue the Taylor expansion in order to refine the error term.
This strategy has been extended to the case of independent Wigner matrices $M_n, M_n'$ with the goal of obtaining similar bounds like: $$\mathbb{E}F(M_n)-\mathbb{E}F(\tilde{M}_n)=o(1/n).$$ 
\noindent Given that up to this point, we have had the ample opportunity of exploring the powerful uses of Green's function, one may ask if there is an equivalent approach for the spectral behavior of random matrices. There is, indeed, such a result, also known as the \textit{Four Moment Theorem for Green's function} \cite{tao2}:


\subsection*{Four Moment Theorem for Green's function}
Let $M_n, M_n'$ be two Wigner random matrices satisfying the following condition: $$P(|(M_n)_{ij}|\geq t^C)\leq \exp(-t)$$
for all $1\leq i, j\leq n$ and $t\geq C'$ for some constants $C, C'$, independent of $i, j$ and $n$. Furthermore, assume that $M_n$ and $M_n'$ match in moments up to order 4 off the diagonal and up to order 2 on the diagonal for some sufficiently large $C_0$. Let $z = E + i\eta$ for some $E \in \mathbb{R}$ and some $\eta > 0$. Additionally, assume the level repulsion hypothesis, i.e that for any $c > 0$, one has with high probability that:
\begin{equation*}
    \inf_{1 \leq i \leq n} |\lambda_i(\sqrt{n}M_n) - nz| \geq n^{-c}
\end{equation*}
\begin{equation*}
    \inf_{1 \leq i \leq n} |\lambda_i(\sqrt{n}M_n') - nz| \geq n^{-c}.
\end{equation*}

\noindent Let $1 \leq p,q \leq n$. Then for any smooth function $G: \mathbb{C} \to \mathbb{C}$ obeying the bounds $\nabla^j G(x) = O(1)$ for all $x \in \mathbb{C}$ and $0 \leq j \leq 5$, the theorem states that for some constant $c_0 > 0$ independent of $n$:
\begin{equation*}
    \mathbb{E} G \left( \left( \frac{1}{\sqrt{n}} M_n - zI \right)^{-1}_{pq} \right)
    - \mathbb{E} G \left( \left( \frac{1}{\sqrt{n}} M_n' - zI \right)^{-1}_{pq} \right)
    = O(n^{-c_0})
\end{equation*}



\subsection{Replacement}
The goal for the remainder of this thesis is to prove a theorem, similar to the one above, for the delocalization of RBMs with a relaxation on the Gaussian assumption. Recall that at the root of the Lindeberg exchange strategy is the stochastic control at each step of replacement. While we are working with matrices instead of vectors of random variables, the notion of replacement naturally extends when we consider a lexicographic ordering of the indices. Similar to how we swapped $X_n$ with $Y_n$, in this context we will also start with $H_{NN}$. We will work our way backwards by replacing consecutively $H_{N, N-1}, \cdots,H_{i, j+1},  H_{ij}, H_{i, j-1}\cdots$.  As such, at step $(i, j)$ of the replacement (accounting for the Hermitian constraint), we will have the following two matrices, represented in lexicographic vector form:
$$H^{(i, j)}=(H_{11}, \cdots H_{i, j-1}, H_{ij}, H_{i, j+1}^G, \cdots, H_{NN}^G)$$
$$H^{(i, j), G}=(H_{11}, \cdots H_{i, j-1}, H_{ij}^G, H_{i, j+1}^G, \cdots, H_{NN}^G),$$
where $H_{N, N}^G, \cdots, H_{i, j+1}^G,  H_{ij}^G$ are the Gaussian replacement entries, matching in the sense of $(8.1.1)$ :

\begin{boxt}{Conditions}
\begin{enumerate}[label=(8.1.\arabic*),align=left, leftmargin=30pt, labelwidth=25pt, labelsep=5pt]\label{8.1.1}
\item $\mathbb{E}\left[H_{ij}^\ell \left(\overline{H}_{ij}\right)^{k-\ell}\right] = \mathbb{E}\left[\left(H_{ij}^G\right)^\ell \left(\overline{H}_{ij}^G\right)^{k-\ell}\right], \quad 0\leq \ell\leq k\leq 5$\label{8.1.2}
\item $\forall q\geq 1, \quad \mathbb{E}\left\vert H_{xy}\right\vert ^{2q}\prec W^{-q}$
\item $\max_x|G(z)_{xx}|\prec 1$ and $\max_{a, b}|G(z)_{ab}-m(z)\delta_{ab}|\prec W^{-\delta}$ for some fixed $\delta>0$. \label{8.1.3}
\end{enumerate}
\end{boxt}
The first condition $(8.1.1)$ is the $\textit{Five Moment}$ matching assumption that gives the name of our result. The second $(8.1.2)$ is a moment decay condition that is standard in the context of RMT ($5.6$ in \cite{dynamic}) and will allow us to relax the Gaussian assumption, while still providing sufficient strength of our bounds. The last condition assumes weak a priori on- and off-diagonal estimates, in line with what the conditions were for Lemma B.3 and B.4. We will use them to verify that each consecutive replacement is manageable by constructing a continuity argument. As a first step, let us define the interpolation for $t\in [0, 1]$ at entries $(ij)$ and $(ji)$, where with slight abuse of notation:
$$tH_{ij} = H^{(i, j)}[H_{ij}, H_{ji}\mapsto t H_{ij}, tH_{ji}]$$
$$tH_{ij}^{G} = H^{(i, j),G}[H^G_{ij}, H^G_{ij}\mapsto t H^G_{ij}, t H^{G}_{ij}].$$
Whenever the replacement step $(i, j)$ is implicit from context, we will drop the superscript. In order to make our theorem applicable to the context of delocalization, we want to consider the following function: 
\begin{equation*}\label{8.3}
F(H) = \left\vert
(H-z)_{xy}^{-1} - m(z)\delta_{xy}
\right\vert ^{2p} = \left\vert 
G_{xy}(z)-m(z)\delta_{xy}
\right\vert ^{2p} = |W_{xy}|^{2p}\tag{8.3}
\end{equation*}
Our reasoning is simple. Having familiarized ourselves with the flow method and the resolvent approach, it makes sense for us to choose a function that quantifies the deviation of resolvent entries form their deterministic approximations. By raising to power $2p$, we have better tools at our disposal, namely Chebyshev' and Markov's inequalities. Let us apply our strategy to $F$ at step $(i,j)$ for $t\in [0, 1]$: 

$$F(tH_{ij}) = F(H_{11}, \cdots, H_{i, j-1}, tH_{ij}, H_{i, j+1}^G, \cdots H_{NN}^G),$$
up to the Hermitian constraint. In order to apply the Taylor series for the error terms, let us calculate the derivatives: $$\partial_t F(tH_{ij}) = \partial_{H_{ij}} F(tH_{ij})H_{ij} + \partial_{H_{ji}} F(tH_{ij})\overline{H_{ij}}$$
We claim that the following represents the $k$th derivative: 
\begin{equation*}\frac{d^k}{dt^k} F(tH_{ij}) = \sum_{\ell=0}^k \binom{k}{\ell}\partial_{H_{ij}}^\ell \partial_{H_{ji}}^{k-\ell} F(tH_{ij})H_{ij}^\ell \left(\overline{H_{ij}}\right)^{k-\ell}\tag{8.4}\end{equation*}
Assume the latter for an induction hypothesis, having already verified the base case. Then: 
$$\frac{d^{k+1}}{dt^{k+1}} F(tH_{ij}) = \frac{d}{dt}\left[\frac{d^k}{dt^k} F(tH_{ij}) \right]=\frac{d}{dt}\left[ \sum_{\ell=0}^k \binom{k}{\ell}\partial_{H_{ij}}^\ell \partial_{H_{ji}}^{k-\ell} F(tH_{ij})H_{ij}^\ell \left(\overline{H_{ij}}\right)^{k-\ell}\right]=$$
$$=\sum_{\ell=0}^k\binom{k}{\ell}\frac{d}{dt}\left[\partial_{H_{ij}}^\ell \partial_{H_{ji}}^{k-\ell} F(tH_{ij})\right]H_{ij}^\ell \left(\overline{H_{ij}}\right)^{k-\ell}$$

\noindent By the chain rule for any $\ell$ and $m$:
$$\frac{d}{dt}\left[\partial_{H_{ij}}^\ell \partial_{H_{ji}}^m F(tH_{ij})\right] = \partial_{H_{ij}}^{\ell+1} \partial_{H_{ji}}^m F(tH_{ij})H_{ij} + \partial_{H_{ij}}^\ell \partial_{H_{ji}}^{m+1} F(tH_{ij})\overline{H_{ij}}$$
Substituting in the derivative formula gives us:
$$
= \sum_{\ell=0}^k \binom{k}{\ell}\left[\partial_{H_{ij}}^{\ell+1} \partial_{H_{ji}}^{k-\ell} F(tH_{ij})H_{ij} + \partial_{H_{ij}}^\ell \partial_{H_{ji}}^{k-\ell+1} F(tH_{ij})\overline{H_{ij}}\right]H_{ij}^\ell \left(\overline{H_{ij}}\right)^{k-\ell}$$
$$
= \sum_{\ell=0}^k \binom{k}{\ell}\partial_{H_{ij}}^{\ell+1} \partial_{H_{ji}}^{k-\ell} F(tH_{ij})H_{ij}^{\ell+1} \left(\overline{H_{ij}}\right)^{k-\ell} + \sum_{\ell=0}^k \binom{k}{\ell}\partial_{H_{ij}}^\ell \partial_{H_{ji}}^{k-\ell+1} F(tH_{ij})H_{ij}^\ell \left(\overline{H_{ij}}\right)^{k-\ell+1}
$$
$$
= \sum_{\ell=1}^{k+1} \binom{k}{\ell-1}\partial_{H_{ij}}^{\ell} \partial_{H_{ji}}^{k+1-\ell} F(tH_{ij})H_{ij}^{k+1-\ell} \left(\overline{H_{ij}}\right)^{k+1-\ell} + \sum_{\ell=0}^k \binom{k}{\ell}\partial_{H_{ij}}^\ell \partial_{H_{ji}}^{k-\ell+1} F(tH_{ij})H_{ij}^\ell \left(\overline{H_{ij}}\right)^{k-\ell+1}
$$
But observe that we can split the sums and group it as follows:
$$=\binom{k}{0}\partial_{H_{ij}}^0\partial_{H_{ji}}^{k+1} F(tH_{ij})\left(\overline{H_{ij}}\right)^{k+1}+\binom{k}{k}\partial_{H_{ij}}^{k+1}\partial_{H_{ji}}^{0} F(tH_{ij})\left(\overline{H_{ij}}\right)^{k+1}+$$
$$=+\sum_{\ell=1}^k \left(\binom{k}{\ell-1}+\binom{k}{\ell}\right)\partial_{H_{ij}}^{\ell} \partial_{H_{ji}}^{k+1-\ell} F(tH_{ij})H_{ij}^{\ell} \left(\overline{H_{ij}}\right)^{k+1-\ell} $$
But since $\binom{k}{k} = \binom{k+1}{k+1}$ and $\binom{k}{\ell-1}+\binom{k}{\ell} = \binom{k+1}{\ell}$, this completes the induction:
\begin{equation*}
\frac{d^{k+1}}{dt^{k+1}} F(tH_{ij}) = \sum_{\ell=0}^{k+1} \binom{k+1}{\ell}\partial_{H_{ij}}^\ell \partial_{H_{ji}}^{k+1-\ell} F(tH_{ij})H_{ij}^\ell \left(\overline{H_{ij}}\right)^{k+1-\ell}
\end{equation*}

\subsection*{Taylor series}
Having derived the explicit derivatives, we can now continue with the next step in the Lindeberg exchange strategy. By the Lagrange form of Taylor's (quintic) theorem we know $\exists t_{ij}\in [0, 1]$, s.t.:
$$F(H_{ij}) =F(0) + \sum_{k=1}^5 \frac{1}{k!}\left.  \frac{d^k}{dt^k}F(tH_{ij})\right\vert_{t=0}+ \left.\frac{1}{720} \frac{d^6}{dt^6}F(tH_{ij})\right\vert_{t=t_{ij}},$$
where $F(sH_{ij})=F(0)$ for $s=0$ represents setting $H_{ij}$ and $H_{ji}$ to zero. Similarly, $\exists t_{ij}^G\in[0, 1]$, s.t.:
$$F(H_{ij}^G) =F(0) + \sum_{k=1}^5 \frac{1}{k!} \left. \frac{d^k}{dt^k}F(tH_{ij}^G)\right\vert_{t=0}+ \frac{1}{720}\left. \frac{d^6}{dt^6}F(tH_{ij}^G)\right\vert_{t=t_{ij}^G}.$$
Let us take their difference:
\begin{equation*}\begin{array}{l}
F(H_{ij}) - F(H_{ij}^G)=
\end{array}
\end{equation*}
\begin{equation*}=\sum_{k=1}^5 \frac{1}{k!} \left[\left. \frac{d^k}{dt^k}F(tH_{ij})\right\vert_{t=0}-\left. \frac{d^k}{dt^k}F(tH_{ij}^G)\right\vert_{t=0}\right\vert +\frac{1}{720} \left[\left.\frac{d^6}{dt^6}F(tH_{ij})\right\vert_{t=t_{ij}}-\left.\frac{d^6}{dt^6}F(tH_{ij}^G)\right\vert_{t=t_{ij}^G}\right] \tag{8.5}
\end{equation*}
Focusing on the first term, we can use formula $(8.4)$ we derived earlier, s.t. the first term becomes:
$$\sum_{k=1}^5\frac{1}{k!} \left[\sum_{\ell=0}^k \left.\binom{k}{\ell}\partial_{H_{ij}}^\ell \partial_{H_{ji}}^{k-\ell} F(tH_{ij})H_{ij}^\ell \left(\overline{H_{ij}}\right)^{k-\ell}\right\vert_{t=0}-\sum_{\ell=0}^k \left.\binom{k}{\ell}\partial_{H_{ij}}^\ell \partial_{H_{ji}}^{k-\ell} F(tH_{ij}^G)(H_{ij}^G)^\ell \left(\overline{H_{ij}}\right)^{k-\ell}\right\vert_{t=0}\right]=$$

$$=\sum_{k=1}^t\frac{1}{k!}\sum_{\ell=0}^k \binom{k}{\ell}\partial_{H_{ij}}^\ell \partial_{H_{ji}}^{k-\ell} F(0) \left\{H_{ij}^\ell \left(\overline{H_{ij}}\right)^{k-\ell}-(H_{ij}^G)^\ell \left(\overline{H_{ij}}\right)^{k-\ell}\right\}$$
Let us take the expectation of $(8.4)$. We can use the independence of $F(0)$ from $H_{ij}$ and $H_{ij}^G$ to split the expectation of the first term as:

$$\mathbb{E}\left[\sum_{k=1}^t\frac{1}{k!}\sum_{\ell=0}^k \binom{k}{\ell}\partial_{H_{ij}}^\ell \partial_{H_{ji}}^{k-\ell} F(0) \left\{H_{ij}^\ell \left(\overline{H_{ij}}\right)^{k-\ell}-(H_{ij}^G)^\ell \left(\overline{H_{ij}}\right)^{k-\ell}\right\}\right]=$$
$$=\sum_{k=1}^t\frac{1}{k!}\sum_{\ell=0}^k \binom{k}{\ell}\mathbb{E}\left[\partial_{H_{ij}}^\ell \partial_{H_{ji}}^{k-\ell} F(0) \right]\mathbb{E}\left[H_{ij}^\ell \left(\overline{H_{ij}}\right)^{k-\ell}-(H_{ij}^G)^\ell \left(\overline{H_{ij}}\right)^{k-\ell}\right]=$$
$$=\sum_{k=1}^t\frac{1}{k!}\sum_{\ell=0}^k \binom{k}{\ell}\mathbb{E}\left[\partial_{H_{ij}}^\ell \partial_{H_{ji}}^{k-\ell} F(0) \right]\left\{\mathbb{E}\left[H_{ij}^\ell \left(\overline{H_{ij}}\right)^{k-\ell}\right]-\mathbb{E}\left[(H_{ij}^G)^\ell \left(\overline{H_{ij}}\right)^{k-\ell}\right]\right\}=0$$
by condition $(8.1.1)$. This means that:
$$\mathbb{E}\left[F(tH_{ij})-F(tH_{ij}^G)\right]=\frac{1}{720} \mathbb{E}\left[\left.\frac{d^6}{dt^6}F(tH_{ij})\right\vert_{t=t_{ij}}\right]-\mathbb{E}\left[\left.\frac{d^6}{dt^6}F(tH_{ij}^G)\right\vert_{t=t_{ij}^G}\right]$$
By the Hölder' and the triangle inequality: 
$$\left\vert \mathbb{E}\left[ F(H_{ij}) - F(H_{ij}^G)\right] \right\vert\leq \frac{1}{720} \left\vert\mathbb{E}\left[\left.\frac{d^6}{dt^6}F(tH_{ij})\right\vert_{t=t_{ij}}\right]\right\vert+\frac{1}{720}\left\vert\mathbb{E}\left[\left.\frac{d^6}{dt^6}F(tH_{ij}^G)\right\vert_{t=t_{ij}^G}\right]\right\vert $$
$$\lesssim \mathbb{E}\left\vert\left.\frac{d^6}{dt^6}F(tH_{ij})\right\vert_{t=t_{ij}}\right\vert+\mathbb{E}\left\vert\left.\frac{d^6}{dt^6}F(tH_{ij}^G)\right\vert_{t=t_{ij}^G}\right\vert = \mathcal{P} + \mathcal{P}^G$$
Applying formula $(8.4)$ to $\mathcal{P}$ (with $\mathcal{P}^G$ following WLOG), we get:
$$\mathcal{P} = \mathbb{E}\left\vert\left.\frac{d^6}{dt^6}F(tH_{ij})\right\vert_{t=t_{ij}}\right\vert  = \mathbb{E}\left\vert  \left.\sum_{\ell=0}^6 \binom{6}{\ell}\partial_{H_{ij}}^\ell \partial_{H_{ji}}^{6-\ell} F(tH_{ij})H_{ij}^\ell \left(\overline{H_{ij}}\right)^{6-\ell}\right\vert_{t=t_{ij}}\right\vert\leq $$
$$\leq \sum_{\ell=0}^6 \binom{6}{\ell}\mathbb{E}\left\vert \left.
\partial_{H_{ij}}^\ell \partial_{H_{ji}}^{6-\ell} F(tH_{ij})H_{ij}^\ell \left(\overline{H_{ij}}\right)^{6-\ell}\right\vert_{t=t_{ij}}
\right\vert=$$
$$ = \sum_{\ell=0}^6 \binom{6}{\ell}\mathbb{E}\left[\left\vert \left.
\partial_{H_{ij}}^\ell \partial_{H_{ji}}^{6-\ell} F(tH_{ij})\right\vert _{t=t_{ij}}\right\vert
\left\vert H_{ij}^\ell \right\vert \left\vert \left(\overline{H_{ij}}\right)^{6-\ell}
\right\vert \right]\leq $$
\begin{equation*}\leq \sum_{\ell=0}^6\binom{6}{\ell} \mathbb{E}\left[ \left\vert \partial_{H_{ij}}^\ell \partial_{H_{ji}}^{6-\ell} F(tH_{ij})|_{t=t_{ij}}\right\vert\right] \mathbb{E}\left[|H_{ij}|^\ell \left\vert \overline{H_{ij}}\right\vert^{6-\ell}\right]\lesssim \sum_{\ell=0}^6 \mathbb{E}\left\vert \partial_{H_{ij}}^\ell \partial_{H_{ji}}^{6-\ell} F(tH_{ij})|_{t=t_{ij}}\right\vert\mathbb{E}|H_{ij}|^6,\tag{8.5.1}\end{equation*}
where the first is the triangle inequality, the second is Hölder's $L^1$ inequality and the last step uses the fact that $|H_{ij}| = \left\vert\overline{H_{ij}}\right\vert$. Let us call $\mathcal{Q}=\left\vert \partial_{H_{ij}}^\ell \partial_{H_{ji}}^{6-\ell} F(tH_{ij})|_{t=t_{ij}}\right\vert$. Observe that WLOG, we have the same result for $\mathcal{P}^G$ and $\mathcal{Q}^G$, respectively. By condition $(8.1.2)$, we have that $\mathbb{E}|H_{ij}|^6\prec W^{-3}$, so our goal will be to establish a bound on $\mathcal{Q}$ for all $t\in [0, 1]$. The delocalization argument will follow readily, so let us now direct our attention to the derivatives on $F(H)$ w.r.t. the entries $H_{ij}$. 

\subsection{Derivatives of $F(tH_{ij})$}
By looking at the general form of our target $\left\vert \partial_{H_{ij}}^\ell \partial_{H_{ji}}^{k-\ell} F(tH_{ij})\right\vert$ , we can glean that a good place to start would be with the calculation of $\partial_{H_{ij}}F(tH_{ij})$. Given our definition of the interpolation $tH_{ij}$, we can tell in advance that term $t$ will only appear as a multiple after we apply the chain rule, so we can focus exclusively on $\partial_{H_{ij}}^\ell F(tH_{xy})$. Recall our definition $(8.3)$ of $F(H)$: 
$$F(H) = \left\vert (H-z)_{xy}^{-1} - m(z)\delta_{xy} \right\vert ^{2p} = \left\vert G_{xy}(z)-m(z)\delta_{xy} \right\vert ^{2p} = |W_{xy}|^{2p}
$$

\noindent Let us apply the chain rule, by recognizing that $|W_{xy}|^{2p} = \left(W_{xy} \cdot \overline{W_{xy}}\right)^p$:

$$\frac{\partial |W_{xy}|^{2p}}{\partial H_{ij}} = \frac{\partial |W_{xy}|^{2p}}{\partial W_{xy}} \cdot \frac{\partial W_{xy}}{\partial G_{xy}} \cdot \frac{\partial G_{xy}}{\partial H_{ij}} + \frac{\partial |W_{xy}|^{2p}}{\partial \overline{W_{xy}}} \cdot \frac{\partial \overline{W_{xy}}}{\partial \overline{G_{xy}}} \cdot \frac{\partial \overline{G_{xy}}}{\partial H_{ij}}$$

\noindent Using complex differentiation, we have: $$\frac{\partial |W_{xy}|^{2p}}{\partial W_{xy}} = p|W_{xy}|^{2p-2} \cdot \overline{W_{xy}}, \quad\frac{\partial |W_{xy}|^{2p}}{\partial \overline{W_{xy}}} = p|W_{xy}|^{2p-2} \cdot W_{xy}$$

\noindent And since $W_{xy} = G_{xy} - m\delta_{xy}$: $$\frac{\partial W_{xy}}{\partial G_{xy}} = 1, \quad \frac{\partial \overline{W_{xy}}}{\partial \overline{G_{xy}}} = 1$$

\noindent Recall the identities we proved in Sec.2: $\partial_{H_{ij}}G_{xy} = -G_{xi}G_{jy}$ and $\partial_{\overline{H_{ij}}}\overline{G_{xy}} = -\overline{G_{xi}}\ \overline{G_{jy}}$. Substituting the derivatives into our chain rule expression yields:

$$\frac{\partial |W_{xy}|^{2p}}{\partial H_{ij}} = p|W_{xy}|^{2p-2} \cdot \overline{W_{xy}} \cdot 1 \cdot (-G_{xi}G_{jy}) + p|W_{xy}|^{2p-2} \cdot W_{xy} \cdot 1 \cdot (-\overline{G_{xi}}\overline{G_{jy}})$$

$$\frac{\partial |W_{xy}|^{2p}}{\partial H_{ij}} = -p|W_{xy}|^{2p-2} \cdot [\overline{W_{xy}} \cdot G_{xi}G_{jy} + W_{xy} \cdot \overline{G_{xi}G_{jy}}]$$

$$\frac{\partial |W_{xy}|^{2p}}{\partial H_{ij}} = -p|W_{xy}|^{2p-2} \cdot [\overline{W_{xy}} \cdot G_{xi}G_{jy} + (\overline{W_{xy}} \cdot G_{xi}G_{jy})^*]$$

$$\frac{\partial |W_{xy}|^{2p}}{\partial H_{ij}} = -2p|W_{xy}|^{2p-2} \cdot \text{Re}(\overline{W_{xy}} \cdot G_{xi}G_{jy})$$

\noindent Hence, by putting everything together, we get: $$\partial_{H_{ij}}F(H) = \partial_{H_{ij}} |W_{xy}|^{2p} = 2p|W_{xy}|^{2p-2}\text{Re}\left(W_{xy}\partial_{H_{ij}}\overline{W_{xy}}\right)= -2p|W_{xy}|^{2p-2}\text{Re}\left(W_{xy}\overline{G_{xi}G_{jy}}\right)$$

\noindent Let us calculate the second derivative:

$$\partial_{H_{ij}}^2F(H) = \partial_{H_{ij}}\left[ 2p|W_{xy}|^{2p-2}\text{Re}\left(W_{xy}\partial_{H_{ij}}\overline{W_{xy}}\right)\right]=$$ $$= -2p \partial_{H_{ij}}\left[|W_{xy}|^{2p-2}\right]\text{Re}\left(W_{xy}\overline{G_{xi}G_{jy}}\right)-2p|W_{xy}|^{2p-2}\partial_{H_{ij}} \left[\text{Re}\left(W_{xy}\overline{G_{xi}G_{jy}}\right)\right]$$

\noindent The first term $\partial_{H_{ij}}\left[|W_{xy}|^{2p-2}\right] $ is: $$ (2p-2)|W_{xy}|^{2p-4}\partial_{H_{ij}}\left[|W_{xy}|^2\right] = (2p-2)|W_{xy}|^{2p-4}\left(\partial_{H_{ij}} W_{xy}\cdot \overline{W}_{xy}+W_{xy}\cdot \partial_{H_{ij}} \overline{W}_{xy}\right)=$$ $$= -(2p-2)|W_{xy}|^{2p-4}\left(G_{xi}G_{jy}\cdot \overline{W}_{xy}+W_{xy}\cdot \overline{G_{xi}G_{jy}}\right) = -(4p-4)|W_{xy}|^{2p-4} \text{Re}\left(W_{xy}\cdot \overline{G_{xi}G_{jy}}\right)$$

\noindent And the second is: $$\partial_{H_{ij}}\left[\text{Re}\left(W_{xy}\overline{G_{xi}G_{jy}}\right)\right] = \text{Re}\left(\partial_{H_{ij}} W_{xy} \overline{G_{xi}G_{jy}}+W_{xy} \partial_{H_{ij}} \overline{G_{xi}G_{jy}}\right)=$$ $$=\text{Re}\left(-G_{xi}G_{jy}\overline{G_{xi}G_{jy}}+W_{xy}\partial_{H_{ij}}\left[\overline{G_{xi}G_{jy}}\right]\right) =-\text{Re}\left(G_{xi}G_{jy}\overline{G_{xi}G_{jy}}+W_{xy}\overline{G_{xi}G_{ii}}\overline{G_{jy}}+\overline{G_{xi}}\overline{G_{jj}G_{jy}}\right) $$

\noindent As such, by combining the terms, we get the following expansion: $$(8p^2-8p)|W_{xy}|^{2p-4}\text{Re}^2\left(W_{xy}\overline{G_{xi}G_{jy}}\right) +2p|W_{xy}|^{2p-2}\left[|G_{xi}G_{jy}|^2+\text{Re}\left(W_{xy}\overline{G_{xi}G_{ii}}\overline{G_{jy}}\right)+\text{Re}\left(\overline{G_{xi}}\overline{G_{jj}G_{jy}}\right)\right]$$

\noindent While this doesn't look as bad, we have already went from one to four terms in our expression (two and eight, respectively, if written in conjugate form). In fact, the growth is factorial, giving that the chain rule generates several new elements for each term: 
$$\partial_{H_{ij}}^3F(H) = -48p(p-1)(p-2)|W_{xy}|^{2p-6}\text{Re}^3\left(W_{xy}\overline{G_{xi}G_{jy}}\right)$$ $$+24p(p-1)|W_{xy}|^{2p-4}\text{Re}\left(W_{xy}\overline{G_{xi}G_{jy}}\right)\text{Re}\left(W_{xy}\overline{G_{xx}G_{ii}G_{jy}}+W_{xy}\overline{G_{xi}G_{jj}G_{yy}}\right)$$ $$-4p(2p-2)|W_{xy}|^{2p-4}\text{Re}\left(W_{xy}\overline{G_{xi}G_{jy}}\right)|G_{xi}G_{jy}|^2-2p|W_{xy}|^{2p-2}\text{Re}\left(W_{xy}\overline{G_{xi}G_{jx}G_{ii}G_{jy}}\right)$$ $$-2p|W_{xy}|^{2p-2}\text{Re}\left(W_{xy}\overline{G_{xx}G_{ii}G_{xx}G_{ii}G_{jy}}\right)-2p|W_{xy}|^{2p-2}\text{Re}\left(W_{xy}\overline{G_{xx}G_{ii}G_{xi}G_{jj}G_{yy}}\right)$$ $$-2p|W_{xy}|^{2p-2}\text{Re}\left(W_{xy}\overline{G_{xi}G_{ji}G_{jj}G_{yy}}\right)-2p|W_{xy}|^{2p-2}\text{Re}\left(W_{xy}\overline{G_{xi}G_{jj}G_{yi}G_{jy}}\right)$$ 
$$-2p|W_{xy}|^{2p-2}\text{Re}\left(W_{xy}\overline{G_{xx}G_{ij}G_{ji}G_{jy}}\right)-2p|W_{xy}|^{2p-2}\text{Re}\left(W_{xy}\overline{G_{xx}G_{ii}G_{jj}G_{yj}}\right)$$ 
$$-2p|W_{xy}|^{2p-2}\text{Re}\left(G_{xi}G_{jy}\overline{G_{xx}G_{ii}G_{jy}}+G_{xi}G_{jy}\overline{G_{xi}G_{jj}G_{yy}}\right)$$
We can see that this quickly becomes unruly. Getting to exponent six, which is what we need for the bound on $\mathcal{Q}$, will simply be impossible without some way of abstracting away the details. Fortunately, we are only concerned with stochastic bounds, and as such, don't need the exact form of the terms. 

\subsection*{Abstract Form}
That said, our work of calculating the third derivative isn't wasted, because now we can make an educated guess as to the form that all subsequent steps would take. Namely, we should expect: 

\begin{equation*}
\partial_{H_{ij}}^{k_1}\partial_{H_{ji}}^{k_2} F(H)= \sum_{\ell=0}^{2p+k_1+k_2} \sum_{\vec\sigma\in \{-1,1\}^\ell}\sum_{\vec a\in \{x, y, i, j\}^\ell}\sum_{\vec b\in \{x, y, i, j\}^\ell} C_{\ell, \vec \sigma, \vec a, \vec b} \prod_{\alpha=1}^\ell \tilde{G}_{a_\alpha b_\alpha}(\sigma_\alpha)\tag{8.6},
\end{equation*}
where $C_{\ell, \vec \sigma \vec a, \vec b}$ is some constant that is polynomial in $p$ (which includes $0$), and $\tilde{G}_{a_\alpha,b_\alpha}(\sigma_\alpha)$ is a function of the resolvent, representing all possible configurations , namely: 
\begin{itemize}
\item If $a_\alpha=x \wedge b_\alpha=y$, then $\tilde{G}_{a_\alpha, b_\alpha}(1) = W_{xy}$ and  $\tilde{G}_{a_\alpha, b_\alpha}(-1) = \overline{W_{xy}}$.
\item For $(a_\alpha=x\wedge b_\alpha=y)^C$, we have $\tilde G_{a_\alpha b_\alpha}(1) = G_{a_\alpha b_\alpha}$ and $\tilde G_{a_\alpha b_\alpha}(-1) = \overline{G_{a_\alpha b_\alpha}}$.
 \end{itemize}
Note that all terms of the form $|W_{xy}|$ that we encountered in our derivative expansion were in fact always raised to an even power $|W_{xy}|^{2p_0}$, which means that they can be represented as multiples of $|W_{xy}|^{2} = W_{xy}\cdot \overline{W}_{xy}$.  We can trivially see that the base case is satisfied by the definition of $F(H)$ and the first derivative we proved. Let us assume $(8.6)$ for the induction hypothesis. Now, we want to prove the cases for $k_1+1$ and $k_2+1$. Starting with the first one, we can differentiate $(8.6)$, s.t. by linearity and the product rule: $$\partial_{H_{ij}}^{k_1+1}\partial_{H_{ji}}^{k_2} F(H) = \partial_{H_{ij}}\left(\partial_{H_{ij}}^{k_1}\partial_{H_{ji}}^{k_2} F(H)\right) = \sum_{\ell=0}^{2p+k_1+k_2} \sum_{\vec\sigma}\sum_{\vec a}\sum_{\vec b} C_{\ell, \vec \sigma, \vec a, \vec b}\  \partial_{H_{ij}}\left(\prod_{\alpha=1}^\ell \tilde{G}_{a_\alpha b_\alpha}(\sigma_\alpha)\right)=$$
$$=\sum_{\ell=0}^{2p+k_1+k_2} \sum_{\vec\sigma}\sum_{\vec a}\sum_{\vec b}C_{\ell, \vec \sigma, \vec a, \vec b}\ \sum_{k=1}^\ell \left(\left(\partial_{H_{ij}}\tilde{G}_{a_k b_k}(\sigma_k)\right) \prod_{\alpha=1,\alpha\neq k}^\ell \tilde{G}_{a_\alpha b_\alpha}(\sigma_\alpha)\right) $$
We now need to calculate $\partial_{H_{ij}}\tilde{G}_{a_k b_k}(\sigma_k)$, depending on the cases, delineated above:
\begin{itemize}

\item For $a_k=x, b_k=y, \sigma_k=1$ we have $\partial_{H_{ij}} W_{xy} = -G_{xi}G_{jy}$

\item Similarly, for $a_k=x, b_k=y, \sigma_k=-1$ we have $\partial_{H_{ij}} \overline{W_{xy}} = -\overline{G_{xj}G_{iy}}$

\item For $(a_k=x\wedge b_k=y)^C$ and $\sigma_k = 1$ the result is the exact identity $\partial_{H_{ij}} G_{a_k b_k} = -G_{a_k i}G_{j b_k}$

\item Lastly, for $(a_k=x\wedge b_k=y)^C$ and $\sigma_k=-1$, we similarly have $\partial_{H_{ij}} \overline{G_{a_k b_k}} = -\overline{G_{a_k j}G_{i b_k}}$
\end{itemize}
All elements are clearly of the form $\tilde{D}_{a_k b_k} = \tilde{G}_{xi}\tilde{G}_{jy}\vee \tilde{G}_{a_k i}\tilde{G}_{j b_k}$ This means that:
$$\partial_{H_{ij}}^{k_1+1}\partial_{H_{ji}}^{k_2} F(H) = \sum_{\ell=0}^{2p+k_1+k_2} \sum_{\vec\sigma}\sum_{\vec a}\sum_{\vec b}C_{\ell, \vec \sigma, \vec a, \vec b}\ \sum_{k=1}^\ell \left(\partial_{H_{ij}} \tilde{G}_{a_k b_k}(\sigma_k) \prod_{\alpha=1,\alpha\neq k}^\ell \tilde{G}_{a_\alpha b_\alpha}(\sigma_\alpha)\right) = $$
$$ = \sum_{\ell=0}^{2p+k_1+k_2} \sum_{\vec\sigma}\sum_{\vec a}\sum_{\vec b}C_{\ell, \vec \sigma, \vec a, \vec b}\ \sum_{k=1}^\ell \left(\tilde{D}_{a_k b_k} \prod_{\alpha=1,\alpha\neq k}^\ell \tilde{G}_{a_\alpha b_\alpha}(\sigma_\alpha)\right)$$
But observe that the term inside the bracket is precisely the product of $\ell+1$ terms of the form $\tilde G_{a_ib_i}$ ($\ell-1$ from the product and $2$ from $\tilde{D}_{a_kb_k}$), whereas the sum over $k$ can be factored in as part of the sums over vectors $\vec a$ and $\vec b$ as an additional index, making the vectors $\ell+1$-dimensional, i.e: 

$$\partial_{H_{ij}}^{k_1+1}\partial_{H_{ji}}^{k_2} F(H) = \sum_{m=1}^{2p+k_1+k_2+1} \sum_{\vec\sigma'}\sum_{\vec a'}\sum_{\vec b'} C'_{m, \vec\sigma', \vec a', \vec b'} \prod_{\alpha=1}^m\ \tilde{G}_{a'_\alpha b'_\alpha}(\sigma'_\alpha).$$
 This clearly matches the form of $(8.6)$. And since the proof for $\partial_{H_{ij}}^{k_1}\partial_{H_{ji}}^{k_2+1} F(H)$ follows WLOG, we have indeed completed the induction. Let us now establish a stochastic bound on $\partial_{H_{ij}}^{k_1}\partial_{H_{ji}}^{k_2} F(H)$. Recall the minimal resolvent condition we assumed $(8.1.3)$, namely $\max_{a, b}|G(z)_{xx}|\prec 1$ that $\max_{a, b}|G(z)_{ab}-m(z)\delta_{ab}|\prec W^{-\delta}$ for some fixed $\delta>0$. By the triangle inequality, applied to $(8.6)$, we get: 

$$\left\vert \partial_{H_{ij}}^{k_1}\partial_{H_{ji}}^{k_2} F(H)\right\vert \leq    \sum_{\ell=0}^{2p+k_1+k_2} \sum_{\vec\sigma\in \{-1,1\}^\ell}\sum_{\vec a\in \{x, y, i, j\}^\ell}\sum_{\vec b\in \{x, y, i, j\}^\ell} |C_{\ell, \vec \sigma, \vec a, \vec b}| \prod_{\alpha=1}^\ell \left\vert \tilde{G}_{a_\alpha b_\alpha}(\sigma_\alpha)\right\vert$$
Observe that sum is finite, where the number of terms is $\sum_{\ell=0}^{2p+k_1+k_2} 2^{\ell}\cdot 4^\ell\cdot 4^\ell = \sum_{\ell=0}^{2p+k_1+k_2}32^\ell=M$. By the same token, we can take $\mathcal{C} = \max_{\ell, \vec\sigma, \vec a, \vec b} |C_{\ell, \vec\sigma, \vec a, \vec b}|$. Using the product rule of stochastic domination (\hyperref[lemma-a6]{Lemma A6}) and the resolvent conditions, we have that: 
$$\mathcal{C}\prod_{\alpha=1}^\ell \left\vert \tilde{G}_{a_\alpha b_\alpha}(\sigma_\alpha)\right\vert \prec 1^{\kappa_1} \cdot W^{-\delta\kappa_2} \prec W^{-C_{\kappa_1,\kappa_2}},$$
 where $\kappa_{1}, \kappa_2=O(1)$ are some integer constants  $0\leq \kappa_1, \kappa_2\leq \ell$ that depend on the particular vectors $\vec a, \vec b, \vec \sigma$, since different terms differ in resolvent combinations, and as such will have different bounds. We don't care for the exact details, since we know that stochastic domination is also preserved under summation (\hyperref[lemma-a7]{Lemma A7}), and such this is enough for us to conclude that $\exists C=O(1)$ (some deterministic function of all the $C_{\kappa_1, \kappa_2}$), s.t. $|\mathcal{Q}|\prec W^{-C}$. That is all we need for the last part of our proof. 

\subsection{Continuity Argument}
\noindent Having established the stochastic bound on the individual derivatives, we can now lay out the continuity argument that will allows us to preserve control throughout the entire replacement. To start, define : 

\begin{equation*}
\begin{array}{c}
G_1 = G\left(H_{11}, \cdots, H_{ij}^G, \cdots, H_{NN}^G\right)\\
G_2 = G\left(H_{11}, \cdots,\tilde{t}_{ij}^G H_{ij}^G=\cdots H_{NN}^G\right) \\
G_3 = G\left(H_{11}, \cdots,\tilde{t}_{ij} H_{ij}, \cdots, H_{NN}^G\right),
\end{array},
\end{equation*}
where $\tilde{t}_{ij}^G$ and $\tilde{t}_{ij}$ are the replacement variables $\in [0, 1]$ at step $(i,j)$ for this argument, different than the ones we used in our Lindeberg argument. Our goal is the following - assuming that $G_1$ satisfies the resolvent condition $(8.1.3)$, we will now prove that $G_2$ and $G_3$ also do. This will be equivalent to verifying that the interpolation at step $(i,j)$ preserves our resolvent bounds, which we need to establish sufficient control throughout the entire replacement of the $NW$ non-zero entries in the band matrix. To start, let us use the matrix identity $(2.4)$ for $A = H_2-z$ and $B = H_1-H_2$:
$$G_1 = G_2 - G_2\left(H_1-H_2\right)G_1$$
$$\Rightarrow G_2 = G_1+G_2\left(H_1-H_2\right)G_1$$

\noindent Given that $H_2$ and $H_1$ differ only in their $(ij)$th and $(ji)$th component, we can represent their difference as the sum of two matrices with the only non-zero entries being the latter entries. Those are, by definition of the orthonormal basis, $e_i^*e_j$ and $e_j^*e_i$, respectively. Then, for $c = 1-\tilde{t}_{ij}^G\in [0, 1]$, we have:
\begin{equation*}
G_2= G_1 + G_2\left(cH_{ij}^Ge_i^*e_j+ cH_{ji}^G e_j^*e_i\right)G_1
\end{equation*}

$$\Rightarrow (G_2)_{xy} = (G_1)_{xy} + cH_{ij}^G (G_2)_{xi}(G_1)_{jy}+{c}H_{ij}^G (G_2)_{xj}(G_1)_{iy}$$
Given $(8.1.2)$ and $(8.1.3)$, we have:
$$|(G_2)_{xy}|\prec 1 + W^{-\frac{1}{2}}|(G_2)_{jy}|+W^{-1/2}|(G_2)_{iy}|$$
Let $\Xi = \max_{x, y}|(G_2)_{xy}|$. We can rewrite the expression above as: 
\begin{equation*}\Xi \prec 1 + W^{-1/2}\Xi\tag{8.7}\end{equation*}

\noindent Now, let $\epsilon, D>0$ be arbitrary. Given that we work in the context of the delocalization, we have $W\gg N^{\frac{1}{2}+\beta} = N^{\frac{1+2\beta}{2}}\Rightarrow W^{-1/2}\ll N^{-\frac{1+2\beta}{4}}$ for any fixed $\beta>0$. Hence, let us choose $\delta \in \left(0, \frac{2\beta+1}{4\epsilon}\right)$ by the density of the rationals. By Def.1 of stochastic domination, from $(8.7)$, we have that: $$\sup_{s\in S_N}\mathbb{P}\left(\Xi_N(s)>N^{\delta\epsilon}(1+W^{-1/2}\Xi_N(s))\right)<N^{-D}$$
$$\sup_{s\in S_N}\mathbb{P}\left(\Xi_N(s)>N^{\delta\epsilon}+N^{\delta\epsilon} W^{-1/2}\Xi_N(s)\right)<N^{-D}$$
$$\sup_{s\in S_N}\mathbb{P}\left(\Xi_N(s)(1-N^{\delta\epsilon} W^{-1/2})>N^{\delta\epsilon}\right)<N^{-D}$$
$$\sup_{s\in S_N}\mathbb{P}\left(\Xi_N(s)>\frac{N^{\delta\epsilon}}{(1-N^{\delta\epsilon} W^{-1/2})}\right)<N^{-D}$$
By the delocalization scale and per our choice $\delta\epsilon <\frac{2\beta+1}{4\epsilon}$, we have that $$N^{\delta\epsilon}W^{-1/2} \ll N^{\delta\varepsilon} \cdot N^{-1/4-\beta/2} = N^{\delta\varepsilon-1/4-\beta/2}$$
$$1-N^{\delta\epsilon} W^{-1/2}\gg 1-N^{\delta\epsilon -\frac{1+2\beta}{4}}$$
$$\frac{N^{\delta\epsilon}}{1-N^{\delta\epsilon} W^{-1/2}}\ll \frac{N^{\delta\epsilon}}{1-N^{{\delta\epsilon}-\frac{1+2\beta}{4}}}$$

\noindent This means that $\exists N_0$, s.t. for all $N\geq N_0$, the latter event is a subset: $$ \left\{\Xi_N(s)>\frac{N^\epsilon}{1-N^{{\delta\epsilon}-\frac{1+2\beta}{4}}}\right\}\subset \left\{\Xi_N(s)>\frac{N^{\delta\epsilon}}{(1-N^{\delta\epsilon} W^{-1/2})}\right\}$$
$$\Rightarrow \mathbb{P} \left\{\Xi_N(s)>\frac{N^{\delta\epsilon}}{1-N^{\delta\epsilon-\frac{1+2\beta}{4}}}\right\}\leq\mathbb{P} \left\{\Xi_N(s)>\frac{N^{\delta\epsilon}}{(1-N^{\delta\epsilon} W^{-1/2})}\right\}\leq N^{-D},$$
by Boole's inequality. But observe that we can take $\delta$ to be small enough, s.t. 
$$N^{\delta\epsilon-\frac{1+2\beta}{4}}\leq \frac{1}{2}\quad \Leftrightarrow \quad 1-N^{\delta\epsilon-\frac{1+2\beta}{4}}\geq \frac{1}{2}\quad \Rightarrow \quad \frac{N^{\delta\epsilon}}{1-N^{\delta\epsilon-\frac{1+2\beta}{4}}}\leq 2N^{\delta\epsilon}$$
$$  \left\{\Xi_N(s)> 2N^{\delta\epsilon}\right\}\subset \left\{\Xi_N(s)>\frac{N^{\delta\epsilon}}{1-N^{\delta\epsilon-\frac{1+2\beta}{4}}}\right\}$$
$$\Rightarrow \mathbb{P} \left\{\Xi_N(s)> 2N^{\delta\epsilon}\right\}\leq\mathbb{P} \left\{\Xi_N(s)>\frac{N^{\delta\epsilon}}{1-N^{\delta\epsilon-\frac{1+2\beta}{4}}}\right\}\leq N^{-D},$$
\noindent by Boole's inequality again. But since $\delta\epsilon$ is small, this is exactly what we need per the definition of stochastic domination (up to trivial rescaling w.r.t to $\epsilon$), meaning that we have verified that $\Xi\prec 1$. Notice that the approach is equivalent for $\Xi_3 = \max_{x, y}|(G_3)_{xy}|$, i.e WLOG, we have that $\Xi_3\prec 1$ as well. What we have essentially established is a sort of a continuity argument. Namely, at every step $(i,j)$ of the replacement, the proof above demonstrates that we preserve the resolvent bounds. Now remains the last step is extending this argument throughout the replacement of the entire matrix. 

\subsection{Five Moment Theorem}
All our work up to this point will allows us to finally bound the the difference: $$\left\vert \mathbb{E}\left[ F(H_{ij} - F(H_{ij}^G)\right]\right\vert \lesssim \mathbb{E}\left\vert\left.\frac{d^6}{dt^6}F(tH_{ij})\right\vert_{t=t_{ij}}\right\vert+\mathbb{E}\left\vert\left.\frac{d^6}{dt^6}F(tH_{ij}^G)\right\vert_{t=t_{ij}^G}\right\vert = \mathcal{P} + \mathcal{P}^G$$
We showed in the last section that $\mathcal{Q}=\left\vert \partial_{H_{ij}}^\ell \partial_{H_{ji}}^{6-\ell} F(tH_{ij})|_{t=t_{ij}}\right\vert\prec W^{-C}$ for some $C\lesssim 1$. By the definition stochastic domination, for any $\epsilon>0$ and $D>0$, there exists an event $\mathcal{A}$, s.t. $\mathcal{Q}\leq W^\epsilon$ on $\mathcal{A}$ with $P(\mathcal{A})\geq 1-W^{-D}$. Hence, we can split $\mathcal{P}$ as in $(8.5.1)$ and use Cauchy-Schwarz: 
$$\mathcal{P}\lesssim \mathbb{E}|H_{ij}|^6W^\epsilon + \mathbb{E}|H_{ij}|^6 \mathbb{E}\left[\mathbf{1}_{A^C}\mathcal{Q}\right]\lesssim W^{-3+\epsilon} + W^{-C} \left\{\mathbb{E}|H_{ij}|^{12}\right\}^{1/2} \left\{P(A^C)\right\}^{1/2}\lesssim$$
$$\lesssim W^{-3+\epsilon} + W^{-C} W^{-3} W^{-D_0},$$
where the last bound follows from $(8.1.2)$ and our probability bound on $\mathcal{A}$. But this means that we can pick $D_0$ large enough so the stochastically dominating term is $W^{-3+\epsilon}$, i.e. $\mathcal{P}+\mathcal{P}^G\lesssim W^{-3+\epsilon}$: $$\left\vert \mathbb{E}\left[F(H_{ij})\right] - \mathbb{E}\left[F(H_{ij}^G)\right]\right\vert\lesssim W^{-3+\epsilon}$$
Then, by applying the triangle inequality to the telescoping sum of replacements, and accounting for the $NW$ non-zero entries in the band, we have: 
 \begin{equation*} \left\vert \mathbb{E}\left[F(H)\right] - \mathbb{E}\left[F(H^G)\right]\right\vert\lesssim NWW^{-3+\epsilon} = NW^{-2+\epsilon}\tag{8.8}
 \end{equation*}
\noindent For the condition of delocalization, we need for any $\beta>0$ the following to hold:
 $$W\gg N^{\frac{1}{2}+\beta} = N^{\frac{1+2\beta}{2}}\Leftrightarrow N\ll W^{\frac{2}{1+2\beta}}$$
\noindent Bounding $(8.8)$ by the latter condition gives us:
  $$NW^{-2+\epsilon}\ll W^{-2+\epsilon+\frac{2}{1+2\beta}}$$
\noindent Hence, in order for our asymptotics to work, we need a negative exponent, i.e: $$-2+\epsilon + \frac{2}{1+2\beta}<0$$
  $$0<\epsilon <2-\frac{2}{1+2\beta} = \frac{4\beta}{1+2\beta}>0,$$
  whereas since $\beta>0\Rightarrow \frac{4\beta}{1+2\beta}$, we can always pick such an $\epsilon>0$ by the density of the rationals, completing the proof. Note, this is also why we needed a five matching moments assumption - if we had four instead, then the bound we would get per condition $(8.1.2)$ would be $\left\vert \mathbb{E}\left(F(H_{ij}\right) - \mathbb{E}\left(F(H_{ij}^G\right)\right\vert\lesssim W^{-\frac{5}{2}+\epsilon}$. By following the same steps as above, we would get $-\frac{3}{2}+\epsilon+\frac{2}{1+2\beta}<0\Leftrightarrow \epsilon< \frac{6\beta-1}{1+2\beta}$.  But then $\epsilon$ would be negative for $\beta<\frac{1}{6}$, meaning the delocalization bound does not hold. 
  







\newpage




\section{Appendix}
\label{sec:appendix}

\begin{boxt}{Lemma A.1.1} \label{lemma-a11} $$ 1-|m|^2 \asymp \eta$$ 
\end{boxt}
\textit{Proof} - can be found in Lemma 3.5 \cite{21} and Lemma 6.2 \cite{dynamic}). 

\begin{boxt}{Lemma A.1.2} \label{lemma-a12} 
\begin{enumerate}
\item $ \im w_s = \im w_t + (t-s)\im m(z)$ 
\item $(t-s)\lesssim \eta_s $
\end{enumerate}

\end{boxt}
\textit{Proof}- Recall the definition $w_s = -\frac{1}{m(z)} - sm(z)$, where $m(z)$ is the Stieltjes transform. By taking its imaginary part, we have: $$\im w_s = -\im \left( \frac{1}{m(z)}\right)-s\im m(z)= -\im \left( \frac{1}{m(z)}\right)-t\im m(z) + (t-s)\im m(z)$$
$$=\im w_t + (t-s)\im m(z),$$
which proves the first identity. Rearrange as: 
$$\im w_s\geq \im w_s\im w_t= (t-s)\im m(z)$$
Since $|\im m(z)|$ is bounded uniformly away from $0$ in the bulk (Lemma 6.2. in \cite{dynamic} and $\im w_t\geq 0$ by def. , we can take the absolute value divide by $\im m(z)$ the result above, s.t. we get $(t-s)\lesssim \eta_s$. 

\begin{boxt}{Lemma A.2} \label{lemma-a2} -
 For $|E|<2$ fixed and $\eta\gtrsim W^2/N^2$, then $$(\Theta_t)_{ab} + |(\Theta_tS^{\frac{1}{2}})_{ab}| + |\left(S^{\frac{1}{2}}\Theta_t S^{\frac{1}{2}}\right)_{ab}|\lesssim W^{-1}\eta_t^{-\frac{1}{2}}, $$
 $$\max_a \sum_b |(\Theta_t)_{ab}|+\max_a \sum_b |(\Theta_tS^{\frac{1}{2}})_{ab}|\lesssim \eta_t^{-1}$$
\end{boxt}
Proof can be found as Lemma 23 in \cite{bandSDE}. From it clearly follows that $\max_{a,b}\Theta_{ab} \lesssim W^{-1}\eta^{-1/2}$. 

\begin{boxt}{Corollary A.3} \label{cor-a3} 
$$\sup_x \sum_y \{\text{Id}+(t-s)\Theta_t\}_{xy} + \sup_y \sum_x \{\text{Id} + (t-s)\Theta_t\}_{xy} = 1+O(\eta_t^{-1}\eta_s)$$
 \end{boxt}
 \textit{Proof} - Can be found as expression $(3.7)$ in the proof of Lemma 9 \cite{bandSDE}. 

\begin{boxt}{Lemma A.4} \label{lemma-a4} - For $|E|<2$ fixed and arbitrary $a, b$, we have for $B_t  =(\text{Id} - tm(z)^2S)^{-1}$: $$\sum_\alpha |(B_t)_{\alpha b}| + \sum_{\beta}|(B_t)_{a\beta}|\lesssim 1$$
 \end{boxt}
 \textit{Proof} - found as Lemma 25 in \cite{bandSDE}. Clearly it implies that $\sum_y |B_{xy}|$ and $\sum_y |S_{xy}|$ are $\lesssim 1$.  


 
  

\begin{boxtheorem}{Lemma A5: (Large Deviation estimates)}    \label{lemma-a5}

  

Let $\left(X_i^{(N)}\right)$ and $\left(Y_i^{(N)}\right)$ are independent families of random variables and $\left(a_{ij}^{(N)}\right)$ and $\left(b_{ij}^{(N)}\right)$ be deterministic. Suppose all entrires $X_{i}^{(N)}$ and $Y_{i}^{(N)}$ are independent and satisfy: $$\E X = 0, \quad \E |X|^2=1, \quad \|X\|_p:=(\E |X|^p)^{1/p}\leq \mu_p$$

for all $p\in \N$ and some constants $\mu_p$. Then, we have the bounds: 

\begin{equation*}\sum_i b_i X_i \prec \left(\sum_i |b_i|^2\right)^{1/2}\tag{1.1}\end{equation*}

\begin{equation*}\sum_{i, j} a_{ij}X_iY_j \prec \left(\sum_{i,j} |a_{ij}|^2\right)^{1/2}\tag{1.2}\end{equation*}

\begin{equation*}\sum_{i\neq j} a_{ij}X_iY_j \prec \left(\sum_{i\neq j} |a_{ij}|^2\right)^{1/2}\tag{1.3}\end{equation*}

\end{boxtheorem}
\textit{Proof} - can be found as Theorem 7.7 in Yau's Dynamic RMT textbook \cite{dynamic}. 


\begin{boxtheorem}{Lemma A6: (Stochastic Domination Product)}\label{lemma-a6}

If $X_i \prec Y_i$ for $i = 1,2,...,\ell$, then $\prod_{i=1}^{\ell} X_i \prec \prod_{i=1}^{\ell} Y_i$.
\end{boxtheorem}
\textit{Proof} - By definition, for any $\epsilon, D > 0$, there exist $N_{\epsilon,D,i}$ such that: $$\sup_{s\in S_N}\mathbb{P}\left(X_i(s) > N^{\epsilon/\ell} Y_i(s)\right) < N^{-(D+\log_N \ell)}$$ for all $N \geq N_{\epsilon,D,i}$. For the product, consider: $$\mathbb{P}\left(\prod_{i=1}^{\ell} X_i(s) > N^{\epsilon} \prod_{i=1}^{\ell} Y_i(s)\right).$$
This event occurs only if at least one $X_i(s) > N^{\epsilon/\ell} Y_i(s)$. Then $\forall N \geq \max_{i=1}^{\ell} N_{\epsilon,D,i}$ by Boole's inequality: $$\mathbb{P}\left(\prod_{i=1}^{\ell} X_i(s) > N^{\epsilon} \prod_{i=1}^{\ell} Y_i(s)\right) \leq \sum_{i=1}^{\ell} \mathbb{P}\left(X_i(s) > N^{\epsilon/\ell} Y_i(s)\right) < \ell \cdot N^{-(D+\log_N \ell)} = N^{-D}$$ 
\begin{boxtheorem}{Lemma A7: (Stochastic Domination Sum)}\label{lemma-a7}

If $X_i \prec Y_i$ for $i = 1,2,...,\ell$, then $\sum_{i=1}^{\ell} X_i \prec \sum_{i=1}^{\ell} Y_i$.
\end{boxtheorem}
\textit{Proof} - Analogous to Lemma A6. 
\subsection*{Resolvent estimates}


\begin{boxtheorem}{Lemma B.1} \label{lemma-b1}


$$\sup_x \sum_y |G_{xy}|\lesssim W^{\frac{\delta_\sto}{20}}W^{\frac{1}{2}+\epsilon} \eta_s^{-\frac{3}{4}}$$
\end{boxtheorem}
\textit{Proof} - can be found as result $4.5$ in \cite{bandSDE}.

\begin{boxtheorem}{Lemma B.2} \label{lemma-b2}

$$\max_{x, y} |G_{xy} - \delta_{xy}m(z)|\lesssim W^{\frac{\delta_\sto}{20}}W^{-\frac{1}{2}} \eta_s^{-\frac{1}{4}}
$$
\end{boxtheorem}
\textit{Proof} - can be found as result $4.6$ in \cite{bandSDE}). 


\begin{boxtheorem}{Lemma B.3 (Off-diagonal $G_{xy}$)} \label{lemma-b3} - Let $\max_x|G_t(z)_{xx}|\prec 1$ and $\max_{x\neq y}|G_t(z)_{xy}|\prec W^{-\delta}$ for some $\delta>0$. Then, for any fixed $t\in [0, 1]$:$$\max_{a\neq b} \frac{|G_t(z)_{ab}|^2}{(S^{1/2}T_t(z)S^{1/2})_{ab}+S_{ab}^{1/2}}\prec 1$$

\end{boxtheorem}
\textit{Proof} - can be found as Lemma 12 in \cite{bandSDE}\label{lemma-b4}
\begin{boxtheorem}{Lemma B.4 (On-diagonal $|G_{xx}-m(z)|$)} - Suppose the following two conditions for some $\delta>0$:$$\max_{a,b}|G_t(z)_{ab}-m(z)\delta_{ab}|\leq W^{-\delta}$$
$$\max_{a\neq b}|G_t(z)_{ab}|\prec \max_{a, b}|(S^{1/2}T_t(z)S^{1/2})_{ab}|+\max_{a, b}|S_{ab}^{1/2}|$$
Then, the following holds: $$\max_x |G_t(z)_{xx} - m(z)|^2\prec \max_{a, b}|S^{1/2}T_t(z)S^{1/2})_{ab}|+\max_{a,b}|S_{ab}^{1/2}|+W^{-1}$$

\end{boxtheorem}
\textit{Proof} - can be found as Lemma 13 in \cite{bandSDE}

\begin{boxtheorem}{Corollary B5}\label{corollary-b5} - Let $\mathcal{A}$ be an even on which the following holds for all $t\in [0, 1]$: $$\max_{a, b}\left(S^{1/2}T_t(z)S^{1/2}\right)_{ab}\lesssim W^{-1}\eta_t^{-1/2}$$
Then, we have the following stochastic bound for any large $C>0$ that is independent of $W$: $$\mathbf{1}_{\mathcal{A}}\max_{t\in [0, 1]} \max_{x, y}\frac{|G_{xy} - \delta_{xy}m(z)|}{(S^{1/2}T_t(z)S^{1/2})_{ab}+S^{1/2}_{ab}+W^{-C}}\prec 1$$
.
\end{boxtheorem}
\textit{Proof} - the result follows by Lemmas B.3 and B.4 from a continuity argument. Can be found as Corollary 14 in \cite{bandSDE}

\newpage 
\addcontentsline{toc}{section}{References}

\begin{thebibliography}{99}

\bibitem{wigner} E. Wigner. Distribution of neutron resonance level spacing. \textit{Columbia University Report CU175}, 1957.

\bibitem{dynamic} L. Erdős, H.-T. Yau. A Dynamical Approach to Random Matrix Theory. \textit{American Mathematical Society}, 2017.

\bibitem{mehta} M. L. Mehta. Random Matrices and the Statistical Theory of Energy Levels. \textit{Academic Press},
New York, NY, 1967.

\bibitem{taoWDM} T. Tao and V. Vu. The Wigner-Dyson-Mehta Bulk Universality Conjecture for Wigner Matrices. \textit{Electronic Journal of Probability} 16:2104-2121, 2011

\bibitem{taoRMT} T. Tao. Topics in random matrix theory. \textit{American Mathematical Society}, 2012.

\bibitem{anderson} P. W. Anderson. Absence of diffusion in certain random lattices. \textit{Physical Review Letters}, 109: 1492-1505, 1958.

\bibitem{randomoperators} M. Aizenman and S. Warzel. Random Operators: Disorder Effects on Quantum Spectra and Dynamics. \textit{American Mathematical Society}, 2015.

\bibitem{24} J. Frölich and T. Spencer. Absence of diffusion in the anderson tight binding model for large disorder or low energy. \textit{Communications in Mathematical Physics}, 83:151-184, 1988.

\bibitem{1} M. Aizenman and S. Molchanov. Localization at large disorder and at extreme energies: an elementary derivation. \textit{Communications in Mathematical Physics}, 157: 245-278, 1993. 

\bibitem{backer2019} A. Bäcker, et al. Multifractal Dimensions for Random Matrices, Chaotic Quantum Maps, and Many-Body Systems. \textit{Physical Review}, 100: 032117–032117, 2019.


\bibitem{bandSDE} S. Dubova and K. Yang. Quantum diffusion and delocalization in one-dimensional band matrices via the flow method. \textit{arXiv.2412.15207}, 2024.


\bibitem{7} P. Bourgade. Random band matrices. \textit{Proceedings of the International Congress of Mathematicians}, p. 2759-2783, 2018.

\bibitem{9} P. Bourgade, H.-T. Yau, and J. Yin. Random band matrices in the delocalized phase, I: Quantum unique ergodicity and universality. \textit{Communications on Pure and Applied Mathematics}, 73: 1526-1596, 2020.

\bibitem{8} P. Bourgade, F. Yang, H.-T. Yau, and J. Yin. Random band matrices in the delocalized phase, II: Generalized resolvent estimates. \textit{Communications on Pure and Applied Mathematics}, 174: 1189-1221, 2019.


\bibitem{40} F.Yang, H.-T. Yau, and J. Yin. Random band matrices in the delocalized phase, III: averaging fluctuations. \textit{Probability Theory and Related Fields}, 179: 451-540, 2021.

\bibitem{highdim} C. Xu et al. Bulk universality and quantum unique ergodicity for random band matrices in high dimensions. \textit{The Annals of Probability} 52.3:765–837, 2024.


\bibitem{mirlin1996} A. D. Mirlin, Y. V. Fyodorov, F. M. Dittes, et al. Transition from localized to extended eigenstates in the ensemble
of power-law random banded matrices.\textit{ Physical Review} 54: 3221, 1996.

\bibitem{yauyin} H.-T. Yau and J. Yin. Delocalization of One-Dimensional Random Band Matrices. \textit{arxiv.2501.01718}, 2025.

\bibitem{twodim} S. Dubova, K. Yang, H.-T. Yau, J. Yin. Delocalization of Two-Dimensional Random Band Matrices. \textit{arxiv.2503.07606}, 2025. 

\bibitem{kravtsov2009}  V. E. Kravtsov. Random Matrix Representations of Critical Statistics. \textit{The Oxford Handbook of Random Matrix Theory}, 2011


\bibitem{knowles2019} F. Beynach-Georges and A. Knowles. Lectures on the local semicircle law for Wigner matrices. \textit{arxiv.1601.04055}, 2016.

\bibitem{arbm1} C. Casati, L. Molinari, and F. Izrailev. Scaling properties of band random matrices. \textit{Physical Review Letters}, 64: 1851-1854, 1990.

\bibitem{arbm2} Y. V. Fyodorov and A.D. Mirlin. Scaling properties of localization in random band matrices: a sigma-model approach . \textit{Physical Review}, 67, 1991.

\bibitem{13} C. Cipolloni, R. Peled, J. Schenker, and J, Shapiro. Dynamical localization for random band matrices up to {$W\ll N^{1/4}$}. \textit{Communications in Mathematical Physics}, 205, 2024. 

\bibitem{porter1960}  N. Rosenzweig and C. E. Porter. Repulsion of Energy Levels in Complex Atomic Spectra. \textit{Physical Review}, 120: 1698-1714, 1960.

\bibitem{RP} P. von Soosten and S. Warzel. Non-Ergodic Delocalization in the Rosenzweig–Porter Model.\textit{Letters in Mathematical Physics}, 109, 2019.

\bibitem{steele} J. M. Steele. Stochastic Calculus and Financial Applications.\textit{ Springer New York}, 2001.

\bibitem{bdg} D.L. Burkholder, B. Davis, R.F. Gundy. Integral inequalities for convex functions of operators on martingales. \textit{Berkeley Symp. on Math. Statistics and Probability}, 2: 223–240, 1972.

\bibitem{21} L. Erdős, A. Knowles, H.-T. Yau, and J.Yin. Delocalization and Diffusion Profile for Random Band Matrices. \textit{Communications in Mathematical Physics}, 323: 367-416, 2013.

\bibitem{yau2011} L. Erdős, H.-T. Yau , and J. Yin. Universality for generalized Wigner matrices with Bernoulli distribution. \textit{Journal of Combinatorics}, Vol. 2, 1: 15–81, 2011.

\bibitem{38} F.Yang, H.-T. Yau, and J. Yin. Delocalization and quantum diffusion of random band matrices in high dimensions I: Self-energy renormalization. \textit{arXiv:2104.12048}, 2021.

\bibitem{39} F.Yang, H.-T. Yau, and J. Yin. Delocalization and quantum diffusion of random band matrices in high dimensions II: T-expansion. \textit{Communications in Mathematical Physics}, 396: 527-622, 2022.

\bibitem{soosten} P. von Soosten and S. Warzel. Random characteristics for Wigner matrices.  \textit{Electronic Communications in Probability} 24: 1-12 , 2019.

\bibitem{17} J. Edwards and D. Thouless. Numerical studies of localization in disordered systems. \textit{Journal of Physics
C: Solid State Physics}, 5:807–820, 1972.

\bibitem{32} T. Spencer. Random banded and sparse matrices. \textit{Oxford handbook of random matrix theory}, pages
471–488, 2015.

\bibitem{33} D. Thouless. Maximum metallic resistance in thin wires. \textit{Physical Review Letters}, pages 1167–1169,
1977.


\bibitem{lindeberg} J.W. Lindeberg. Eine neue Herleitung des Exponentialgesetzes in der Wahrscheinlichkeitsrechnung, Math. Z. 15 (1922), 211-225. 

\bibitem{chatterjee} S. Chatterjee. A generalization of the Lindenberg principle. \textit{The Annals of Probability}, 34: 2061-2076, 2006.
\bibitem{tao} T. Tao and V. Vu. Random Matrices: The Universality Phenomenon for Wigner Ensembles. \textit{ Proceedings of Symposia in Applied Mathematics} 72:01, 2011.

\bibitem{tao2} T. Tao and V. Vu. Random Matrices: Universal Properties of Eigenvectors. \textit{Random Matrices: Theory and Applications} 01:01, 2012.

%\bibitem{spectralgap} Find reference


\end{thebibliography}


\end{document}
